\chapter{Finite-energy measures and variational solutions}
\index{finite-energy measure!BFJ}\index{Monge--Ampere equation!variational solution}
\label{chap:bfj-finite-energy-measures}

The previous two chapters put the two sides of the Monge--Amp\`ere equation
on an equal footing.  A potential $\varphi\in\mathcal E^1$ has an energy
and a full-mass measure, while differentiability identifies the derivative of
energy with integration against that measure.  We now formulate the inverse
problem variationally: given a positive measure $\mu$ of the correct total
mass, find $\varphi$ with $\MA_\omega(\varphi)=\mu$.

The method is the one every reader of the complex theory will recognize.
One writes down a concave functional
$F_\mu(\varphi)=E_\omega(\varphi)-\int\varphi\,\rd\mu$ whose
Euler--Lagrange equation is the equation to be solved, one shows that
$F_\mu$ is bounded above and attains its supremum, and one differentiates.
Each of the three steps has a non-Archimedean cost.  Boundedness above is
exactly the statement that $\mu$ lies in the domain of a Legendre
transform, and that transform---the \emph{dual energy}
$E^\vee_\omega$---is the principal new object of the chapter.  Attainment
of the supremum uses compactness in its subnet form, without a
metrizability assumption on $\PSH_0(X,\omega)$
(\Cref{rem:ch28-net-form}).  Differentiation is
\Cref{prop:bfj-euler-lagrange}, and it consumes the orthogonality property.

Three results have different inputs and conclusions.  The chapter first
distinguishes them, then proves the implications needed for existence.
The first is a \emph{duality} statement: $E^\vee_\omega$ is
the concave conjugate of the energy, and the supremum of $F_\mu$ equals
$E^\vee_\omega(\mu)$ whenever the latter is finite.  This costs nothing
but the definitions.  The second is a \emph{completion} statement: the
Monge--Amp\`ere measures of model potentials are dense, in a strong measure
topology, in the space $\mathcal M^1_V$ of measures of finite dual energy.
Completion alone concerns measures.  In the present BFJ setting,
its quantitative estimates can be combined with the pointwise energy
compactness of \Cref{chap:bfj-energy} to obtain a finite-energy maximizer
for every $\mu\in\mathcal M^1_V$; the complete deduction is
\Cref{thm:ch28-finite-energy-existence}.  The third result is the
\emph{continuous} existence theorem of Boucksom--Favre--Jonsson for
measures supported on one SNC dual complex.  Support is an existence
input in their original proof and supplies the additional regularity in
\Cref{chap:bfj-continuity}.  It is not part of the definition of finite
energy or a necessary condition for a finite-energy solution.

\newpage
\begin{warning}{The scope of the existence theorem}
\Cref{thm:bfj-supported-cy} is a theorem about a smooth projective variety
over a complete \emph{discretely valued} field of residue characteristic
zero, with a semipositive model form of ample de Rham class having the
orthogonality property, and about a measure supported on the dual complex of
an SNC model.  Every one of these hypotheses is used.  It is
\emph{not} a theorem over an arbitrary complete non-Archimedean field: over
such a field the class $\cE^1$ can still be defined, but the extension of
the Monge--Amp\`ere operator to it is explicitly left open in
\cite[Remark~4.4]{PS23Hybrid}, and no assertion of this chapter presupposes
that extension.  \Cref{sec:ch28-scope} displays the full hypothesis
signature and lists, separately and with dates, the later theorems that
weaken individual hypotheses.
\end{warning}

\paragraph{Standing hypotheses.}
Unless a statement says otherwise we retain the geometric hypotheses fixed
at the beginning of \Cref{chap:bfj-model-forms} and used throughout
\Cref{chap:bfj-psh,chap:bfj-compactness,chap:bfj-envelopes,chap:bfj-ma,chap:bfj-capacity,chap:bfj-energy,chap:bfj-differentiability}:
$K$ is complete and discretely valued with residue characteristic zero and
$v_K(\pi)=1$; $X$ is smooth, connected and projective of dimension $n$
over $K$; and $\omega$ is a closed model $(1,1)$-form which is
semipositive and whose de Rham class $\{\omega\}$ is ample.  We write
\begin{equation}
\label{eq:ch28-standing}
 \begin{aligned}
 V&:=(\{\omega\}^n)>0,\\
 \mathcal M_V(X)
 &:=\Bigl\{\mu\ \Bigm|\
   \begin{gathered}
    \mu\ \text{is a positive Radon measure on }X^{\An},\\[-1mm]
    \mu(X^{\An})=V
   \end{gathered}\Bigr\},
 \end{aligned}
\end{equation}
and $\PSH_0(X,\omega)$ for the normalized psh space
\eqref{eq:ch22-normalized-space}.  All Monge--Amp\`ere measures below have
total mass $V$, never mass one; the probability normalization is recorded
in \eqref{eq:appJ-dual-energy} and is used only when a source is being
transported, as in \Cref{thm:finite-energy-measure-completion}.

Two further hypotheses are named rather than assumed globally.

\begin{itemize}
\item \emph{Orthogonality} (\Cref{def:orthogonality-property}) is assumed in
every statement that differentiates an envelope, and it is flagged in the
statement.  The general standing setting is covered by the volume route of
\Cref{thm:orthogonality-provenance}.  The elementary curve statement
\Cref{prop:ch27-curve-orthogonality} has its additional model-envelope
hypothesis; \Cref{thm:ch27-toric-orthogonality} concerns the invariant
envelope with a rational invariant model obstacle.
\item \emph{Support} hypotheses on $\mu$ are always displayed.  No theorem
here gives a finite-energy solution for every element of
$\mathcal M_V(X)$: finite dual energy is necessary.  It is sufficient under
orthogonality by \Cref{thm:ch28-finite-energy-existence}; continuity of
that solution is a further conclusion for skeletal-supported measures.
\end{itemize}

\paragraph{What the reader should learn.}
By the end of the chapter the reader should be able to: write down the
Legendre dual energy of a measure and explain why the total-mass condition
is forced rather than imposed; prove that the supremum of the variational
functional equals the dual energy, and compute both in closed form on a
metric graph; distinguish a maximizing net from a maximizer, and say
precisely which compactness statement upgrades the first to the second;
state the Boucksom--Favre--Jonsson existence theorem with its complete
hypothesis signature and identify the role of each hypothesis in the proof;
explain why uniqueness is not a corollary of the comparison principle;
compute the dual energy of an explicit measure on a curve skeleton and
exhibit a measure of infinite dual energy; and quantify the dependence of
the solution on the measure.

\paragraph{Prerequisites.}
From \Cref{chap:topological-tools} we use the Riesz interface
\Cref{thm:riesz-interface} and its uniqueness clause
\Cref{cor:riesz-uniqueness}, weak convergence and the portmanteau theorem
\Cref{prop:portmanteau}, the joint convergence lemma
\Cref{lem:joint-weak-uniform}, and the concave-conjugate formalism of
\Cref{subsec:legendre-finite}, in particular
\eqref{eq:dual-energy-abstract}, \Cref{lem:dual-mass},
\eqref{eq:fenchel-energy} and \Cref{prop:concave-biconjugation}.  From
\Cref{chap:bfj-compactness} we use skeletal equicontinuity
\Cref{thm:bfj-skeletal-equicontinuity} with its explicit modulus
\Cref{cor:ch22-explicit-modulus}, compactness of $\PSH_0$
\Cref{thm:bfj-psh-compactness} in the net form
\Cref{cor:ch22-net-compactness}, and decreasing regularization
\Cref{thm:bfj-decreasing-regularization}.  From \Cref{chap:bfj-energy} we
use the energy cocycle \eqref{eq:energy-cocycle}, the extension
\Cref{thm:extended-energy-basic}, the class $\cE^1$ of \Cref{def:e1}, the
mass tail \eqref{eq:e1-mass-tail} and the integrability bound
\eqref{eq:e1-integrability}, the mixed-energy identities
\Cref{thm:ch26-mixed-energy-identities}, the finite-energy sandwich
\eqref{eq:ch26-e1-sandwich}, monotone continuity
\Cref{prop:e1-monotone-continuity}, and the compactness of normalized
energy superlevels \Cref{prop:e1-compact-superlevels} in the net form
\Cref{cor:ch26-net-superlevel}.  From
\Cref{chap:bfj-differentiability} we use exactly two statements, the
finite-energy projected variation \Cref{cor:e1-projected-variation} and the
Euler--Lagrange proposition \Cref{prop:bfj-euler-lagrange}.

\section{Orientation}
\label{sec:ch28-orientation}

\subsection{The inverse problem, and why a functional appears}
\label{sec:ch28-inverse}
\index{variational method!BFJ existence}

Every construction of \Cref{chap:bfj-ma} and \Cref{chap:bfj-energy} runs in
one direction: from a potential to a measure.  A model potential determines
a vertical divisor, hence $n$-fold intersection numbers on a model, hence
an atomic measure supported at divisorial points
(\Cref{def:mixed-model-measure}); monotone limits extend this to bounded and
then to finite-energy potentials.  The inverse problem asks for the
preimage.  It has no direct answer, because there is no formula: the
model operator is polynomial of degree at most $n$ in the
potential.  For $n\geq2$ its equations can be nonlinear, and positivity
must also be checked on all vertical curves.  A finite collection of vertex
masses alone does not describe the general nef constraint.

The variational method replaces inversion by optimization.  Because
$E_\omega$ is concave with derivative $\MA_\omega$, the equation
$\MA_\omega(\varphi)=\mu$ says exactly that the affine functional
$\psi\mapsto\int\psi\,\rd\mu$ supports the concave functional
$E_\omega$ at $\varphi$; in the language of
\Cref{prop:concave-biconjugation}\,(3), $\mu$ is a supergradient of
$E_\omega$ at $\varphi$.  That is a statement about a maximum, namely
that
\[
 F_\mu(\psi):=E_\omega(\psi)-\int\psi\,\rd\mu
\]
is maximal at $\psi=\varphi$; and this reformulation, unlike the equation
itself, is stable under the only limiting operations the theory provides.
\Cref{cor:ch27-euler-lagrange-supergradient} already recorded the
equivalence.  What remains, and what this chapter supplies, is the analysis
of the maximization problem: for which $\mu$ is $F_\mu$ bounded above,
for which $\mu$ is its supremum attained, and how the maximizer moves when
$\mu$ does.

One dimension of the problem is visible already on a curve, and it is worth
recording at the outset because the fixed-skeleton problem then has a
complete elementary answer.  For $n=1$ the operator is affine on a fixed
skeleton: by \eqref{eq:curve-ma-laplacian},
$\MA_\omega(u)=\mu_0+\Delta_\Gamma u$ on a skeleton, so the equation is
the linear Poisson equation \eqref{eq:curve-poisson}, the functional
$F_\mu$ is a strictly concave quadratic form, and the dual energy is a
classical Dirichlet integral.  \Cref{sec:ch28-curves} carries this out with
numbers.  In dimension $n\geq2$ nothing of the sort is available, and each
of the three questions above becomes a separate theorem.

\subsection{Three theorems that must not be merged}
\label{sec:ch28-three}

\begin{center}
\small
\begin{tabular}{p{.20\textwidth}|p{.36\textwidth}|p{.36\textwidth}}
\toprule
\textbf{Name}&\textbf{Statement}&\textbf{Status here}\\
\midrule
Duality
 &$\sup F_\mu=E^\vee_\omega(\mu)$, and equality in the Fenchel
  inequality identifies maximizers, and under orthogonality solutions
 &proved:
  \Cref{prop:ch28-sup-is-dual-energy},
  \Cref{cor:finite-energy-variational-characterization}\\
Completion
 &model Monge--Amp\`ere measures are strongly dense in
  $\mathcal M^1_V$; model maximizing sequences exist for every
  $\mu\in\mathcal M^1_V$
 &quoted with its source:
  \Cref{thm:finite-energy-measure-completion}; the maximizing-sequence density assertion is
  proved here in \Cref{thm:finite-energy-measure-stability}\\
Existence
 &the supremum of $F_\mu$ is \emph{attained} in $\cE^1$, and the
  maximizer solves $\MA_\omega(\varphi)=\mu$
 &proved for all $\mu\in\mathcal M^1_V$ under orthogonality:
  \Cref{thm:ch28-finite-energy-existence};
  supported regularity in \Cref{thm:bfj-supported-cy}\\
\bottomrule
\end{tabular}
\end{center}

Passing from the second row to the third requires control of the
energy and the pairing, beyond compactness of $\PSH_0$ alone.
\Cref{sec:ch28-coercivity} gives a short argument for admissible
measures, and \Cref{sec:ch28-general-finite-energy-existence} supplies
the energy-level argument for every finite-energy measure.

\begin{warning}{``Finite energy'' is used of two different objects}
A \emph{finite-energy potential} is an element of $\cE^1(X,\omega)$, that
is, a psh function whose Monge--Amp\`ere energy is finite
(\Cref{def:e1}).  A \emph{finite-energy measure} is an element of
$\mathcal M^1_V(X,\omega)$, that is, a positive measure of mass $V$
whose Legendre dual energy is finite
(\Cref{def:dual-energy-measure}).  The two conditions are Legendre-dual to
one another, but they are not interchangeable, and neither implies a
statement about the other without a theorem.  In particular:
$\MA_\omega(\varphi)\in\mathcal M^1_V$ for every $\varphi\in\cE^1$
(\Cref{prop:ch28-dual-energy-of-ma}), but the converse inclusion
$\mathcal M^1_V\subseteq\MA_\omega(\cE^1)$ requires an existence
theorem.  Under the standing hypotheses and orthogonality it follows from
\Cref{thm:ch28-finite-energy-existence}, not from the definitions.
\end{warning}

\subsection{The hypothesis signature, and the later extensions}
\label{sec:ch28-scope}

\index{hypothesis signature!BFJ existence theorem}
The following table is the scope discipline of the chapter in one place.
The first column is the axis, the second the hypothesis actually used in
\Cref{thm:bfj-supported-cy}, the third the place in the proof where it is
used, and the fourth the later results that weaken it.  A later result is
never silently substituted for the original one; each is stated separately,
with its date, its source, and its own additional hypotheses.

\begin{center}
\small
\begin{tabular}{p{.15\textwidth}|p{.22\textwidth}|p{.28\textwidth}|p{.26\textwidth}}
\toprule
\textbf{Axis}&\textbf{Hypothesis used}&\textbf{Where it enters}
 &\textbf{Later extensions}\\
\midrule
base field
 &complete, discretely valued
 &cofinality of SNC models \Cref{prop:snc-cofinal}; the whole model
  calculus of \Cref{chap:bfj-model-forms}
 &ample curve metric envelopes and skeletal-supported continuous
  solutions over arbitrary complete nontrivially valued fields,
  \Cref{thm:curves-arbitrary-field-continuity}; the general case is
  \emph{open}, \Cref{prob:frontier-general-ma}\\
residue char.
 &zero
 &resolution of singularities; multiplier ideals in
  \Cref{thm:bfj-envelope-regularity}
 &equal characteristic $p$ with geometric origin,
  \Cref{thm:gjkm-positive-characteristic-cy}; mixed characteristic,
  \Cref{thm:fgk-mixed-characteristic-cy}\\
space
 &$X$ smooth projective
 &SNC models; skeletal equicontinuity
  \Cref{thm:bfj-skeletal-equicontinuity}
 &none in \Cref{part:bfj}; the unibranch obstruction in the trivially
  valued theory is \Cref{chap:bj-strong-topologies}\\
class
 &$\{\omega\}$ ample, $\omega$ semipositive
 &positivity of the model measure \Cref{prop:model-ma-basic}; the Hodge
  index input \Cref{thm:ch20-mixed-hodge}
 &big classes are not treated here\\
orthogonality
 &assumed
 &\Cref{prop:variational-maximizer-solves} only
 &algebraizability, \cite[Appendix~A]{BFJ15}; removed in residue
  characteristic zero, \Cref{thm:bggjkm-remove-algebraizability}\\
measure
 &mass $V$, supported on one SNC dual complex
 &coercivity \Cref{lem:ch28-skeleton-admissible}; Dini in
  \Cref{thm:bfj-supported-regularity}
 &all finite-energy measures for finite-energy existence,
  \Cref{thm:ch28-finite-energy-existence}; continuity still uses support\\
\bottomrule
\end{tabular}
\end{center}

Two entries deserve emphasis, because they are the ones most often blurred
in summaries of the theory.

First, \emph{algebraizability is not a hypothesis of
\Cref{thm:bfj-supported-cy}}.  The historical Theorem~A of \cite{BFJ15}
assumed that $X$ descends from a smooth projective variety over a
one-variable function field whose completion is $K$; that assumption was
used only to prove orthogonality, in \cite[Appendix~A]{BFJ15}.  Our theorem
assumes orthogonality outright and is therefore insensitive to the route by
which it is obtained.  In residue characteristic zero, the differentiability
theorem for non-Archimedean volumes removes algebraizability altogether
\cite[Theorem~6.3.3 and Corollary~6.3.4]{BGGJKM20}; this is
\Cref{thm:bggjkm-remove-algebraizability}.

Second, \emph{discreteness is not removed by any theorem proved in this
book, except on curves}.  The class $\cE^1$ admits a definition over an
arbitrary complete non-Archimedean field, but the Monge--Amp\`ere operator
is constructed in the discretely valued residue-characteristic-zero
case treated here and in the trivially valued case of \Cref{part:bj}.
The quoted positive- and mixed-characteristic results have their own
additional resolution hypotheses.  A general extension is stated as
unresolved in
\cite[Remark~4.4]{PS23Hybrid} and is recorded in this book as
\Cref{prob:frontier-general-ma}.  Nothing in this chapter should be read as
asserting the general extension, and in particular the sentence ``the
finite-energy theory works over any complete field'' is not a theorem.

\subsection{What is proved here and what is imported}
\label{sec:ch28-imported}

The arguments below use
\Cref{chap:bfj-model-forms,chap:bfj-psh,chap:bfj-compactness,chap:bfj-envelopes,chap:bfj-ma,chap:bfj-capacity,chap:bfj-energy,chap:bfj-differentiability},
together with the following source inputs, flagged at their points of use.

\begin{enumerate}
\item[\textup{(I1)}]
\emph{Orthogonality} in dimension $n\geq2$, quoted with its two
provenances in \Cref{thm:orthogonality-provenance}.  Its Euler--Lagrange use is through
\Cref{cor:e1-projected-variation} and \Cref{prop:bfj-euler-lagrange}, that
is, in \Cref{prop:variational-maximizer-solves} and its consequences.
The general finite-energy argument also uses the quantitative
orthogonality estimates through input \textup{(I3)}.
\item[\textup{(I2)}]
\emph{Uniqueness} of the finite-energy solution modulo constants.  What is
proved here is the structure of the solution set,
\Cref{prop:ch28-solution-set}, which has an elementary continuous-curve uniqueness consequence
(\Cref{prop:ch28-curve-uniqueness}).  Finite-energy uniqueness in all
dimensions is the non-Archimedean Calabi--Yau uniqueness theorem, quoted as
\cite[Proposition~8.2]{BFJ15} and recorded in
\Cref{cor:bfj-original-continuous-solution}.  It is \emph{not} a consequence
of the comparison principle, for the reason given in
\Cref{rem:ch25-no-uniqueness}.
\item[\textup{(I3)}]
The \emph{strong completion} theorem for finite-energy measures,
\Cref{thm:finite-energy-measure-completion}, quoted from
\cite{BJ25FiniteEnergy}.  It is used nowhere in the proof of
\Cref{thm:bfj-supported-cy}; it is recorded because it supplies strong density and the quantitative pairing estimates used
in \Cref{thm:ch28-finite-energy-existence}.  The density argument is
expanded in \Cref{thm:finite-energy-measure-stability}.
\item[\textup{(I4)}]
The divisorial Green-envelope identification
\cite[Lemma~8.5 and \textup{(8.7)}]{BFJ15}, and the non-$\mathbf Q$-PL divisorial examples and the non-retraction-invariant
finite-atomic example of \cite{BJGreen24}.
\end{enumerate}

Nothing in this chapter is used in
\Cref{chap:bfj-envelopes,chap:bfj-ma,chap:bfj-capacity,chap:bfj-energy,chap:bfj-differentiability},
which is what makes the dependency order
\eqref{eq:ch27-dependency-diagram} of \Cref{sec:ch27-dependency} consistent:
the present chapter is step \textup{(D11)} of that order and nothing else
depends on it inside \Cref{part:bfj} except
\Cref{chap:bfj-continuity,chap:bfj-examples}.

\subsection{Plan}
\label{sec:ch28-plan}

\Cref{sec:ch28-dual-energy} defines the dual energy, proves that its
definition does not depend on the test class used, computes it at a
Monge--Amp\`ere measure, and identifies the Fenchel defect with the
functional $J$ of \Cref{def:ch26-ij}.
\Cref{sec:ch28-variational} sets up the variational functional, proves that
its supremum is the dual energy, and proves that a maximizer solves the
equation.  \Cref{sec:ch28-coercivity} isolates coercivity, proves the
existence theorem for admissible measures with the compactness argument in
net form, and gives the counterexamples that show the hypotheses are not
decorative.  \Cref{sec:ch28-strong-measures} records the strong topology on
measures and the completion theorem, and separates them carefully from the
gauge on potentials of \Cref{sec:ch26-e1-gauge}.
\Cref{sec:ch28-general-finite-energy-existence} combines the quantitative
measure estimates with compactness on an energy level to prove existence
for every finite-energy measure in the standing BFJ framework, under
orthogonality.
\Cref{sec:ch28-bfj-theorem} states the original BFJ existence theorem with
its full hypothesis signature and proves it.  \Cref{sec:ch28-calabi-yau}
assembles the Calabi--Yau correspondence.  \Cref{sec:ch28-stability} proves
stability with respect to the measure, in three forms.
\Cref{sec:ch28-curves} carries everything to explicit numbers on curves, and
\Cref{sec:ch28-dictionary} is the chapter dictionary.

\section{The Legendre dual energy}
\label{sec:ch28-dual-energy}

\subsection{Definition and the forced mass condition}
\label{sec:ch28-dual-def}

\begin{definition}[Energy of a measure]
\label{def:dual-energy-measure}
\index{dual energy!of a measure}\index{Legendre transform!dual energy}
For $\mu\in\mathcal M_V(X)$, set
\begin{equation}
\label{eq:dual-energy-measure}
 E_\omega^\vee(\mu)
 :=\sup_{\varphi\ \mathrm{psh\ model}}
       \left\{E_\omega(\varphi)-\int\varphi\,\rd \mu\right\}
 \in[0,+\infty].
\end{equation}
This is a Legendre transform on the psh test/model locus; no integral of an
arbitrary unbounded potential is used in the definition.  We call $\mu$ a
\emph{finite-energy measure} when
$E_\omega^\vee(\mu)<+\infty$, and write
$\mathcal M^1_V(X,\omega)$ for the resulting class.
\end{definition}

Three features of \eqref{eq:dual-energy-measure} deserve comment before any
property is proved.

\emph{The test class is small.}  The supremum is taken over
$\omega$-psh \emph{model} functions.  These are continuous and bounded on the compact analytification, so
$\int\varphi\,\rd\mu$ is a finite integral of a bounded continuous function
against a finite measure, and no convention about $-\infty$ is needed.
This is deliberate: it makes $E^\vee_\omega$ manifestly a supremum of
weak-continuous affine functionals of $\mu$, which is what
\Cref{prop:dual-energy-elementary} exploits.
\Cref{prop:ch28-test-class} shows that enlarging the test class to bounded,
continuous, or finite-energy potentials does not change the value; the small
class is used in the definition and the large one in the applications.

\emph{The sign convention is the concave one.}  $E_\omega$ is concave and
is to be maximized, so \eqref{eq:dual-energy-measure} is the concave
conjugate \eqref{eq:concave-conjugate} of $E_\omega$ with respect to the
pairing $\varphi\mapsto\int\varphi\,\rd\mu$, and
$E^\vee_\omega$ takes values in $[0,+\infty]$ rather than
$[-\infty,0]$.  \Cref{app:normalizations} records this at
\eqref{eq:appJ-dual-energy}.

\emph{The mass is not normalized.}  For probability measures the usual
convention divides both the energy and the Monge--Amp\`ere measure by $V$.
Formula \eqref{eq:dual-energy-measure} keeps the intersection-theoretic mass
and is therefore compatible with the rest of this book; the dictionary
between the two conventions is
\begin{equation}
\label{eq:ch28-probability-dictionary}
 E_\omega^{\vee,\mathrm{prob}}(V^{-1}\mu)=V^{-1}E_\omega^\vee(\mu),
\end{equation}
which follows by dividing both terms of \eqref{eq:dual-energy-measure} by
$V$.  Mixing a probability measure with the unnormalized energy destroys
translation invariance and is therefore immediately detectable, as
\eqref{eq:appJ-dual-energy} warns.

\begin{lemma}
\label{lem:ch28-mass-forced}
Let $\mu$ be a positive Radon measure on $X^{\An}$ and define
$E^\vee_\omega(\mu)$ by \eqref{eq:dual-energy-measure}.  Then
$E^\vee_\omega(\mu)<+\infty$ forces $\mu(X^{\An})=V$.  Moreover
$E^\vee_\omega(\mu)\geq0$ for every $\mu\in\mathcal M_V(X)$, with
equality if and only if the constant $0$ maximizes $F_\mu$ among model
psh functions.
\end{lemma}

\begin{proof}
Constants are $\omega$-psh model functions, so they are admissible test
objects.  By the translation identity \eqref{eq:bfj-energy-translation},
\begin{equation}
\label{eq:dual-energy-constant-test}
 E_\omega(\varphi+c)-\int(\varphi+c)\,\rd \mu
 =E_\omega(\varphi)-\int\varphi\,\rd \mu
   +c\bigl(V-\mu(X^{\An})\bigr)
\end{equation}
for every model psh $\varphi$ and every $c\in\mathbf R$.  If
$\mu(X^{\An})\neq V$, let $c$ run to $+\infty$ or to $-\infty$
according to the sign of $V-\mu(X^{\An})$; the right-hand side is
unbounded above, so $E^\vee_\omega(\mu)=+\infty$.  This is
\Cref{lem:dual-mass} in the present notation, with $M=V$ and
$\varphi_0=0$, for which $E_\omega(0)=0$ is finite.  Testing
\eqref{eq:dual-energy-measure} against $\varphi=0$ gives
$E^\vee_\omega(\mu)\geq E_\omega(0)-0=0$, and equality means precisely that
no model psh competitor does better than $0$.
\end{proof}

\begin{warning}{A positive measure of the wrong mass has no dual energy}
\Cref{lem:ch28-mass-forced} is the reason why $\mathcal M^1_V$ is defined
inside $\mathcal M_V$ rather than inside the cone of all positive
measures.  The value $E^\vee_\omega(\mu)=+\infty$ at the wrong mass is not
an artefact of a bad normalization; it records that the variational problem
has no solution, since $\MA_\omega(\varphi)$ has mass $V$ for every
$\varphi\in\cE^1$ by \Cref{prop:e1-integrability-tail}.  Consequently the
words \emph{finite-energy measure} always include the mass condition.
\end{warning}

\subsection{Elementary properties}
\label{sec:ch28-dual-elementary}

\begin{proposition}
\label{prop:dual-energy-elementary}
The functional $E_\omega^\vee$ is convex and weakly lower
semicontinuous on $\mathcal M_V(X)$.  For every psh model function $u$,
\begin{equation}
\label{eq:measure-fenchel-energy}
       \int u\,\rd \mu
       \geq E_\omega(u)-E_\omega^\vee(\mu).
\end{equation}
\end{proposition}

\begin{proof}
For fixed model $u$, the expression
$E_\omega(u)-\int u\,\rd \mu$ is continuous and affine in $\mu$.
Its supremum over model psh functions is therefore convex and lower
semicontinuous.  Inequality \eqref{eq:measure-fenchel-energy} is the
definition rearranged.
\end{proof}

We expand the two assertions of the proof, since both are used repeatedly.
Affineness in $\mu$ is clear, the first term being independent of $\mu$
and the second linear.  Weak continuity in $\mu$ holds because $u$ is a
model function, hence continuous on the compact space $X^{\An}$, so
$\mu\mapsto\int u\,\rd\mu$ is by definition continuous for the weak
topology of \Cref{prop:portmanteau}.  A pointwise supremum of a family of
continuous affine functions is convex and lower semicontinuous, by
\Cref{prop:net-semicontinuity} applied to each member of the family; this is
the same mechanism as in \Cref{prop:concave-biconjugation}\,(1), where the
conjugate of a concave function is convex and lower semicontinuous for the
same reason.  Nothing about $E_\omega$ is used, and in particular no
concavity, monotonicity or semicontinuity of the energy enters.

\begin{proposition}
\label{prop:ch28-dual-order}
Let $\mu,\nu\in\mathcal M_V(X)$ and $s\in[0,1]$.
\begin{enumerate}
\item\label{itm:ch28-dual:convex}
 $E^\vee_\omega\bigl((1-s)\mu+s\nu\bigr)\leq
 (1-s)E^\vee_\omega(\mu)+sE^\vee_\omega(\nu)$; in particular
 $\mathcal M^1_V(X,\omega)$ is a convex subset of $\mathcal M_V(X)$.
\item\label{itm:ch28-dual:sublevel}
 For every $C\geq0$ the set
 $\{\mu\in\mathcal M_V(X):E^\vee_\omega(\mu)\leq C\}$ is convex and weakly
 compact.
\item\label{itm:ch28-dual:reference}
 If $f$ is an $\omega$-psh model function and $\omega'=\omega+\ddc f$,
 then $E^{\vee}_{\omega'}(\mu)=E^\vee_\omega(\mu)-E_\omega(f)
 +\int f\,\rd\mu$ whenever either side is finite.
\end{enumerate}
\end{proposition}

\begin{proof}
\cref{itm:ch28-dual:convex} is \Cref{prop:dual-energy-elementary}; convexity
of $\mathcal M^1_V$ follows because the right-hand side is finite when
both dual energies are.

\cref{itm:ch28-dual:sublevel}.  The set $\mathcal M_V(X)$ is weakly compact:
it is a weak-$*$ closed subset of the ball of radius $V$ in the dual of
$C^0(X^{\An})$, and $X^{\An}$ is compact by
\Cref{thm:berk-proper-compact}, so the Banach--Alaoglu theorem quoted
in \Cref{thm:riesz-interface} applies; closedness is the statement that
$\mu\mapsto\mu(X^{\An})=\int 1\,\rd\mu$ is weakly continuous, the constant
$1$ being a legitimate test function on a compact space.  A sublevel set of
a lower semicontinuous convex function on a compact convex set is closed and
convex, hence compact.

\cref{itm:ch28-dual:reference}.  By the change-of-reference formula
\eqref{eq:ch26-reference-change}, an $\omega'$-psh model function is
exactly $\varphi-f$ for $\varphi$ an $\omega$-psh model function, and
$E_{\omega'}(\varphi-f)=E_\omega(\varphi)-E_\omega(f)$.  Substituting into
\eqref{eq:dual-energy-measure} and using
$\int(\varphi-f)\,\rd\mu=\int\varphi\,\rd\mu-\int f\,\rd\mu$ gives the
stated identity, the two correction terms being finite real numbers
independent of $\varphi$.
\end{proof}

\begin{remark}[Weak compactness is not strong compactness]
\label{rem:ch28-weak-not-strong}
\Cref{prop:ch28-dual-order}\,\cref{itm:ch28-dual:sublevel} is a genuine
compactness statement, and it is the reason the measure side of the theory
is better behaved than the potential side: the corresponding superlevel sets
of the energy on $\cE^1$ are weakly compact
(\Cref{prop:e1-compact-superlevels}) and need not be strongly compact
(\Cref{prop:ch26-no-strong-compactness}).  The measure-side sublevels
likewise need not be strongly compact.  In the projective-line example of
that proposition, the measures
$\nu_N=\MA_\omega(\tau_N^\varepsilon)$ converge weakly to
$\mu_0=\delta_\eta$: their atoms lie in pairwise distinct branches,
and continuity at $\eta$ gives convergence against every continuous test.
Their dual energies equal $\varepsilon/2>0$, while
$E_\omega^\vee(\mu_0)=0$.  Choosing $\varepsilon/2\leq C$ disproves
strong compactness of that sublevel for every $C>0$ in this example.
The quoted measure-side completeness is a different property.
\end{remark}

\subsection{The test class does not matter}
\label{sec:ch28-test-class}

The definition uses model psh functions.  In every application one wants to
test against a bounded, a continuous, or an arbitrary finite-energy
potential.  The following proposition says that all four suprema agree.  Its
proof is the first place where decreasing regularization
\textup{(D3)} is used.

\begin{proposition}
\label{prop:ch28-test-class}
Let $\mu\in\mathcal M_V(X)$.  Then
\begin{equation}
\label{eq:ch28-test-class}
 E^\vee_\omega(\mu)
 =\sup_{\varphi\in\mathcal T}
   \Bigl\{E_\omega(\varphi)-\int\varphi\,\rd\mu\Bigr\}
\end{equation}
for each of the following test classes $\mathcal T$:
\begin{enumerate}
\item\label{itm:ch28-test:model}
 the $\omega$-psh model functions;
\item\label{itm:ch28-test:cont}
 the continuous $\omega$-psh functions;
\item\label{itm:ch28-test:bdd}
 the bounded $\omega$-psh functions;
\item\label{itm:ch28-test:e1}
 the $\varphi\in\cE^1(X,\omega)$ with $\int\varphi\,\rd\mu>-\infty$.
\end{enumerate}
In \cref{itm:ch28-test:e1} the test class is nonempty, since it contains the
constants.
\end{proposition}

\begin{proof}
Write $S_1,\ldots,S_4$ for the four suprema.  Model psh functions are
continuous and bounded and lie in $\cE^1$ by
\Cref{thm:ch26-e1-structure}\,\cref{itm:ch26-e1:bounded}, and for such
$\varphi$ the integral $\int\varphi\,\rd\mu$ is finite; hence
$S_1\leq S_2\leq S_3\leq S_4$, the last two inclusions of test classes
being immediate.  It therefore suffices to prove $S_4\leq S_1$.

Let $\varphi\in\cE^1(X,\omega)$ with $\int\varphi\,\rd\mu>-\infty$, and
let $\varepsilon>0$.  By \Cref{thm:bfj-decreasing-regularization} there is
a decreasing net $(\varphi_\alpha)$ of $\omega$-psh model functions with
$\varphi_\alpha\searrow\varphi$.  Each $\varphi_\alpha$ is an admissible
test object for $S_1$, so
\[
 S_1\ \geq\ E_\omega(\varphi_\alpha)-\int\varphi_\alpha\,\rd\mu
 \qquad\text{for every }\alpha .
\]
We pass to the limit in each term separately, both limits existing in
$\mathbf R$.

For the energy, \Cref{thm:extended-energy-basic} gives continuity of
$E_\omega$ along decreasing nets, so
$E_\omega(\varphi_\alpha)\to E_\omega(\varphi)$, a finite number because
$\varphi\in\cE^1$.

For the integral, the net $(\varphi_\alpha)$ is decreasing with pointwise
infimum $\varphi$ --- the regularization theorem produces pointwise
convergence, not merely convergence on dual complexes --- and on the cofinal tail
$\alpha\geq\alpha_0$ every member is bounded above by the bounded
function $\varphi_{\alpha_0}$.  The monotone convergence theorem
for a decreasing net of upper semicontinuous functions against a fixed finite
measure, \Cref{cor:usc-net-monotone-integral}, therefore gives
\[
 \int\varphi_\alpha\,\rd\mu\ \longrightarrow\ \int\varphi\,\rd\mu ,
\]
the limit being finite by hypothesis.  Note that this is a statement about a
\emph{fixed} measure and a monotone net of integrands; the portmanteau
inequality \Cref{prop:portmanteau} would give only an inequality in the
wrong direction, and \Cref{rem:no-reverse-inequality} explains why no
reverse inequality is available in general.

Combining the two limits, $S_1\geq E_\omega(\varphi)-\int\varphi\,\rd\mu$.
Taking the supremum over admissible $\varphi$ gives $S_1\geq S_4$.
\end{proof}

\begin{remark}[Where the enlargement is used, and where it is not]
\label{rem:ch28-test-class-use}
The enlarged test class is used to identify the supremum of $F_\mu$
and its value at a solution.  The original model class already gives
a direct proof of convexity and weak lower semicontinuity of dual
energy, since each test is bounded and continuous.  Alternatively,
for a fixed $\cE^1$ potential $\varphi$, upper semicontinuity and its
global upper bound imply that
$\mu\mapsto E_\omega(\varphi)-\int\varphi\,\rd\mu$ is weakly lower
semicontinuous, allowing the value $+\infty$ when the integral is
$-\infty$.  The minus sign is essential.  The small test class avoids
this extended-value convention in the definition.
\end{remark}

\subsection{Dual energy at a Monge--Amp\`ere measure}
\label{sec:ch28-dual-at-ma}

The next proposition is the basic computation of the chapter.  It says that
the dual energy of a measure that already \emph{is} a Monge--Amp\`ere
measure can be written down, and that the supremum defining it is attained.
It uses only the finite-energy sandwich \eqref{eq:ch26-e1-sandwich}, that is
step \textup{(D6)} of \Cref{sec:ch27-dependency}, and in particular it does
not use orthogonality.

\begin{proposition}
\label{prop:ch28-dual-energy-of-ma}
Let $\varphi\in\cE^1(X,\omega)$ and put $\mu:=\MA_\omega(\varphi)$.  Then
$\mu\in\mathcal M^1_V(X,\omega)$, the supremum in
\eqref{eq:ch28-test-class} is attained at $\varphi$, and
\begin{equation}
\label{eq:ch28-dual-energy-of-ma}
 E^\vee_\omega\bigl(\MA_\omega(\varphi)\bigr)
 =E_\omega(\varphi)-\int\varphi\,\MA_\omega(\varphi).
\end{equation}
If moreover $\sup_{X^{\An}}\varphi=0$, then
\begin{equation}
\label{eq:ch28-dual-energy-bracket}
 0\ \leq\ E^\vee_\omega\bigl(\MA_\omega(\varphi)\bigr)
 \ \leq\ n\bigl(-E_\omega(\varphi)\bigr).
\end{equation}
\end{proposition}

\begin{proof}
\emph{Step 1: the right-hand side of \eqref{eq:ch28-dual-energy-of-ma} is a
finite real number.}  Normalize $\varphi\leq0$ by subtracting a constant;
by \eqref{eq:dual-energy-constant-test} with $\mu(X^{\An})=V$, which holds
by \Cref{prop:e1-integrability-tail}, neither side of
\eqref{eq:ch28-dual-energy-of-ma} changes.  Then $E_\omega(\varphi)$ is
finite by definition of $\cE^1$ and
$0\leq\int(-\varphi)\MA_\omega(\varphi)\leq(n+1)(-E_\omega(\varphi))$ by
\eqref{eq:e1-integrability}.

\emph{Step 2: every competitor is dominated.}  Let
$\psi\in\cE^1(X,\omega)$ with $\int\psi\,\rd\mu>-\infty$.  The
right-hand inequality of the two-sided sandwich
\eqref{eq:ch27-two-sided-sandwich}, applied to the pair
$(\varphi,\psi)$ with no ordering hypothesis, gives
\[
 E_\omega(\psi)-E_\omega(\varphi)
 \ \leq\ \int(\psi-\varphi)\,\MA_\omega(\varphi)
 =\int\psi\,\rd\mu-\int\varphi\,\rd\mu ,
\]
all four integrals being finite by
\Cref{lem:ch27-mutual-integrability} and Step~1.  Rearranging,
\[
 E_\omega(\psi)-\int\psi\,\rd\mu
 \ \leq\ E_\omega(\varphi)-\int\varphi\,\rd\mu .
\]

\emph{Step 3: conclusion.}  By Step~2 the supremum over the test class
\cref{itm:ch28-test:e1} of \Cref{prop:ch28-test-class} is attained at
$\varphi$ and equals the right-hand side of
\eqref{eq:ch28-dual-energy-of-ma}, which is finite; by
\Cref{prop:ch28-test-class} that supremum equals $E^\vee_\omega(\mu)$.
Hence $\mu\in\mathcal M^1_V$ and \eqref{eq:ch28-dual-energy-of-ma} holds.

\emph{Step 4: the bracket.}  Assume $\sup\varphi=0$.  The lower bound is
\Cref{lem:ch28-mass-forced}.  For the upper bound, combine
\eqref{eq:ch28-dual-energy-of-ma} with the right-hand inequality of
\eqref{eq:e1-integrability}:
\[
 E^\vee_\omega(\mu)
 =E_\omega(\varphi)+\int(-\varphi)\MA_\omega(\varphi)
 \leq E_\omega(\varphi)+(n+1)\bigl(-E_\omega(\varphi)\bigr)
 =n\bigl(-E_\omega(\varphi)\bigr).
\]
This is \eqref{eq:ch28-dual-energy-bracket}.
\end{proof}

\begin{remark}[Reading \eqref{eq:ch28-dual-energy-of-ma}]
\label{rem:ch28-reading-dual-of-ma}
Three consequences are worth isolating.

First, $\MA_\omega$ maps $\cE^1(X,\omega)$ into
$\mathcal M^1_V(X,\omega)$.  This is the easy half of the Calabi--Yau
correspondence and is proved unconditionally; the hard half, surjectivity,
is proved under orthogonality in
\Cref{thm:ch28-finite-energy-existence}.

Second, \eqref{eq:ch28-dual-energy-of-ma} is a formula for the dual energy,
not merely a bound, and it is the formula one computes with.  On a curve it
becomes the classical Dirichlet integral of the Green potential of $\mu$,
\Cref{prop:ch28-curve-dual-energy}, and the computation there is completely
explicit.

Third, \eqref{eq:ch28-dual-energy-bracket} bounds the dual energy of the
image by the energy of the source, with the dimensional constant $n$.  The
upper bound is attained for $n=1$, by
\Cref{ex:ch28-unbounded-dual-energy}, so the dimensional bound is optimal when $n=1$.
No optimality assertion for each higher dimension follows from that example.  The lower bound $0$ cannot be replaced by any positive multiple
of $-E_\omega(\varphi)$: the affine family of
\Cref{ex:ch28-affine-family} has
\[
 \frac{E^\vee_\omega(\MA_\omega(u_\varepsilon))}
      {-E_\omega(u_\varepsilon)}
 =\frac{\varepsilon^2/2}{\varepsilon+\varepsilon^2/2}
 \ \xrightarrow[\varepsilon\downarrow0]{}\ 0 .
\]
The two sides of \eqref{eq:ch28-dual-energy-bracket} are therefore genuinely
different quantities, and no equivalence of the two ``energies'' is
available.
\end{remark}

\subsection{The Fenchel defect and the functional \texorpdfstring{$J$}{J}}
\label{sec:ch28-fenchel-defect}

The Fenchel inequality \eqref{eq:measure-fenchel-energy} has a defect, and
that defect is an object already named in \Cref{chap:bfj-energy}.  Recording
the identification saves a great deal of later bookkeeping and explains the
shape of \eqref{eq:measure-dirichlet-quasimetric}.

For psh model functions $\varphi,\psi$, introduce the directed Dirichlet
functional
\begin{equation}
\label{eq:directed-dirichlet-functional}
 J_\omega(\varphi,\psi)
 :=E_\omega(\varphi)-E_\omega(\psi)
   +\int(\psi-\varphi)\MA_\omega(\varphi)\geq0.
\end{equation}

\begin{proposition}
\label{prop:ch28-fenchel-defect}
\index{Fenchel inequality!defect}
Let $\varphi,\psi\in\cE^1(X,\omega)$ and put
$\mu=\MA_\omega(\varphi)$.  Then all the quantities below are finite,
\begin{equation}
\label{eq:ch28-defect-identity}
 J_\omega(\varphi,\psi)
 =E^\vee_\omega(\mu)-\Bigl(E_\omega(\psi)-\int\psi\,\rd\mu\Bigr)
 \ \geq\ 0 ,
\end{equation}
and $J_\omega(\varphi,\psi)=J_\varphi(\psi)$ in the notation of
\Cref{def:ch26-ij}.  If $\varphi$ and $\psi$ are bounded, then
\begin{equation}
\label{eq:ch28-defect-ak}
 \begin{aligned}
 J_\omega(\varphi,\psi)
 &=\frac1{n+1}\sum_{k=0}^{n-1}(k+1)\,a_k(\varphi,\psi),\\
 \frac1{n+1}I(\varphi,\psi)
 &\leq J_\omega(\varphi,\psi)
 \leq\frac n{n+1}I(\varphi,\psi),
 \end{aligned}
\end{equation}
with $a_k$ as in \eqref{eq:ch26-ak} and $I$ as in
\eqref{eq:ch26-I-definition}.  In particular
$J_\omega(\varphi,\psi)=0$ if and only if $I(\varphi,\psi)=0$, if and
only if $a_k(\varphi,\psi)=0$ for every $k$.
\end{proposition}

\begin{proof}
Finiteness is \Cref{lem:ch27-mutual-integrability} together with
$\varphi,\psi\in\cE^1$.  Nonnegativity in
\eqref{eq:directed-dirichlet-functional} is the right-hand inequality of
\eqref{eq:ch27-two-sided-sandwich}, which says
$E_\omega(\psi)-E_\omega(\varphi)\leq\int(\psi-\varphi)\MA_\omega(\varphi)$.

For \eqref{eq:ch28-defect-identity}, substitute
\eqref{eq:ch28-dual-energy-of-ma}:
\[
 E^\vee_\omega(\mu)-E_\omega(\psi)+\int\psi\,\rd\mu
 =E_\omega(\varphi)-\int\varphi\,\rd\mu-E_\omega(\psi)
  +\int\psi\,\rd\mu ,
\]
which is exactly \eqref{eq:directed-dirichlet-functional} once
$\int(\psi-\varphi)\,\rd\mu$ is written as a difference of two finite
integrals.  Comparing \eqref{eq:directed-dirichlet-functional} with
\eqref{eq:ch26-J-definition}, in which the roles of the two arguments are
interchanged, gives $J_\omega(\varphi,\psi)=J_\varphi(\psi)$; and
\eqref{eq:ch28-defect-ak} is then
\Cref{thm:ch26-mixed-energy-identities}\,\cref{itm:ch26-ij:J} and
\cref{itm:ch26-ij:comparison}, whose equality clause is
\cref{itm:ch26-ij:positive}.
\end{proof}

\begin{warning}{$J_\omega$ is not symmetric, and where rigidity enters}
Two cautions.  First, $J_\omega(\varphi,\psi)\neq J_\omega(\psi,\varphi)$
in general: by \eqref{eq:ch28-defect-ak} the two are weighted sums of the
same nonnegative numbers $a_k$ with the reversed weights $k+1$ and
$n-k$.  They agree for $n=1$, where both equal $\tfrac12a_0$, and they
always satisfy
$J_\omega(\varphi,\psi)+J_\omega(\psi,\varphi)=I(\varphi,\psi)$ by
\Cref{thm:ch26-mixed-energy-identities}\,\cref{itm:ch26-ij:sum}.  This
asymmetry is the reason
\eqref{eq:measure-dirichlet-quasimetric} produces a \emph{quasi}-metric.

Second, the elementary intersection argument proving
$J_\omega(\varphi,\psi)=0$ does not by itself prove that
$\varphi-\psi$ is constant.  It first shows that all $n$ Dirichlet
coefficients $a_k(\varphi,\psi)$ vanish, hence, by
\Cref{cor:ch26-cauchy-schwarz}, that the $n+1$ mixed measures
$\omega_\varphi^{\,j}\wedge\omega_\psi^{\,n-j}$ coincide; this is
\Cref{prop:ch28-solution-set}\,\cref{itm:ch28-sol:mixed} below.  Passing from
there to $\varphi-\psi\in\mathbf R$ uses the imported uniqueness theorem
\textup{(I2)} of \Cref{sec:ch28-imported}, applied to the two equal endpoint
measures.  Thus vanishing \emph{does} imply rigidity in the full theory;
what is unavailable is a proof from the relatively ample equality case of
\Cref{thm:ch20-mixed-hodge}, since $\omega_\varphi$ and
$\omega_\psi$ are merely relatively nef.  This is the same source
boundary recorded in \Cref{rem:ch25-no-uniqueness} and in the warning after
\Cref{cor:ch26-gauge-nondegenerate}.
\end{warning}

\begin{remark}[Measure energy versus potential energy]
\label{rem:ch28-two-energies}
The adjective \emph{finite energy} occurs on both sides, but the meanings
are not interchangeable by definition.  A measure in
$\mathcal M^1_V$ is a strong limit of model Monge--Amp\`ere measures by
\Cref{thm:finite-energy-measure-completion} below.  Finite dual energy does
already imply that it integrates every $\varphi\in\mathcal E^1$: otherwise
the bounded truncations $\max\{\varphi,-j\}$ would have energies converging
to the finite number $E_\omega(\varphi)$ and integrals decreasing to
$-\infty$, contradicting the bounded-test characterization
\Cref{prop:ch28-test-class}\,\cref{itm:ch28-test:bdd}.  Identifying the measure with $\MA(\varphi)$ for a pointwise
$\varphi\in\mathcal E^1$ requires an additional argument.  In the
standing BFJ framework that argument is
\Cref{thm:ch28-finite-energy-existence}; the abstract measure theory of
\cite{BJ25FiniteEnergy} also applies in settings where this pointwise
identification has not been supplied.
\end{remark}

\section{The variational functional}
\label{sec:ch28-variational}

\subsection{Domain, normalization, and the value of the supremum}
\label{sec:ch28-domain}

For $\mu\in\mathcal M_V(X)$ and a model potential define
\begin{equation}
\label{eq:variational-functional}
                  F_\mu(\varphi)
                  :=E_\omega(\varphi)-\int\varphi\,\rd \mu.
\end{equation}
It is invariant under addition of constants.  We use the same notation on a
larger psh class only after the integral is known to be finite; this is the
case on $\cE^1$ for every finite-dual-energy measure, and on all
$\PSH$ for measures supported on a fixed dual complex.  We generally
normalize by $\sup_{X^{\An}}\varphi=0$.

The three conventions listed in \Cref{sec:ch27-variational-functional}
remain in force and are repeated here because they are used without further
comment.  The \emph{domain} of $F_\mu$ is the set of
$\varphi\in\cE^1(X,\omega)$ with $\int\varphi\,\rd\mu$ finite; the
integral is always defined in $[-\infty,+\infty)$ because $\varphi$ is
upper semicontinuous and bounded above, so finiteness is the only issue.
For $\mu\in\mathcal M^1_V$ finiteness is automatic, by the truncation
argument of \Cref{rem:ch28-two-energies}: it is
\Cref{prop:ch28-test-class}\,\cref{itm:ch28-test:bdd} applied to
$\max\{\varphi,-j\}$, not \eqref{eq:measure-fenchel-energy}, which by
itself covers only model functions.  The domain is therefore all of
$\cE^1(X,\omega)$ when $\mu\in\mathcal M^1_V$, and may be a proper subset
when $\mu$ has infinite dual energy.  The
\emph{sign} is the concave one, and $F_\mu$ is to be maximized.  The
\emph{normalization} is $\sup_{X^{\An}}\varphi=0$, which by
\Cref{thm:bfj-psh-compactness} places the competitors in a compact set and
which loses nothing, since $F_\mu$ is constant on
$\mathbf R$-orbits.

\begin{proposition}
\label{prop:ch28-sup-is-dual-energy}
Let $\mu\in\mathcal M_V(X)$.  Then
\begin{equation}
\label{eq:ch28-sup-is-dual-energy}
 \sup\bigl\{F_\mu(\varphi):\varphi\in\cE^1(X,\omega),\
   \textstyle\int\varphi\,\rd\mu>-\infty\bigr\}
 =E^\vee_\omega(\mu),
\end{equation}
and the same value is obtained if the supremum is restricted to the
normalized competitors $\sup_{X^{\An}}\varphi=0$.  In particular $F_\mu$
is bounded above on its domain if and only if
$\mu\in\mathcal M^1_V(X,\omega)$.
\end{proposition}

\begin{proof}
The first equality is \Cref{prop:ch28-test-class}, case
\cref{itm:ch28-test:e1}, since $F_\mu(\varphi)$ is by definition the
expression whose supremum appears in \eqref{eq:ch28-test-class}.  For the
second assertion, note that if $\varphi$ is in the domain then so is
$\varphi-\sup_{X^{\An}}\varphi$: the supremum is a finite real number,
because $\varphi$ is upper semicontinuous on the compact space
$X^{\An}$, hence attains its maximum by \Cref{cor:usc-attains-max}, and
$\varphi\not\equiv-\infty$ since $E_\omega(\varphi)>-\infty$.  The value
of $F_\mu$ is unchanged by \eqref{eq:dual-energy-constant-test}, the two
masses being equal.  Finally, $F_\mu$ is bounded above on its domain
exactly when the left-hand side of \eqref{eq:ch28-sup-is-dual-energy} is
finite, and $\mathcal M^1_V$ was defined by finiteness of the right-hand
side.
\end{proof}

\begin{warning}{Boundedness and attainment require separate arguments}
\Cref{prop:ch28-sup-is-dual-energy} says that $F_\mu$ is bounded above
precisely on $\mathcal M^1_V$.  Attainment requires the compactness argument of
\Cref{sec:ch28-coercivity} or
\Cref{sec:ch28-general-finite-energy-existence}.  A model maximizing
sequence exists once the supremum is finite, and
\Cref{thm:finite-energy-measure-stability} extracts real information from
it, but attainment still needs a proof.  In the present setting
\Cref{thm:ch28-finite-energy-existence} supplies that proof and bounds the
energies of normalized model maximizing sequences.  The supported
theorem additionally proves continuity, which is a separate step.
\end{warning}

\subsection{Maximizers solve the equation}
\label{sec:ch28-maximizers-solve}

\begin{proposition}
\label{prop:variational-maximizer-solves}
Assume orthogonality.  If $\varphi\in\mathcal E^1(X,\omega)$ maximizes
$F_\mu$ and the pairing with $\mu$ is finite, then
\begin{equation}
\label{eq:variational-ma-equation}
                         \MA_\omega(\varphi)=\mu.
\end{equation}
\end{proposition}

\begin{proof}
Let $f$ be a model function and put
$\varphi_t=P_\omega(\varphi+tf)$.  The envelope satisfies
$\varphi_t\leq\varphi+tf$, so maximality gives
\[
 E_\omega(\varphi_t)-E_\omega(\varphi)
 \leq\int(\varphi_t-\varphi)\,\rd \mu
 \leq t\int f\,\rd \mu
\]
for $t>0$.  The envelope differentiability theorem of
\Cref{chap:bfj-differentiability} yields
\[
                    \int f\,\MA_\omega(\varphi)
                    \leq\int f\,\rd \mu.
\]
Use $-f$ to obtain equality.  Model functions are uniformly dense in the
continuous test functions, and both measures have mass $V$; hence they
are equal.  The use of $P_\omega$ is essential: $\varphi+tf$ need not
be psh.
\end{proof}

The proof just given is the argument in outline.  Because
\Cref{prop:variational-maximizer-solves} is the pivot of the whole chapter,
and because every one of its four steps quotes a statement whose hypotheses
must be checked, we repeat it in full detail, in the exact form in which the
dependency order \eqref{eq:ch27-dependency-diagram} makes it available.  The
reader who is satisfied with the outline may skip to
\Cref{sec:ch28-coercivity}.

\begin{proof}[Second proof, with the hypotheses displayed]
Throughout, $\mu\in\mathcal M_V(X)$ and
$\varphi\in\cE^1(X,\omega)$ with $\int\varphi\,\rd\mu$ finite, and
$\varphi$ maximizes $F_\mu$ over the domain of $F_\mu$.

\emph{Step 1: the competitors are admissible.}  Fix
$f\in C^0(X^{\An})$ and $t\in\mathbf R$.  The obstacle
$\varphi+tf$ is upper semicontinuous and bounded above, and it dominates
the $\omega$-psh function $\varphi-|t|\lVert f\rVert_\infty$, so the
competitor family in \eqref{eq:ch23-envelope} is nonempty and
$\varphi_t:=P_\omega(\varphi+tf)$ is $\omega$-psh by
\Cref{cor:ch23-usc-obstacle}.  By
\Cref{lem:ch27-e1-envelope-calculus}\,(4) one has
$\varphi_t\in\cE^1(X,\omega)$; by
\Cref{lem:ch27-e1-envelope-calculus}\,(2),
\begin{equation}
\label{eq:ch28-envelope-two-sided}
 \varphi-|t|\lVert f\rVert_\infty\ \leq\ \varphi_t\ \leq\ \varphi+tf .
\end{equation}
The left-hand inequality shows that $\int\varphi_t\,\rd\mu$ is finite,
since $\int\varphi\,\rd\mu$ is; so $\varphi_t$ lies in the domain of
$F_\mu$.

\emph{Step 2: the auxiliary function and its maximum.}  Put
\[
 H(t):=E_\omega(\varphi_t)-\int(\varphi+tf)\,\rd\mu
      =E_\omega\bigl(P_\omega(\varphi+tf)\bigr)
       -\int\varphi\,\rd\mu-t\int f\,\rd\mu .
\]
It is real valued, the first term by Step~1 and the last two because
$\int\varphi\,\rd\mu$ is finite and $f$ is bounded with
$\mu(X^{\An})=V<\infty$.  The right-hand inequality of
\eqref{eq:ch28-envelope-two-sided} and $\mu\geq0$ give
$-\int(\varphi+tf)\,\rd\mu\leq-\int\varphi_t\,\rd\mu$, so
\[
 H(t)\ \leq\ F_\mu(\varphi_t)\ \leq\ F_\mu(\varphi)=H(0),
\]
the middle inequality by maximality of $\varphi$ and the last equality
because $P_\omega(\varphi)=\varphi$, $\varphi$ being already
$\omega$-psh.  Thus $H$ attains its maximum over $\mathbf R$ at the
interior point $t=0$.

\emph{Step 3: differentiation.}  This is the only step that uses
orthogonality.  By \Cref{cor:e1-projected-variation}, which assumes
orthogonality and applies because $\varphi\in\cE^1$ and
$f\in C^0(X^{\An})$, the function
$t\mapsto E_\omega(P_\omega(\varphi+tf))$ is differentiable at $t=0$
with derivative $\int f\,\MA_\omega(\varphi)$; note that the measure on the
right is that of $\varphi$ itself and carries no envelope correction, as
\Cref{rem:ch27-e1-derivative-measure} explains.  The remaining part of $H$
is affine in $t$.  Hence $H$ is differentiable at $0$ with
\[
 H'(0)=\int f\,\MA_\omega(\varphi)-\int f\,\rd\mu ,
\]
and a differentiable function with an interior maximum has vanishing
derivative there, so
\begin{equation}
\label{eq:ch28-tested-equality}
 \int f\,\MA_\omega(\varphi)=\int f\,\rd\mu
 \qquad\text{for every }f\in C^0(X^{\An}).
\end{equation}
No separate appeal to $-f$ is needed: \eqref{eq:ch28-tested-equality} is
an equality, obtained from a two-sided derivative, and the outline proof
above obtains it from two one-sided estimates instead.

\emph{Step 4: from tested equality to equality of measures.}  Both
$\MA_\omega(\varphi)$ and $\mu$ are positive Radon measures on the
compact Hausdorff space $X^{\An}$ of total mass $V$: the first by
\Cref{prop:e1-integrability-tail}, the second by hypothesis.  Two positive
Radon measures agreeing on $C^0(X^{\An})$ coincide, by the uniqueness
clause \Cref{cor:riesz-uniqueness} of the Riesz representation theorem.
This gives \eqref{eq:variational-ma-equation}.
\end{proof}

\begin{remark}[Why the envelope cannot be omitted]
\label{rem:ch28-envelope-essential}
The naive variation $\varphi+tf$ is not admissible: for $t\neq0$ it need
not be $\omega$-psh, so $E_\omega(\varphi+tf)$ is not defined by
\Cref{def:extended-energy}, and the inequality $F_\mu(\varphi+tf)\leq
F_\mu(\varphi)$ has no meaning.  Replacing it by $P_\omega(\varphi+tf)$
repairs the domain at the cost of making the $t$-dependence nonlinear, and
the entire difficulty of \Cref{chap:bfj-differentiability} is to show that
the composite is still differentiable.  In the supported existence proof, orthogonality enters when converting
the maximizer into a solution.  In the general finite-energy proof it
also supplies the strong-density and quantitative pairing estimates
used to obtain the maximizer.
\end{remark}

\begin{remark}[The converse, and the two directions of the equivalence]
\label{rem:ch28-converse-direction}
The converse of \Cref{prop:variational-maximizer-solves} is the second half
of \Cref{prop:bfj-euler-lagrange}: if $\MA_\omega(\varphi)=\mu$ then
$\varphi$ maximizes $F_\mu$.  That direction does \emph{not} use
orthogonality; it is the right-hand inequality of the two-sided sandwich
\eqref{eq:ch27-two-sided-sandwich}, exactly as in Step~2 of the proof of
\Cref{prop:ch28-dual-energy-of-ma}.  The asymmetry matters: the class of
measures for which a maximizer is known to exist depends on compactness, and
is enlarged in \Cref{sec:ch28-coercivity}, while the class of measures which
are known to be Monge--Amp\`ere measures is by definition
$\MA_\omega(\cE^1)$, and the two enlargements are logically independent.
\end{remark}

\section{Coercivity and the existence of maximizers}
\label{sec:ch28-coercivity}

\Cref{prop:ch28-sup-is-dual-energy} makes $F_\mu$ bounded above exactly on
$\mathcal M^1_V$.  We now ask when the supremum is attained.  The
mechanism is the classical one---compactness of the competitor set plus
upper semicontinuity of the functional---but each of the two ingredients
must be supplied by a hypothesis on $\mu$, and neither is automatic.

\subsection{Two properties of a measure}
\label{sec:ch28-two-properties}

\begin{definition}[Uniformly integrable and continuously paired measures]
\label{def:ch28-admissible-measure}
\index{admissible measure!BFJ existence}
Let $\mu\in\mathcal M_V(X)$.  We say that $\mu$ is
\begin{enumerate}
\item\label{itm:ch28-adm:coercive}
 \emph{uniformly integrable} for $\omega$ if
 \begin{equation}
 \label{eq:ch28-uniform-integrability}
  M(\mu):=\sup\Bigl\{-\int\varphi\,\rd\mu\ \Bigm|\
   \varphi\in\PSH_0(X,\omega)\Bigr\}<+\infty ;
 \end{equation}
\item\label{itm:ch28-adm:paired}
 \emph{continuously paired} with $\omega$ if the map
 $\varphi\mapsto\int\varphi\,\rd\mu$ is continuous on
 $\PSH_0(X,\omega)$ for the BFJ topology.
\end{enumerate}
A measure with both properties is called \emph{admissible}.  We write
$\mathcal M^{\mathrm{adm}}_V(X,\omega)$ for the set of admissible
measures.
\end{definition}

Here \emph{uniformly integrable} denotes the displayed uniform
$L^1$ bound over normalized psh functions, rather than a separate
uniform-tail-integrability assertion.
The two conditions have different characters.  Condition
\cref{itm:ch28-adm:coercive} is a uniform bound; it makes $F_\mu$
\emph{coercive} in the sense of \Cref{lem:ch28-coercivity}.  Condition
\cref{itm:ch28-adm:paired} is a semicontinuity statement, and it is needed
because the BFJ topology controls restrictions on dual complexes
and does not by definition make pairing with an arbitrary measure continuous.
The portmanteau theorem concerns varying measures against fixed
semicontinuous test functions; it does not by itself settle this
potential-variable pairing.

\begin{lemma}
\label{lem:ch28-coercivity}
\index{coercivity!of the variational functional}
Let $\mu\in\mathcal M_V(X)$ be uniformly integrable, with constant
$M(\mu)$.  Then $\mu\in\mathcal M^1_V(X,\omega)$ with
\begin{equation}
\label{eq:ch28-dual-energy-bound-M}
 E^\vee_\omega(\mu)\ \leq\ M(\mu),
\end{equation}
the domain of $F_\mu$ contains all of $\cE^1(X,\omega)$, and for every
$\varphi\in\cE^1(X,\omega)$ with $\sup_{X^{\An}}\varphi=0$,
\begin{equation}
\label{eq:ch28-coercive-inequality}
 F_\mu(\varphi)\ \leq\ E_\omega(\varphi)+M(\mu)
 \ =\ -\bigl(-E_\omega(\varphi)\bigr)+M(\mu).
\end{equation}
Extend $F_\mu$ to all of $\PSH(X,\omega)$ by the same formula;
the pairing is finite under uniform integrability, so this extension has
values in $[-\infty,+\infty)$.  Consequently, for every $c\leq M(\mu)$,
\begin{equation}
\label{eq:ch28-sublevel-inclusion}
 \bigl\{\varphi\in\PSH_0(X,\omega)\ :\ F_\mu(\varphi)\geq c\bigr\}
 \ \subseteq\ \mathcal K_{M(\mu)-c},
\end{equation}
with $\mathcal K_C$ the normalized energy superlevel
\eqref{eq:normalized-energy-superlevel}; for $c>M(\mu)$ the set on the
left is empty.
\end{lemma}

\begin{proof}
Let $\varphi\in\PSH(X,\omega)$ with $\sup_{X^{\An}}\varphi=0$.  Then
$\varphi\leq0$, so $-\int\varphi\,\rd\mu\in[0,M(\mu)]$; in particular
the integral is finite and $\varphi$ lies in the domain of $F_\mu$ as
soon as $E_\omega(\varphi)>-\infty$.  Since every
$\varphi\in\cE^1(X,\omega)$ differs from a normalized one by a constant,
and both terms of $F_\mu$ shift by the same amount
\eqref{eq:dual-energy-constant-test}, the domain of $F_\mu$ is all of
$\cE^1(X,\omega)$.

Inequality \eqref{eq:ch28-coercive-inequality} is now immediate from
\eqref{eq:variational-functional}:
$F_\mu(\varphi)=E_\omega(\varphi)-\int\varphi\,\rd\mu\leq
E_\omega(\varphi)+M(\mu)$.  Taking the supremum over normalized
$\varphi\in\cE^1$ and using $E_\omega(\varphi)\leq E_\omega(0)=0$
(monotonicity, \Cref{thm:extended-energy-basic}) gives
$E^\vee_\omega(\mu)\leq M(\mu)$ by
\Cref{prop:ch28-sup-is-dual-energy}, which is
\eqref{eq:ch28-dual-energy-bound-M}.

For \eqref{eq:ch28-sublevel-inclusion}, if $\varphi\in\PSH_0(X,\omega)$
satisfies $F_\mu(\varphi)\geq c$, then
\eqref{eq:ch28-coercive-inequality} gives
$E_\omega(\varphi)\geq c-M(\mu)$; in particular
$E_\omega(\varphi)>-\infty$, so $\varphi\in\cE^1$, and $\varphi$ lies
in $\mathcal K_{M(\mu)-c}$ by \eqref{eq:normalized-energy-superlevel},
whose index $M(\mu)-c$ is nonnegative by hypothesis.  If $c>M(\mu)$, the
same inequality together with $E_\omega(\varphi)\leq0$ gives
$c\leq F_\mu(\varphi)\leq M(\mu)$, a contradiction, so the set is empty.
\end{proof}

\begin{remark}[The shape of the coercivity inequality]
\label{rem:ch28-coercivity-shape}
The uniform bound on the linear term in
\eqref{eq:ch28-coercive-inequality} transfers energy coercivity to
$F_\mu$.  For a general $\mu\in\mathcal M^1_V$, Fenchel's inequality alone
does not give this bound: rearranging it makes the two energy terms
cancel.  The stronger quantitative estimate
\cite[Proposition~2.2, equation~(2.7)]{BJ25FiniteEnergy}, however, bounds
the reference Dirichlet energy of a model maximizing sequence.  Together
with the fixed-skeleton bound for $\MA_\omega(0)$, this gives the
normalized energy bound proved in
\Cref{thm:ch28-finite-energy-existence}.  Thus uniform integrability on
all of $\PSH_0$ is a sufficient hypothesis for the shorter argument here,
rather than a necessary hypothesis for finite-energy existence.
\end{remark}

\begin{theorem}
\label{thm:ch28-existence-admissible}
\index{maximizer!existence of}
Let $\mu\in\mathcal M_V(X)$ be admissible in the sense of
\Cref{def:ch28-admissible-measure}.  Then:
\begin{enumerate}
\item\label{itm:ch28-exist:max}
 the functional $F_\mu$ attains its supremum on $\PSH_0(X,\omega)$, at
 some $\varphi\in\cE^1(X,\omega)$ with
 $E_\omega(\varphi)\geq-M(\mu)$;
\item\label{itm:ch28-exist:value}
 the maximum value is $E^\vee_\omega(\mu)\in[0,M(\mu)]$;
\item\label{itm:ch28-exist:equation}
 if in addition $\omega$ has the orthogonality property, then every such
 maximizer satisfies $\MA_\omega(\varphi)=\mu$.
\end{enumerate}
\end{theorem}

\begin{proof}
\emph{Step 1: the competitor set is compact and the functional is upper
semicontinuous.}  By \Cref{thm:bfj-psh-compactness} the normalized space
$\PSH_0(X,\omega)$ is compact for the BFJ topology.  On it, $F_\mu$ is
real valued or $-\infty$, by \Cref{lem:ch28-coercivity}, and it is the sum
of two functions: $E_\omega$, which is upper semicontinuous by
\Cref{thm:extended-energy-basic}, and
$\varphi\mapsto-\int\varphi\,\rd\mu$, which is continuous by
\Cref{def:ch28-admissible-measure}\,\cref{itm:ch28-adm:paired}.  A sum of an
upper semicontinuous function and a continuous one is upper semicontinuous;
no cancellation of infinities occurs, because the second summand takes
values in the bounded interval $[0,M(\mu)]$.

\emph{Step 2: attainment.}  An upper semicontinuous function on a nonempty
compact space attains its maximum, by \Cref{cor:usc-attains-max}.  The space
is nonempty since $0\in\PSH_0(X,\omega)$, and $F_\mu(0)=0$, so the
maximum is at least $0$; by \eqref{eq:ch28-coercive-inequality} it is at
most $M(\mu)$.  Let $\varphi$ be a maximizer.  Then
$F_\mu(\varphi)\geq0$, so \eqref{eq:ch28-sublevel-inclusion} with $c=0$
gives $\varphi\in\mathcal K_{M(\mu)}$; in particular
$\varphi\in\cE^1(X,\omega)$ and $E_\omega(\varphi)\geq-M(\mu)$.  This is
\cref{itm:ch28-exist:max}.

\emph{Step 3: the value.}  The maximum over $\PSH_0$ equals the supremum
over all of $\cE^1(X,\omega)$ by translation invariance, and that supremum
is $E^\vee_\omega(\mu)$ by \Cref{prop:ch28-sup-is-dual-energy}.  Together
with \eqref{eq:ch28-dual-energy-bound-M} and $E^\vee_\omega\geq0$ this is
\cref{itm:ch28-exist:value}.

\emph{Step 4: the equation.}  A maximizer of $F_\mu$ over
$\PSH_0(X,\omega)$ is a maximizer over the whole domain of $F_\mu$, by
Step~3, and its pairing with $\mu$ is finite by
\Cref{lem:ch28-coercivity}.  \Cref{prop:variational-maximizer-solves}
therefore applies and gives \cref{itm:ch28-exist:equation}.
\end{proof}

\begin{remark}[The compactness argument in net form]
\label{rem:ch28-net-form}
Step~2 is stated as an abstract compactness statement, and this is
deliberate.  If one prefers to argue with a maximizing family, the correct
formulation is the following, and it must be run with nets.  Let
$(\varphi_\alpha)$ be a net in $\PSH_0(X,\omega)$ with
$F_\mu(\varphi_\alpha)\to E^\vee_\omega(\mu)$.  Discarding an initial
segment we may assume $F_\mu(\varphi_\alpha)\geq-1$, so by
\eqref{eq:ch28-sublevel-inclusion} the net lies in
$\mathcal K_{M(\mu)+1}$.  By \Cref{cor:ch26-net-superlevel} some subnet
converges in the BFJ topology, uniformly on every dual complex, to a limit
$\varphi\in\cE^1(X,\omega)$ with
$E_\omega(\varphi)\geq
\varlimsup_\beta E_\omega(\varphi_{\alpha_\beta})$ along that subnet; and
$\int\varphi_\alpha\,\rd\mu\to\int\varphi\,\rd\mu$ along the subnet by
\Cref{def:ch28-admissible-measure}\,\cref{itm:ch28-adm:paired}.  Hence
$F_\mu(\varphi)\geq\varlimsup F_\mu(\varphi_\alpha)=E^\vee_\omega(\mu)$,
so $\varphi$ is a maximizer.

This proof uses subnets without assuming sequential compactness.
The space $\PSH_0(X,\omega)$ is
compact but is not known to be metrizable, and
\Cref{lem:ch21-no-cofinal-sequence} shows that the underlying system of
models admits no cofinal sequence; the sequential reductions available in
this theory are \Cref{cor:psh-sequential-regularization} and
\Cref{prop:capacity-subsequence}, and neither applies to an arbitrary
maximizing family.  The policy is that of
\Cref{rem:net-vs-sequence-checklist}.  When $\mu$ is supported on one
dual complex $\Delta$, a sequential argument \emph{is} available on
$\Delta$ itself, because $C^0(\Delta)$ is separable and
\Cref{cor:ch22-arzela-ascoli} gives a compact metrizable family of
restrictions; but the limit so obtained is a function on $\Delta$, not a
psh function on $X^{\An}$, and a restriction on one skeleton alone does not give the compatible
family required by \Cref{thm:bfj-skeletal-characterization}.  We therefore keep the net form
throughout.
\end{remark}

\subsection{Measures supported on a dual complex are admissible}
\label{sec:ch28-skeleton-admissible}

\begin{lemma}
\label{lem:ch28-skeleton-admissible}
\index{dual complex!support of a measure on}
Let $\mathcal X$ be an SNC model on which $\omega$ is determined, with
dual complex $\Delta_{\mathcal X}$, and let $\mu\in\mathcal M_V(X)$ be
supported on $\Delta_{\mathcal X}$.  Then $\mu$ is admissible, and
\begin{equation}
\label{eq:ch28-skeleton-M-bound}
 M(\mu)\ \leq\ C\,\delta(\Delta_{\mathcal X})\,V ,
\end{equation}
where $C=C(\mathcal X,\omega)$ is the equicontinuity constant of
\Cref{thm:bfj-skeletal-equicontinuity} and
$\delta(\Delta_{\mathcal X})$ is the diameter of
\Cref{lem:ch22-path-metric}.
\end{lemma}

\begin{proof}
\emph{Uniform integrability.}  Let $\varphi\in\PSH_0(X,\omega)$.  By
\Cref{thm:bfj-skeletal-equicontinuity}\,(2) --- whose normalization
\eqref{eq:ch22-normalization} is exactly
$\sup_{X^{\An}}\varphi=0$, and which applies because $\omega$ is
determined on $\mathcal X$ --- one has $-C\leq\varphi\leq0$ on
$\Delta_{\mathcal X}$.  The sharper form
\Cref{cor:ch22-explicit-modulus} gives
$\operatorname{osc}_{\Delta_{\mathcal X}}\varphi\leq
C\delta(\Delta_{\mathcal X})$, and
\Cref{prop:ch21-sup-on-skeleton} identifies
$\max_{\Delta_{\mathcal X}}\varphi$ with
$\sup_{X^{\An}}\varphi=0$; hence
\[
 -C\,\delta(\Delta_{\mathcal X})\ \leq\ \varphi\ \leq\ 0
 \qquad\text{on }\Delta_{\mathcal X}.
\]
Since $\mu$ is carried by $\Delta_{\mathcal X}$ and has mass $V$,
\[
 0\ \leq\ -\int\varphi\,\rd\mu
 =-\int_{\Delta_{\mathcal X}}\varphi\,\rd\mu
 \ \leq\ C\,\delta(\Delta_{\mathcal X})\,V .
\]
Taking the supremum over $\varphi$ gives
\eqref{eq:ch28-skeleton-M-bound}, and in particular $M(\mu)<+\infty$.

\emph{Continuous pairing.}  The BFJ topology is, by
\Cref{rem:ch22-topology-comparison}, the topology of uniform convergence on
every dual complex.  Let $\varphi_\alpha\to\varphi$ in
$\PSH_0(X,\omega)$; then
$\varphi_\alpha|_{\Delta_{\mathcal X}}\to\varphi|_{\Delta_{\mathcal X}}$
uniformly, and all these restrictions are continuous functions on the
compact set $\Delta_{\mathcal X}$.  Therefore
\[
 \Bigl|\int\varphi_\alpha\,\rd\mu-\int\varphi\,\rd\mu\Bigr|
 \ \leq\ V\,\bigl\lVert
   \varphi_\alpha-\varphi\bigr\rVert_{C^0(\Delta_{\mathcal X})}
 \ \longrightarrow\ 0 .
\]
By the net criterion \Cref{prop:net-criterion}\,(2) the pairing is
continuous.
\end{proof}

\begin{proposition}
\label{prop:ch28-admissible-calculus}
The set $\mathcal M^{\mathrm{adm}}_V(X,\omega)$ is convex, and it is
closed under the following operations, in the sense stated.
\begin{enumerate}
\item\label{itm:ch28-admcalc:convex}
 If $\mu,\nu$ are admissible and $s\in[0,1]$, then
 $(1-s)\mu+s\nu$ is admissible with
 $M((1-s)\mu+s\nu)\leq(1-s)M(\mu)+sM(\nu)$.
\item\label{itm:ch28-admcalc:dominated}
 If $\mu\in\mathcal M_V(X)$ and $\mu\leq c\,\nu$ for some admissible
 $\nu$ and $c>0$, then $\mu$ is uniformly integrable with
 $M(\mu)\leq cM(\nu)$; it is admissible provided it is also continuously
 paired, for instance if $\mu$ is carried by the support of $\nu$ and
 that support lies in a dual complex.
\item\label{itm:ch28-admcalc:union}
 A measure of mass $V$ supported on a finite union of dual complexes is
 admissible.
\end{enumerate}
\end{proposition}

\begin{proof}
\cref{itm:ch28-admcalc:convex}.  Both defining properties are stable under
convex combinations: the bound \eqref{eq:ch28-uniform-integrability} because
$-\int\varphi\,\rd((1-s)\mu+s\nu)=(1-s)(-\int\varphi\,\rd\mu)
+s(-\int\varphi\,\rd\nu)$, and continuity because a convex combination of
two continuous functions is continuous.  Convexity of the set follows.

\cref{itm:ch28-admcalc:dominated}.  For $\varphi\in\PSH_0$ one has
$-\varphi\geq0$, so $\mu\leq c\nu$ gives
$-\int\varphi\,\rd\mu\leq-c\int\varphi\,\rd\nu\leq cM(\nu)$.  The
statement about continuous pairing is then
\Cref{lem:ch28-skeleton-admissible} applied to the ambient dual complex.

\cref{itm:ch28-admcalc:union}.  The system of SNC models is directed by
domination and cofinal among models, by \Cref{prop:snc-cofinal}, and the
dual complex of a dominating model contains the dual complexes of the models
it dominates, by \eqref{eq:ch21-retraction-composition}.  A finite union of
dual complexes is therefore contained in a single dual complex, and
\Cref{lem:ch28-skeleton-admissible} applies.
\end{proof}

\subsection{Measures in the image are automatically solvable}
\label{sec:ch28-image-solvable}

For a measure already of the form $\MA_\omega(\varphi)$, attainment
follows directly without orthogonality.  This observation also supplies
the nonskeletal example in \Cref{ex:ch28-unbounded-dual-energy}.
Under orthogonality, \Cref{thm:ch28-finite-energy-existence} subsequently
identifies the full image with $\mathcal M^1_V$.

\begin{proposition}
\label{prop:ch28-image-is-solvable}
Let $\varphi\in\cE^1(X,\omega)$ and $\mu=\MA_\omega(\varphi)$.  Then
$\mu\in\mathcal M^1_V(X,\omega)$, the supremum of $F_\mu$ is attained at
$\varphi$, and its value is
$E_\omega(\varphi)-\int\varphi\,\rd\mu$.  No orthogonality hypothesis and
no support hypothesis are used.
\end{proposition}

\begin{proof}
This is \Cref{prop:ch28-dual-energy-of-ma} together with
\Cref{prop:ch28-sup-is-dual-energy}: Step~2 of the former exhibits
$\varphi$ as a maximizer over the test class
\cref{itm:ch28-test:e1}, and the latter identifies the value of the supremum
with $E^\vee_\omega(\mu)$.
\end{proof}

\begin{warning}{Admissibility and solvability are different hypotheses}
Admissibility is an a priori compactness hypothesis on a measure, whereas
solvability means membership in the image of the Monge--Amp\`ere operator.
Under orthogonality, \Cref{thm:ch28-existence-admissible} proves that every
admissible measure is solvable; this inclusion is a theorem, not a
definition.  The converse fails:
\Cref{ex:ch28-unbounded-dual-energy} produces a Monge--Amp\`ere measure of
finite dual energy on $\mathbf P^1$ carried by an infinite set of divisorial points accumulating at a
rational point, which also belongs to its topological support, hence is contained in no
dual complex; the same example proves directly that its
uniform-integrability constant is infinite, so it is solvable but not
admissible in the sense of \Cref{def:ch28-admissible-measure}.  The general finite-energy proof uses the quantitative estimates on
energy levels in place of admissibility on all normalized psh functions.
\end{warning}

\subsection{What can go wrong}
\label{sec:ch28-what-goes-wrong}

We record the two failures that the hypotheses of
\Cref{thm:ch28-existence-admissible} exclude.  Both occur already on
$\mathbf P^1$.

\begin{example}[A measure of infinite dual energy]
\label{ex:ch28-infinite-dual-energy}
\index{dual energy!infinite, example}
Keep the setting of \Cref{ex:ch26-infinite-energy}: $K$ is complete
discretely valued, $X=\mathbf P^1_K$, $L=\cO(1)$, and $\omega$ is the
curvature of the model metric of $(\mathbf P^1_R,\cO(1))$, so that
$\MA_\omega(0)=\delta_\eta$ at the Gauss point and $V=1$.  Let
$x_0\in X^{\An}$ be the rational point $z=0$ and put
\[
 \mu:=\delta_{x_0}\in\mathcal M_1(X).
\]
Then $E^\vee_\omega(\mu)=+\infty$; that is,
$\mu\notin\mathcal M^1_1(\mathbf P^1,\omega)$.

Indeed, let $\varphi_0=\log\lVert z_1\rVert_0=\min\{0,\log|z|\}$ be the
potential of \Cref{ex:ch21-p1-refinement}, which on the ray
$R_0=[\eta,0]$ equals $-t$ in the parameter $t=-\log r$, and let
$\varphi_0^{\langle m\rangle}=\varphi_0\circ p_{\mathcal X_m}$ be its
truncations.  These are $\omega$-psh model functions, and
\Cref{ex:ch26-infinite-energy} computes
\[
 E_\omega\bigl(\varphi_0^{\langle m\rangle}\bigr)=-\frac m2 .
\]
The point $x_0$ lies in the closure of the ray $R_0$ and
$\varphi_0^{\langle m\rangle}(x_0)=-m$, because the retraction
$p_{\mathcal X_m}$ sends $x_0$ to the endpoint $t=m$ of the skeleton.
Hence
\[
 F_\mu\bigl(\varphi_0^{\langle m\rangle}\bigr)
 =E_\omega\bigl(\varphi_0^{\langle m\rangle}\bigr)
  -\varphi_0^{\langle m\rangle}(x_0)
 =-\frac m2+m=\frac m2
 \ \xrightarrow[m\to\infty]{}\ +\infty .
\]
Since each $\varphi_0^{\langle m\rangle}$ is an admissible test object in
\eqref{eq:dual-energy-measure}, this proves
$E^\vee_\omega(\delta_{x_0})=+\infty$.

Three comments.  First, the mass is correct: $\mu(X^{\An})=1=V$.  So the
failure is not the one excluded by \Cref{lem:ch28-mass-forced}; it is a
genuine failure of the Legendre transform to be finite.  Second, the same
computation with $\lambda\varphi_0^{\langle m\rangle}$ for
$\lambda\in(0,1]$ gives $-\lambda^2m/2+\lambda m\to+\infty$ as well, so
no rescaling repairs it.  Third, $\delta_{x_0}$ is supported at a single
point of $X^{\An}$, which shows that \emph{compact support is not the
issue}: what fails is uniform integrability
\eqref{eq:ch28-uniform-integrability}, since
$-\varphi_0^{\langle m\rangle}(x_0)=m$ is unbounded over the normalized
psh space.  The point $x_0$ is rational, hence of type~1, and lies on no
dual complex; \Cref{lem:ch28-skeleton-admissible} does not apply, and by the
present computation it cannot be extended to cover it.
\end{example}

\begin{example}[A positive atom at a rational point]
\label{ex:ch28-whole-space-not-admissible}
The same computation shows that no analogue of
\Cref{lem:ch28-skeleton-admissible} holds with $\Delta_{\mathcal X}$
replaced by all of $X^{\An}$.  Concretely, on the $\mathbf P^1$ of
\Cref{ex:ch28-infinite-dual-energy} put
\[
 \nu:=\tfrac12\delta_\eta+\tfrac12\delta_{x_0}\in\mathcal M_1(X).
\]
Testing against $\varphi_0^{\langle m\rangle}$, whose value at the Gauss
point $\eta$ is $0$, gives
$F_\nu(\varphi_0^{\langle m\rangle})=-m/2+m/2=0$, which is bounded; so this
crude test does not decide finiteness of $E^\vee_\omega(\nu)$.  Testing
instead against $\lambda\varphi_0^{\langle m\rangle}$ gives
$-\lambda^2m/2+\lambda m/2=\lambda(1-\lambda)m/2\to+\infty$ for every
$\lambda\in(0,1)$.  Hence $E^\vee_\omega(\nu)=+\infty$ as well.

More generally, if a positive measure of mass one assigns mass $a>0$ to
$x_0$, then testing with $\lambda\varphi_0^{\langle m\rangle}$ for any
fixed $0<\lambda\leq1$ with $\lambda<2a$ gives a contribution
$(a\lambda-\lambda^2/2)m$ from the energy and the atom, while the remaining
mass contributes nonnegatively because
$\varphi_0^{\langle m\rangle}\leq0$.  Thus a positive atom at this type-1
point destroys finiteness of the dual energy, whatever else the measure
does.  Merely assigning positive mass to each neighbourhood of $x_0$ is
not asserted to do so.
\end{example}

\begin{warning}{Uniform integrability is a hypothesis on $\mu$ and
$\omega$ jointly}
The constant $M(\mu)$ of \eqref{eq:ch28-uniform-integrability} depends on
$\omega$ through the normalized psh space, and the bound
\eqref{eq:ch28-skeleton-M-bound} depends on the model $\mathcal X$ through
both the equicontinuity constant $C$ and the diameter
$\delta(\Delta_{\mathcal X})$.  Refining the model enlarges the dual
complex and, in general, increases the right-hand side; the bound is
therefore not uniform over models, and it is not claimed to be.  What is
uniform, for a fixed model, is the estimate over all normalized psh
functions, and that is exactly what the existence proof needs.
\end{warning}

\section{The strong topology on finite-energy measures}
\label{sec:ch28-strong-measures}

\subsection{Strong convergence and the completion theorem}
\label{sec:ch28-completion}

The set $\mathcal M^1_V(X,\omega)$ carries two topologies, and the
distinction between them repeats, on the measure side, the distinction
between the BFJ topology and the strong topology on $\cE^1$ drawn in
\Cref{sec:ch26-strong}.

\begin{definition}[Strong convergence of measures]
\label{def:ch28-strong-measure-convergence}
\index{strong topology!on finite-energy measures}
A net $\mu_i$ in $\mathcal M^1_V(X,\omega)$ converges
\emph{strongly} to $\mu\in\mathcal M^1_V(X,\omega)$ if it converges to
$\mu$ weakly and
\begin{equation}
\label{eq:ch28-strong-measure-convergence}
                  E_\omega^\vee(\mu_i)\longrightarrow
                  E_\omega^\vee(\mu).
\end{equation}
The \emph{strong topology} on $\mathcal M^1_V(X,\omega)$ is the coarsest
topology refining the weak topology for which $E^\vee_\omega$ is
continuous.
\end{definition}

Since $E^\vee_\omega$ is already weakly lower semicontinuous
(\Cref{prop:dual-energy-elementary}), the strong topology is generated by
the weakly open sets together with the sets
$\{E^\vee_\omega<c\}$, $c\in\mathbf R$; and a net converges strongly if
and only if it converges in that topology, by the net criterion
\Cref{prop:net-criterion}.  Strong convergence implies weak convergence by
definition.  The converse fails even when the weak limit has finite
dual energy: the distinct-branch sequence in
\Cref{rem:ch28-weak-not-strong} converges weakly to $\delta_\eta$ while
its dual energies stay at a fixed positive value.
Finite dual energy also fails to be preserved by arbitrary weak limits:
on the $\mathbf P^1$ of \Cref{ex:ch28-infinite-dual-energy}, the measures
$\MA_\omega(\varphi_0^{\langle m\rangle})=\delta_{t=m}$ converge weakly to
$\delta_{x_0}$, which has infinite dual energy, whereas
$E^\vee_\omega(\delta_{t=m})$ is finite for each $m$ --- by
\Cref{prop:ch28-dual-energy-of-ma} it equals
$E_\omega(\varphi_0^{\langle m\rangle})
-\varphi_0^{\langle m\rangle}(t=m)=-m/2+m=m/2$.  So the weak limit is not
even in $\mathcal M^1_V$, and the dual energies diverge.  Weak limits do
not preserve finite dual energy; this is the exact measure-side analogue of
\Cref{thm:ch26-e1-structure}\,\cref{itm:ch26-e1:not-closed}.

\begin{theorem}
\label{thm:finite-energy-measure-completion}
Assume orthogonality.  The Monge--Amp\`ere measures of psh model functions
are strongly dense in the space $\mathcal M^1_V(X,\omega)$.  There is a
complete quasi-metric $\delta_\omega$, defining its strong topology, such
that
\begin{equation}
\label{eq:measure-dirichlet-quasimetric}
 \delta_\omega\bigl(\MA_\omega(\varphi),\MA_\omega(\psi)\bigr)
       =V^{-1}J_\omega(\varphi,\psi)
\end{equation}
for psh model functions $\varphi,\psi$.
\end{theorem}

\begin{proof}[Proof architecture]
Choose a model sequence computing the Legendre supremum for $\mu$.
Uniform differentiability of the equilibrium energy, which is equivalent to
orthogonality, shows that its Monge--Amp\`ere measures converge weakly to
$\mu$.  Mixed-energy Hölder estimates upgrade this convergence to the
strong topology and show that \eqref{eq:measure-dirichlet-quasimetric}
extends uniquely.  The same estimates prove the quasi-triangle inequality
and completeness.  This combines Theorems~2.22, 2.23 and~2.25 of
\cite{BJ25FiniteEnergy}, equivalently its Theorem~A, and is transported
from probability measures by the bijection
$\mu\mapsto V^{-1}\mu$.  Under this bijection both the potential energy and
the measure energy are divided by $V$, which accounts for the factor on the
right of \eqref{eq:measure-dirichlet-quasimetric}.
\end{proof}

The theorem is imported; it is item \textup{(I3)} of
\Cref{sec:ch28-imported}.  Four reading notes make precise what it does and
does not say, and what part of it is proved in this book.

\begin{remark}[What is proved here]
\label{rem:ch28-what-of-completion-is-proved}
The density half of \Cref{thm:finite-energy-measure-completion} is proved
here, in the following form: for every $\mu\in\mathcal M^1_V(X,\omega)$
and every sequence $(\varphi_j)$ of $\omega$-psh model functions computing
the supremum in \eqref{eq:dual-energy-measure}, the measures
$\MA_\omega(\varphi_j)$ converge to $\mu$ strongly in the sense of
\Cref{def:ch28-strong-measure-convergence}.  That statement is
\Cref{thm:finite-energy-measure-stability}, and its proof, given in
\Cref{sec:ch28-stability}, is complete except for the two quantitative
inputs \eqref{eq:maximizing-uniform-differentiability} and
\eqref{eq:maximizing-mixed-energy-control}, which are quoted.  The
completeness half, and the fact that $\delta_\omega$ defines the strong
topology, are not proved here.
\end{remark}

\begin{remark}[Why the right-hand side of
\eqref{eq:measure-dirichlet-quasimetric} depends only on the two measures]
\label{rem:ch28-quasimetric-well-defined}
The identity \eqref{eq:ch28-defect-identity} shows that
$J_\omega(\varphi,\psi)$ depends on its first argument only through
$\mu=\MA_\omega(\varphi)$:
\[
 J_\omega(\varphi,\psi)
   =E^\vee_\omega(\mu)-E_\omega(\psi)+\int\psi\,\rd\mu.
\]
If $\MA_\omega(\psi)=\MA_\omega(\psi')=\nu$, Fenchel equality gives
$J_\omega(\psi,\psi')=J_\omega(\psi',\psi)=0$.  Choose a finite
reference-energy bound containing the three model potentials.  The
H\"older continuity estimate
\cite[Theorem~1.33, equation~(1.38)]{BJ25FiniteEnergy} then gives
$J_\omega(\varphi,\psi)=J_\omega(\varphi,\psi')$.  Hence the formula
descends to the two model measures using the energy estimates, without
requiring injectivity of the potential-to-measure map.  Strong density
and the continuity estimates extend it to all finite-energy measures.
Uniqueness of pointwise potentials is a separate statement,
\textup{(I2)}.
\end{remark}

\begin{remark}[This is not a metric on potentials]
\label{rem:ch28-not-a-metric-on-potentials}
The quoted completion theorem equips measures with $\delta_\omega$.
The maximum-based potential gauge \eqref{eq:ch26-gauge} fails the
triangle inequality, as \Cref{ex:ch26-triangle-fails} shows; no triangle
inequality is established here for the rooftop gauge
\eqref{eq:ch26-rooftop-gauge} either.  Once
\Cref{thm:ch28-finite-energy-existence} identifies
$\cE^1/\mathbf R$ with $\mathcal M^1_V$ under orthogonality, one can
transport the complete quasi-metric along that bijection.  This does
not identify the transported object with either of those gauges, or
by itself identify its topology with the potential strong topology.
The metric theory for a trivially valued field is developed in
\Cref{chap:bj-strong-topologies}.
\end{remark}

\begin{remark}[Completion of measures is not surjectivity of $\MA_\omega$]
\label{rem:ch28-completion-not-surjectivity}
The measure-side completion theorem alone supplies no pointwise
potential for an arbitrary member of the completion.  In the BFJ
framework of this chapter, the additional compactness and pairing
arguments in \Cref{thm:ch28-finite-energy-existence} supply that
potential.  In particular, the required normalized energy bound follows
from quantitative concavity rather than from a skeletal support
assumption on the prescribed measure.  Extending the pointwise singular
Monge--Amp\`ere operator to more general fields remains a separate
question; it must not be confused with surjectivity in the present
discretely valued framework.
\end{remark}


\section{Finite-energy existence beyond skeletal support}
\label{sec:ch28-general-finite-energy-existence}

The quantitative estimates in the completion theorem bridge the gap
between measure convergence and pointwise potential existence in the
BFJ setting.  We give the bridge explicitly.  This argument uses the
standing assumptions of \Cref{eq:ch28-standing}, including the ample
class and the established pointwise $\cE^1$ calculus, and additionally
assumes orthogonality.

\begin{theorem}[Finite-energy existence]
\label{thm:ch28-finite-energy-existence}
Assume orthogonality.  Every
$\mu\in\mathcal M^1_V(X,\omega)$ has a unique normalized solution
$\varphi_\mu\in\cE^1(X,\omega)$ of
$\MA_\omega(\varphi_\mu)=\mu$.  It maximizes $F_\mu$, and every
normalized model maximizing sequence converges to it in the strong
potential topology.  Continuity of $\varphi_\mu$ is not asserted.
\end{theorem}

\begin{proof}
Write $\mu_0:=\MA_\omega(0)$.  This model measure is carried by a fixed
dual complex.  By \Cref{lem:ch28-skeleton-admissible} there is $M_0<\infty$
such that
\[
             -\int\varphi\,\rd\mu_0\leq M_0
             \qquad(\varphi\in\PSH_0(X,\omega)).
\]
Choose normalized model potentials $\psi_j$ with
$F_\mu(\psi_j)\to E_\omega^\vee(\mu)$, and set
$\varepsilon_j:=E_\omega^\vee(\mu)-F_\mu(\psi_j)\to0$.
These defects are nonnegative; no monotonicity of the sequence is needed.
The quantitative concavity estimate
\cite[Proposition~2.2, equation~(2.7)]{BJ25FiniteEnergy}, applied with
the two model potentials $\psi_j,0$, gives
\[
 J_\omega(\psi_j,0)
 \leq A'\bigl(\varepsilon_j+E_\omega^\vee(\mu)\bigr).
\]
Quasi-symmetry absorbs $A'$ into a constant $A$ depending only on $n$.
The source uses probability measures and energies divided by $V$;
multiplication by $V$ gives this same inequality in our normalization.
On the tail where $\varepsilon_j\leq1$, therefore,
\begin{equation}
\label{eq:ch28-general-energy-bound}
 -E_\omega(\psi_j)
 =J_\omega(0,\psi_j)-\int\psi_j\,\rd\mu_0
 \leq A\bigl(E_\omega^\vee(\mu)+1\bigr)+M_0=:C.
\end{equation}
Thus the sequence lies eventually in the compact set $\mathcal K_C$
of \Cref{prop:e1-compact-superlevels}.

We next prove continuity of the pairing
$\varphi\mapsto\int\varphi\,\rd\mu$ on each $\mathcal K_C$.
Finite dual energy implies that $\mu$ integrates every $\cE^1$ potential,
by \Cref{rem:ch28-two-energies}.  By strong density choose model
Monge--Amp\`ere measures $\nu_k\to\mu$ strongly.  Their dual energies
are bounded on a tail, and $\delta_\omega(\mu,\nu_k)\to0$.
For $\varphi\in\mathcal K_C$,
\[
 0\leq J_\omega(0,\varphi)
       =-E_\omega(\varphi)+\int\varphi\,\rd\mu_0\leq C.
\]
To apply the quoted model estimates, take decreasing model
regularizations $\chi_\alpha\searrow\varphi$.  Monotone energy
continuity and finite pairing with $\mu_0$ show
$J_\omega(0,\chi_\alpha)\to J_\omega(0,\varphi)$.
Consequently a tail has $J_\omega(0,\chi_\alpha)\leq C+1$.
The estimate
\cite[Theorem~2.23(iv), equation~(2.25)]{BJ25FiniteEnergy},
with test potentials $\chi_\alpha,0$, gives a constant $B$ independent
of $\varphi$, $k$, and that tail such that
\[
 \left|\int\chi_\alpha\,\rd(\mu-\nu_k)\right|
       \leq B\,\delta_\omega(\mu,\nu_k)^{1/2}.
\]
Here the reference energies of the measures and tests are bounded;
quasi-symmetry bounds the directed test energy, and the fixed constants
absorb the probability normalization.  Both measures integrate
$\varphi$.  The decreasing-net integral theorem applied separately to
their finite integrals, using \Cref{cor:usc-net-monotone-integral},
therefore passes this estimate to $\varphi$:
\begin{equation}
\label{eq:ch28-energy-level-pairing}
 \sup_{\varphi\in\mathcal K_C}
 \left|\int\varphi\,\rd(\mu-\nu_k)\right|
       \leq B\,\delta_\omega(\mu,\nu_k)^{1/2}\longrightarrow0.
\end{equation}
Each $\nu_k$ is carried by a fixed dual complex, so its pairing is
continuous for the BFJ topology by uniform convergence of psh
restrictions there.  The uniform limit
\eqref{eq:ch28-energy-level-pairing} proves continuity of the pairing
with $\mu$ on $\mathcal K_C$.

Energy is upper semicontinuous on $\mathcal K_C$, hence $F_\mu$ is
upper semicontinuous there and attains its maximum.  By
\eqref{eq:ch28-general-energy-bound} this maximum equals
$E_\omega^\vee(\mu)$.  The Euler--Lagrange implication
\Cref{prop:variational-maximizer-solves} gives a solution; uniqueness
modulo constants is \cite[Proposition~8.2]{BFJ15}.
Finally every cluster point of the normalized maximizing sequence in
$\mathcal K_C$ is a maximizer and hence this same solution.  Compactness
and uniqueness of cluster points imply convergence of the entire
sequence in the BFJ topology, without sequential compactness.
Continuity of the pairing and convergence of $F_\mu(\psi_j)$ then imply
$E_\omega(\psi_j)\to E_\omega(\varphi_\mu)$, which is precisely strong
potential convergence.
\end{proof}

\subsection{The Fenchel equality at a solution}
\label{sec:ch28-fenchel-equality}

\begin{corollary}
\label{cor:finite-energy-variational-characterization}
If $\varphi_\mu\in\mathcal E^1$ satisfies
$\MA_\omega(\varphi_\mu)=\mu$, then it maximizes $F_\mu$, modulo
constants, among the potentials for which the pairing with $\mu$ is
finite.  Moreover,
\begin{equation}
\label{eq:dual-energy-attained}
 E_\omega^\vee(\mu)
 =E_\omega(\varphi_\mu)-
      \int\varphi_\mu\,\rd \mu.
\end{equation}
\end{corollary}

\begin{proof}
Concavity of energy and its first variation imply, for every
$\psi\in\mathcal E^1$ for which $\int\psi\,\rd \mu$ is finite,
\[
 E_\omega(\psi)-E_\omega(\varphi_\mu)
 \leq\int(\psi-\varphi_\mu)\MA_\omega(\varphi_\mu).
\]
Substitute $\mu=\MA_\omega(\varphi_\mu)$.  Thus
$F_\mu(\psi)\leq F_\mu(\varphi_\mu)$.  Under orthogonality, any other
maximizer solves the same equation by
\Cref{prop:variational-maximizer-solves}, and
\cite[Proposition~8.2]{BFJ15} then gives equality modulo constants.

It remains to compare this extended maximum with the model-test supremum
in \Cref{def:dual-energy-measure}.  Choose a decreasing model
regularization $\varphi_j\searrow\varphi_\mu$.  The mixed
finite-energy integrability theorem and the monotone energy cocycle give
\[
 E_\omega(\varphi_j)\longrightarrow E_\omega(\varphi_\mu),
 \qquad
 \int\varphi_j\,\rd \MA_\omega(\varphi_\mu)
   \longrightarrow
 \int\varphi_\mu\,\rd \MA_\omega(\varphi_\mu).
\]
Consequently $F_\mu(\varphi_j)\to F_\mu(\varphi_\mu)$.  Since each
$\varphi_j$ is an admissible test in
\eqref{eq:dual-energy-measure}, the displayed value is exactly
$E_\omega^\vee(\mu)$.  Thus ``attained'' here refers to the natural
finite-energy extension of the model-test functional.
\end{proof}

The corollary repeats, in the language of the dual energy, what
\Cref{prop:ch28-dual-energy-of-ma} established; the two statements are the
same fact seen from opposite sides, and it is worth saying which parts of
the proof just given are unconditional.  The inequality
$F_\mu(\psi)\leq F_\mu(\varphi_\mu)$ is the right-hand branch of the
two-sided sandwich \eqref{eq:ch27-two-sided-sandwich} and uses nothing but
step \textup{(D6)}.  The identity \eqref{eq:dual-energy-attained} is
\eqref{eq:ch28-dual-energy-of-ma}, and the passage from the enlarged test
class to the model test class is \Cref{prop:ch28-test-class}, which uses
decreasing regularization \textup{(D3)}.  The assertion about equality modulo constants additionally uses
orthogonality to identify a maximizer with a solution and the imported
uniqueness statement \textup{(I2)}, discussed in
\Cref{sec:ch28-uniqueness}.  Neither is needed for the displayed value
of the maximum at a given solution.

\begin{remark}[The Fenchel picture]
\label{rem:ch28-fenchel-picture}
The pair $(E_\omega,E^\vee_\omega)$ satisfies Fenchel's inequality,
and its defect at $(\psi,\MA_\omega(\varphi))$ is
$J_\omega(\varphi,\psi)$.  Under orthogonality, equality for a prescribed
finite-energy measure identifies a maximizer with a solution through
\Cref{prop:variational-maximizer-solves}.
There is also an unconditional biconjugation identity on $\cE^1$:
\[
 E_\omega(\varphi)=
 \min_{\mu\in\mathcal M^1_V}
       \left(E^\vee_\omega(\mu)+\int\varphi\,\rd\mu\right).
\]
All the pairings are finite.  Fenchel's inequality gives one direction,
and $\mu=\MA_\omega(\varphi)$ attains equality by
\Cref{prop:ch28-dual-energy-of-ma}.  This identity does not itself
construct a potential for a prescribed measure; that construction uses
the compactness and continuity argument in
\Cref{thm:ch28-finite-energy-existence}.
\end{remark}

\section{The original BFJ existence theorem}
\label{sec:ch28-bfj-theorem}

We now return to the historical theorem, where no dual-energy assumption is
imposed explicitly.  The logical shape is: the support hypothesis implies
admissibility (\Cref{lem:ch28-skeleton-admissible}), admissibility implies
existence of a maximizer (\Cref{thm:ch28-existence-admissible}),
orthogonality converts the maximizer into a solution
(\Cref{prop:variational-maximizer-solves}), and the support hypothesis
enters a second time, through a Dini argument, to make the solution
continuous.

\begin{lemma}
\label{lem:skeletal-measure-finite-energy}
Let $\Delta_{\mathcal X}$ be the dual complex of an SNC model and let
$\mu\in\mathcal M_V(X)$ be supported on $\Delta_{\mathcal X}$.  Then
$E_\omega^\vee(\mu)<+\infty$.
\end{lemma}

\begin{proof}
Refine $\mathcal X$ to an SNC model that also determines $\omega$.
Its dual complex contains $\Delta_{\mathcal X}$, so it suffices to work
on that refined complex.  Normalize $\sup\varphi=0$.
The restrictions to this fixed compact complex of all normalized $\omega$-psh functions form an
equicontinuous, uniformly bounded family.  This is the skeletal estimate of
\Cref{thm:bfj-skeletal-equicontinuity}.  Hence
$-\int\varphi\,\rd \mu$ is bounded above uniformly in $\varphi$.  Since
$E_\omega(\varphi)\leq E_\omega(0)=0$, the supremum
\eqref{eq:dual-energy-measure} is finite.
\end{proof}

Quantitatively, the proof gives more than finiteness:
\Cref{lem:ch28-skeleton-admissible} bounds
$E^\vee_\omega(\mu)$ by $M(\mu)\leq C\delta(\Delta_{\mathcal X'})V$,
where $\mathcal X'$ is such a common refining model, with
$C$ the equicontinuity constant of
\Cref{thm:bfj-skeletal-equicontinuity} and
$\delta(\Delta_{\mathcal X'})$ the diameter of the refined dual complex
in its path metric.  The constant is explicit once a model and a relatively ample class
have been fixed; \Cref{ex:ch28-chain-dual-energy} evaluates both sides on a
four-component chain and measures the loss.

\subsection{Statement, with the hypothesis signature displayed}
\label{sec:ch28-bfj-statement}

\begin{center}
\fbox{\begin{minipage}{0.91\textwidth}
\textbf{Hypothesis signature of \Cref{thm:bfj-supported-cy}.}
\begin{itemize}
\item \emph{Base field:} $K$ complete, \emph{discretely} valued, with
$v_K(\pi)=1$; residue field of characteristic $0$.
\item \emph{Space:} $X$ smooth, connected, projective over $K$, of
dimension $n$.
\item \emph{Class:} $\omega$ a closed model $(1,1)$-form, semipositive,
with ample de Rham class $\{\omega\}$; $V=(\{\omega\}^n)>0$.
\item \emph{Positivity input:} $\omega$ has the orthogonality property
\Cref{def:orthogonality-property}.  This is a hypothesis, not a
consequence; its provenances are \Cref{thm:orthogonality-provenance}.
\item \emph{Measure:} $\mu\geq0$ with $\mu(X^{\An})=V$, supported on the
dual complex $\Delta_{\mathcal X}$ of some SNC model $\mathcal X$.
\item \emph{Not assumed:} algebraizability of $X$ over a one-variable
function field; continuity, density, or atomicity of $\mu$;
factorization of the solution through any retraction.
\end{itemize}
\end{minipage}}
\end{center}

\begin{theorem}
\label{thm:bfj-supported-cy}
Assume $\omega$ has the orthogonality property.  If
$\mu\in\mathcal M_V(X)$ is supported on the dual complex of an SNC model,
there is a unique continuous $\omega$-psh function $\varphi$, normalized
by $\sup\varphi=0$, such that
\begin{equation}
\label{eq:bfj-supported-cy}
                         \MA_\omega(\varphi)=\mu.
\end{equation}
\end{theorem}

\begin{proof}[The three BFJ steps]
First, the normalized psh space is weakly compact.  On the fixed dual
complex, psh restrictions converge uniformly, so
$\varphi\mapsto\int\varphi\,\rd \mu$ is continuous.  Energy is upper
semicontinuous.  Thus $F_\mu$ attains its maximum; since
$F_\mu(0)=0$, the maximizer has finite energy.

Second, \Cref{prop:variational-maximizer-solves}, which is where
orthogonality enters, proves the equation in $\mathcal E^1$.

Third, let $\varphi_j\searrow\varphi$ be a model regularization.  The
restriction $\varphi|_{\Delta_{\mathcal X}}$ is continuous, so Dini's
lemma gives uniform convergence there.  The capacity domination estimate
of Lemma~8.4 of \cite{BFJ15} propagates the resulting almost-everywhere
inequality from $\mu=\MA(\varphi)$ to all of $X^{\An}$.  Hence
$\varphi_j\to\varphi$ uniformly globally, proving continuity.  Uniqueness
is the non-Archimedean Calabi--Yau uniqueness theorem, quoted in the
expanded proof below; it does not use orthogonality, and it is not a
consequence of the comparison principle.
\end{proof}

\subsection{The proof, step by step}
\label{sec:ch28-bfj-proof}

We now give the same argument with every quotation identified.  The reader
should compare the four steps below with the four rows of the second column
of the table in \Cref{sec:ch28-scope}.

\begin{proof}[Detailed proof of \Cref{thm:bfj-supported-cy}]
Let $\mathcal X$ be an SNC model with $\mu$ supported on
$\Delta:=\Delta_{\mathcal X}$.  Enlarging $\mathcal X$ if necessary --- SNC
models are cofinal by \Cref{prop:snc-cofinal}, and the dual complex of a
dominating model contains $\Delta$ --- we may assume that $\omega$ is
determined on $\mathcal X$.

\emph{Step 1: $\mu$ is admissible, and $F_\mu$ attains its maximum.}
By \Cref{lem:ch28-skeleton-admissible}, $\mu$ is uniformly integrable with
$M(\mu)\leq C\delta(\Delta)V<\infty$ and continuously paired.  By
\Cref{thm:ch28-existence-admissible}\,\cref{itm:ch28-exist:max} the
functional $F_\mu$ attains its maximum on $\PSH_0(X,\omega)$ at some
$\varphi$ with
\[
 0=F_\mu(0)\ \leq\ F_\mu(\varphi)=E^\vee_\omega(\mu)\ \leq\ M(\mu),
 \qquad
 E_\omega(\varphi)\ \geq\ -M(\mu)>-\infty ,
\]
so $\varphi\in\cE^1(X,\omega)$.  The compactness used is that of
\Cref{thm:bfj-psh-compactness}, in the net form of
\Cref{cor:ch26-net-superlevel}; see \Cref{rem:ch28-net-form} for why no
sequence is extracted.  This step uses the support hypothesis and no
orthogonality.

\emph{Step 2: the maximizer solves the equation.}  By
\Cref{thm:ch28-existence-admissible}\,\cref{itm:ch28-exist:equation}, which
is \Cref{prop:variational-maximizer-solves}, one has
$\MA_\omega(\varphi)=\mu$.  This is the only step that uses
orthogonality, and it uses it only through
\Cref{cor:e1-projected-variation}, hence only at \emph{continuous}
obstacles through model regularization and passage to the
finite-energy limit; it does not apply the continuous-obstacle theorem
directly to the singular obstacle $\varphi+tf$.
See \Cref{rem:ch27-e1-hypotheses}.

\emph{Step 3: the solution is continuous.}  This is
\Cref{thm:bfj-supported-regularity}, proved in
\Cref{chap:bfj-continuity}, and we recall its mechanism because it is the
second and last use of the support hypothesis.  Subtracting a constant,
which changes neither the equation nor the continuity to be proved, assume
$\sup\varphi=-1$, and choose, by
\Cref{thm:bfj-decreasing-regularization}, a decreasing net
$(\varphi_j)$ of $\omega$-psh model functions with
$\varphi_j\searrow\varphi$.  The closed sets $\{\varphi_j\geq0\}$
form a filtered decreasing family with empty intersection because
$\varphi\leq-1$.  Compactness therefore makes one such set empty,
so on a cofinal tail $\varphi_j<0$.  The restriction $\varphi|_\Delta$ is continuous by
\Cref{thm:bfj-skeletal-equicontinuity}\,(1), and the continuous functions
$\varphi_j|_\Delta$ decrease to it on the compact polyhedral complex
$\Delta$; Dini's theorem for nets therefore gives, for each
$\varepsilon>0$, an index beyond which
\[
 \varphi_j-\varepsilon\ \leq\ \varphi
 \qquad\text{on }\Delta ,
\]
hence $\mu$-almost everywhere, since $\mu$ is carried by $\Delta$.
Because $\mu=\MA_\omega(\varphi)$ by Step~2, the domination lemma
\Cref{lem:bfj-supported-domination}, applied with the continuous psh
function $h=\varphi_j-\varepsilon$ and $\psi=\varphi$, propagates the
inequality to all of $X^{\An}$:
$0\leq\varphi_j-\varphi\leq\varepsilon$ everywhere.  Thus
$\varphi_j\to\varphi$ uniformly on $X^{\An}$ and $\varphi$ is
continuous.  The propagation rests on the capacity estimate
\Cref{lem:bfj-kolodziej-sublevel} and on the fact that a nonempty open set
has positive capacity, \Cref{lem:divisorial-positive-capacity}.

\emph{Step 4: uniqueness.}  Let $\varphi,\psi$ be two continuous
$\omega$-psh solutions with $\sup\varphi=\sup\psi=0$.
\Cref{prop:ch28-solution-set} shows unconditionally that
$I(\varphi,\psi)=0$, that all the Dirichlet coefficients
$a_k(\varphi,\psi)$ vanish, and that the $n+1$ mixed measures
$\omega_\varphi^{\,j}\wedge\omega_\psi^{\,n-j}$ all equal $\mu$.  For
$n=1$ this already gives $\varphi=\psi$,
\Cref{prop:ch28-curve-uniqueness}.  For $n\geq2$ the remaining step is the
non-Archimedean Calabi--Yau uniqueness theorem, which we quote:
$\varphi-\psi$ is constant by \cite[Proposition~8.2]{BFJ15}, resting on
the arithmetic Hodge index theorem of Yuan and Zhang.  The normalization
$\sup\varphi=\sup\psi=0$ then forces $\varphi=\psi$.  This step is
\emph{not} a consequence of the comparison principle
(\Cref{rem:ch25-no-uniqueness}) and does not use orthogonality.
\end{proof}

\begin{warning}{The four hypotheses do four different jobs}
Step~1 uses the support hypothesis to obtain compactness \emph{and}
continuity of the pairing; Step~2 uses orthogonality and nothing else;
Step~3 uses the support hypothesis a second time, now to run Dini on a
compact set carrying $\mu$, together with capacity; Step~4 uses neither
the support hypothesis nor orthogonality, but does use an imported
intersection-theoretic theorem.  It is a common error to compress Steps~1
and~3 into ``$\mu$ lives on a skeleton, so everything is finite
dimensional''.  It is not: the solution need not factor through the
retraction to $\Delta$, and Boucksom--Jonsson construct a finite atomic
measure supported on a dual complex whose solution is invariant under no
SNC retraction \cite[Theorem~B]{BJGreen24}.  What the support hypothesis
gives is a compact set on which all normalized psh functions are uniformly
bounded and uniformly convergent, not a reduction of the problem to that
set.
\end{warning}

\subsection{The structure of the solution set, and uniqueness}
\label{sec:ch28-uniqueness}

The following proposition isolates exactly how much of uniqueness is
elementary.  It uses no orthogonality, no support hypothesis, and no
imported theorem.

\begin{proposition}
\label{prop:ch28-solution-set}
\index{uniqueness!Monge--Ampere solutions}
Let $\mu\in\mathcal M_V(X)$ and let
$\varphi,\psi\in\cE^1(X,\omega)$ satisfy
$\MA_\omega(\varphi)=\MA_\omega(\psi)=\mu$.  Then:
\begin{enumerate}
\item\label{itm:ch28-sol:I}
 $I(\varphi,\psi)=0$ and $J_\omega(\varphi,\psi)=
 J_\omega(\psi,\varphi)=0$, where $I$ is
 \eqref{eq:ch26-I-definition} and $J_\omega$ is
 \eqref{eq:directed-dirichlet-functional};
\item\label{itm:ch28-sol:energy}
 $E_\omega(\varphi)-E_\omega(\psi)=\int(\varphi-\psi)\,\rd\mu$, and both
 potentials give the same value $E^\vee_\omega(\mu)$ in
 \eqref{eq:dual-energy-attained};
\item\label{itm:ch28-sol:mixed}
 if $\varphi$ and $\psi$ are bounded, then all the Dirichlet
 coefficients $a_k(\varphi,\psi)$ of \eqref{eq:ch26-ak} vanish and the
 $n+1$ mixed measures
 $\omega_\varphi^{\,j}\wedge\omega_\psi^{\,n-j}$, $0\leq j\leq n$,
 coincide with $\mu$;
\item\label{itm:ch28-sol:convex}
 the set of maximizers of $F_\mu$ in $\cE^1(X,\omega)$ is convex, and
 $E_\omega$ is affine on every segment joining two of its points;
\item\label{itm:ch28-sol:max}
 $\chi:=\max\{\varphi,\psi\}$ lies in $\cE^1(X,\omega)$, and
 $\MA_\omega(\chi)$ agrees with $\mu$ on the set
 $\{\varphi\neq\psi\}$; moreover
 $\MA_\omega(\chi)(\{\varphi=\psi\})=\mu(\{\varphi=\psi\})$.
\end{enumerate}
\end{proposition}

\begin{proof}
\cref{itm:ch28-sol:I}.  By \eqref{eq:ch26-I-definition},
\[
 I(\varphi,\psi)=\int(\varphi-\psi)
   \bigl(\MA_\omega(\psi)-\MA_\omega(\varphi)\bigr)
 =\int(\varphi-\psi)(\mu-\mu)=0 ,
\]
the integrals being finite by \Cref{lem:ch27-mutual-integrability}.  By
\Cref{prop:ch28-fenchel-defect} and
\Cref{thm:ch26-mixed-energy-identities}\,\cref{itm:ch26-ij:sum} one has
$J_\omega(\varphi,\psi)+J_\omega(\psi,\varphi)=I(\varphi,\psi)=0$, and
both summands are nonnegative, so both vanish.

\cref{itm:ch28-sol:energy}.  Vanishing of $J_\omega(\psi,\varphi)$ is,
by \eqref{eq:directed-dirichlet-functional},
$E_\omega(\psi)-E_\omega(\varphi)+\int(\varphi-\psi)\MA_\omega(\psi)=0$,
which is the stated identity since $\MA_\omega(\psi)=\mu$.  The second
assertion follows by substituting into \eqref{eq:dual-energy-attained}.

\cref{itm:ch28-sol:mixed}.  For bounded potentials
\Cref{thm:ch26-mixed-energy-identities}\,\cref{itm:ch26-ij:I} gives
$\sum_{k=0}^{n-1}a_k(\varphi,\psi)=I(\varphi,\psi)=0$, and each $a_k$ is
nonnegative by \Cref{lem:ch26-dirichlet-coefficients}; so all vanish.  Let
$h$ be a model function.  By
\cite[Corollary~1.13]{BJ25FiniteEnergy}, a positive multiple of $h$ is
a difference of bounded $\omega$-psh model functions, possibly on a
further model.  The Cauchy--Schwarz inequality
\eqref{eq:ch26-cauchy-schwarz}, applied to $u=\varphi-\psi$ and that
difference, gives
\[
 \Bigl|\int h\,\ddc u\wedge
   \omega_\varphi^{\,k}\wedge\omega_\psi^{\,n-1-k}\Bigr|^2
 \leq a_k(\varphi,\psi)\cdot(\text{a finite number})=0 .
\]
Hence the signed measure
$\ddc u\wedge\omega_\varphi^{\,k}\wedge\omega_\psi^{\,n-1-k}$, which is
the difference of two consecutive mixed measures, integrates to zero against
every model function; model functions are uniformly dense in
$C^0(X^{\An})$ by \Cref{lem:ch24-dense-tests}, so the difference vanishes.
Running $k=0,\ldots,n-1$ identifies all $n+1$ mixed measures with one
another, hence with $\MA_\omega(\varphi)=\mu$.

\cref{itm:ch28-sol:convex}.  $F_\mu$ is concave on its domain, being the
sum of the concave $E_\omega$ (\Cref{thm:extended-energy-basic}) and an
affine functional; the set where a concave function attains its maximum is
convex.  If $\varphi,\psi$ are two maximizers and
$\varphi_s=(1-s)\varphi+s\psi$, then
$F_\mu(\varphi_s)\geq(1-s)F_\mu(\varphi)+sF_\mu(\psi)=\max F_\mu$, so
equality holds; subtracting the affine term, which is exactly affine in
$s$, leaves
$E_\omega(\varphi_s)=(1-s)E_\omega(\varphi)+sE_\omega(\psi)$.

\cref{itm:ch28-sol:max}.  Membership is
\Cref{thm:ch26-e1-structure}\,\cref{itm:ch26-e1:max}.  The maximum principle
\eqref{eq:ch26-e1-maximum-principle} gives
\[
 \one_{\{\varphi>\psi\}}\MA_\omega(\chi)
 =\one_{\{\varphi>\psi\}}\MA_\omega(\varphi)
 =\one_{\{\varphi>\psi\}}\mu ,
\]
and the same identity with the two potentials exchanged gives
$\one_{\{\psi>\varphi\}}\MA_\omega(\chi)
=\one_{\{\psi>\varphi\}}\mu$.  Adding, the two measures agree on
$\{\varphi\neq\psi\}$.  Both have total mass $V$, by
\Cref{prop:e1-integrability-tail}, so their masses on the complementary set
$\{\varphi=\psi\}$ agree as well.
\end{proof}

\begin{remark}[Why \cref{itm:ch28-sol:max} stops where it does]
\label{rem:ch28-max-on-contact}
The maximum principle \eqref{eq:ch26-e1-maximum-principle} is an identity on
the Borel strict loci $\{\varphi>\psi\}$ and $\{\psi>\varphi\}$, and
it says nothing about the contact set $\{\varphi=\psi\}$, which is where
$\chi$ coincides with both potentials but where locality of the operator
gives no information: the Monge--Amp\`ere measure is not determined by the
restriction of the potential to a set that is neither open nor of full
measure.  Equality $\MA_\omega(\chi)=\mu$ would follow if one knew that
$\mu$ does not charge $\{\varphi=\psi\}$, or, in the opposite direction,
from uniqueness.  We therefore do not use \cref{itm:ch28-sol:max} in the proof
of any statement below; it is recorded because it shows that the solution set
is stable under maxima away from the contact set, which is the form in which
the corresponding complex statement is usually quoted.
\end{remark}

\begin{proposition}
\label{prop:ch28-curve-uniqueness}
Let $n=1$, that is, let $X$ be a smooth projective curve, and let
$\mu\in\mathcal M_V(X)$.  If $\varphi,\psi\in\cE^1(X,\omega)$ are
\emph{continuous} and satisfy
$\MA_\omega(\varphi)=\MA_\omega(\psi)=\mu$, then $\varphi-\psi$ is
constant.
\end{proposition}

\begin{proof}
Put $u:=\varphi-\psi$.  Since the Monge--Amp\`ere operator is affine in
dimension one,
\[
 \ddc u=\MA_\omega(\varphi)-\MA_\omega(\psi)=0.
\]
Thus $u$ is harmonic, hence subharmonic, on the smooth Berkovich curve
$X^{\An}$.  The continuous function $u$ attains a maximum $M$ because
$X^{\An}$ is compact.  Thuillier's maximum principle, recorded in
\Cref{thm:ch15-thuillier-approximation} and sourced there from
\cite[Proposition~3.1.11]{Thu05}, makes $u$ locally constant near every
point of the nonempty closed set $U:=\{u=M\}$.  Thus $U$ is also open.
Connectedness forces $U=X^{\An}$, so $u=\varphi-\psi$ is constant.

This proof uses neither a cofinal family of reduced semistable models nor a
retraction formula for Monge--Amp\`ere measures.  Those extra inputs would
unnecessarily narrow the standing hypothesis from connectedness to
geometric connectedness and from SNC to semistable models.
\end{proof}

\begin{remark}[What is missing for $n\geq2$, precisely]
\label{rem:ch28-uniqueness-gap}
The vanishing coefficients for bounded solutions do not by themselves
permit application of the relatively ample Hodge equality clause:
the mixed classes $\omega_\varphi,\omega_\psi$ can be merely nef.
The proof of \cite[Proposition~8.2]{BFJ15} first iterates
Cauchy--Schwarz to transfer vanishing to the reference Dirichlet form,
also for finite-energy potentials.  It then replaces the reference by
a positive model representative on sufficiently high models and uses
the Hodge index theorem with an ample real class to identify the
values at all divisorial points.  This supplies uniqueness in all
dimensions.  Thus the additional input is an established
intersection-theoretic argument, rather than an assumed equality
clause for arbitrary nef mixed classes.  The elementary continuous
curve proof above avoids that input.
\end{remark}

\subsection{Beyond the original scope}
\label{sec:ch28-surjectivity}

In \cite{BFJ15}, algebraizability of $X$ over a one-variable function
field was used only to prove orthogonality.  The statement
\Cref{thm:bfj-supported-cy} itself assumes orthogonality, not
algebraizability.  The latter was removed in residue characteristic zero by
\cite[Theorem~6.3.3 and Corollary~6.3.4]{BGGJKM20}.

Three further questions of scope arise naturally at this point, and we state
where each stands.

\emph{Surjectivity of $\MA_\omega$ onto $\mathcal M^1_V$.}
In the present BFJ framework, \Cref{thm:ch28-finite-energy-existence}
proves surjectivity under orthogonality by combining quantitative
concavity, energy-level compactness, and continuity of the pairing.
The conclusion is stronger than the measure-side density statement
alone.  The trivially valued theorem \cite[Theorem~12.8]{BJ22} relates
surjectivity to the envelope property and completeness of the potential
space; the corresponding theory is developed in
\Cref{chap:bj-finite-energy,chap:bj-strong-topologies}.

\emph{Other residue characteristics.}  Equal characteristic $p$, under
geometric-origin and resolution hypotheses, is
\Cref{thm:gjkm-positive-characteristic-cy}, from \cite{GJKM19} (2019);
mixed characteristic, under resolution hypotheses, is
\Cref{thm:fgk-mixed-characteristic-cy}, from \cite{FGK25} (2025).  Both
replace multiplier ideals by test ideals in the envelope step and leave the
support hypothesis on $\mu$ untouched.

\emph{Other ground fields.}  On a smooth projective geometrically
connected curve over an arbitrary complete nontrivially valued field,
\Cref{thm:curves-arbitrary-field-continuity} gives ample metric-envelope
continuity and continuous semipositive solutions for positive measures of
matching mass supported on a semistable skeleton.  The envelope input is
\cite[Theorem~1.1]{GJKM19}.  This does not by itself transport every
finite-energy or singular-psh assertion of this chapter.
In dimension $n\geq2$ over a
general complete field the situation is different in kind.  The class
$\cE^1$ can be defined, but the extension of the Monge--Amp\`ere operator
to it is left open by the 2023 source
\cite[Remark~4.4]{PS23Hybrid}.  The present chapter treats the discretely valued residue-characteristic-zero
framework, and \Cref{chap:bj-finite-energy} treats the trivially valued
framework.  The broader operator question is
\Cref{prob:frontier-general-ma}.  The conditional extensions in other
residue characteristics must be checked under their own hypotheses;
these examples of established frameworks are not an exhaustive
classification of the operator's domain.

\section{The Calabi--Yau correspondence}
\label{sec:ch28-calabi-yau}

\index{Calabi--Yau theorem!non-Archimedean, BFJ}
The name \emph{Calabi--Yau theorem} is used in this book for a statement of
the form: the Monge--Amp\`ere operator is a bijection from a class of
potentials, modulo constants, onto a class of measures.  In the complex
K\"ahler setting this is Yau's theorem, and in the finite-energy formulation
of Guedj--Zeriahi and Boucksom--Eriksson--Guedj--Zeriahi it says that
$\MA$ is a bijection from $\cE^1$ modulo constants onto the measures of
finite energy.  In the standing BFJ setting, orthogonality and
\Cref{thm:ch28-finite-energy-existence} give this finite-energy
bijection.  The original supported theorem additionally gives continuous
potentials for measures carried by one dual complex.

\subsection{Statement}
\label{sec:ch28-cy-statement}

It is convenient to assemble the results of the chapter into one statement
about the map
\begin{equation}
\label{eq:ch28-ma-map}
 \MA_\omega:\ \cE^1(X,\omega)/\mathbf R\ \longrightarrow\
 \mathcal M^1_V(X,\omega),
\end{equation}
which is well defined because $\MA_\omega(\varphi+c)=\MA_\omega(\varphi)$
and because of \Cref{prop:ch28-dual-energy-of-ma}.

\begin{theorem}
\label{thm:ch28-calabi-yau}
Retain the standing hypotheses \eqref{eq:ch28-standing} and assume that
$\omega$ has the orthogonality property.  Then the map
\eqref{eq:ch28-ma-map} has the following properties.
\begin{enumerate}
\item\label{itm:ch28-cy:into}
 It is well defined: $\MA_\omega(\varphi)$ has mass $V$ and finite dual
 energy, and
 \[
  E^\vee_\omega(\MA_\omega(\varphi))
  =E_\omega(\varphi)-\int\varphi\,\MA_\omega(\varphi).
 \]
\item\label{itm:ch28-cy:variational}
 For $\varphi\in\cE^1(X,\omega)$ and $\mu\in\mathcal M_V(X)$ such
 that $\int\varphi\,\rd\mu> -\infty$, the following are equivalent:
 $\MA_\omega(\varphi)=\mu$; $\varphi$
 maximizes $F_\mu$; $\mu$ is a supergradient of $E_\omega$ at
 $\varphi$; and equality holds in the Fenchel inequality
 \eqref{eq:measure-fenchel-energy} at the pair $(\varphi,\mu)$.
\item\label{itm:ch28-cy:onto}
 Its image is all of $\mathcal M^1_V$.  If $\mu$ is carried by one
 dual complex, its preimage contains a continuous representative.
\item\label{itm:ch28-cy:injective}
 It is injective on $\cE^1/\mathbf R$ in every dimension by
 \textup{(I2)}.  A separate elementary proof for continuous potentials
 on a curve is \Cref{prop:ch28-curve-uniqueness}.
\item\label{itm:ch28-cy:usc}
 It is monotone-continuous: a decreasing net
 $\varphi_\alpha\searrow\varphi$ in $\cE^1$ has
 $\MA_\omega(\varphi_\alpha)\to\MA_\omega(\varphi)$ weakly.
\end{enumerate}
\end{theorem}

\begin{proof}
\cref{itm:ch28-cy:into} is \Cref{prop:ch28-dual-energy-of-ma}.
\cref{itm:ch28-cy:variational} is
\Cref{prop:variational-maximizer-solves} together with
\Cref{cor:finite-energy-variational-characterization} and
\Cref{cor:ch27-euler-lagrange-supergradient}; the equality clause in the
Fenchel inequality is equivalently attainment of the supremum
by \Cref{prop:ch28-sup-is-dual-energy}; it does not require assuming
that the prescribed measure is already in the image.
\cref{itm:ch28-cy:onto} is
\Cref{thm:ch28-finite-energy-existence}, with continuity from
\Cref{thm:bfj-supported-cy}.  \cref{itm:ch28-cy:injective} is
\cite[Proposition~8.2]{BFJ15}.  \cref{itm:ch28-cy:usc} is
\Cref{prop:e1-monotone-continuity}.
\end{proof}

\begin{warning}{Regularity and topology of the correspondence}
Finite dual energy characterizes the image, but does not imply continuity
of the solution; \Cref{ex:ch28-unbounded-dual-energy} gives an unbounded
solution.  This theorem does not assert a homeomorphism for the strong
topologies.  Maximizing-sequence convergence is
\Cref{thm:ch28-finite-energy-existence}; fixed-skeleton solution stability
is \Cref{prop:fixed-skeleton-stability}.
\end{warning}

\begin{remark}[The correspondence in one line]
\label{rem:ch28-correspondence-line}
The content of \Cref{thm:ch28-calabi-yau} may be summarized as the
commutativity of the Legendre picture:
\[
 \varphi\ \longmapsto\ \MA_\omega(\varphi),
 \qquad
 E_\omega(\varphi)-E^\vee_\omega\bigl(\MA_\omega(\varphi)\bigr)
 =\int\varphi\,\MA_\omega(\varphi),
\]
the second identity being \eqref{eq:ch28-dual-energy-of-ma} rearranged.  It
is the exact analogue, for the pair (energy, dual energy), of the
Fenchel--Young equality \eqref{eq:fenchel-young-concave} at a supergradient,
and of the entropy computation of \Cref{ex:legendre-entropy-pair}, where the
dual of the free energy is the relative entropy and the equality case is the
Gibbs measure.
\end{remark}

\subsection{The correspondence, item by item}
\label{sec:ch28-cy-table}

\begin{center}
\small
\begin{tabular}{p{.24\textwidth}|p{.24\textwidth}|p{.20\textwidth}|p{.22\textwidth}}
\toprule
\textbf{Potential side}&\textbf{Measure side}&\textbf{Status}
 &\textbf{Reference}\\
\midrule
$\varphi\in\cE^1(X,\omega)$
 &$\MA_\omega(\varphi)\in\mathcal M^1_V$
 &proved
 &\Cref{prop:ch28-dual-energy-of-ma}\\
$E_\omega(\varphi)$
 &$E^\vee_\omega(\MA_\omega\varphi)$
 &Legendre pair
 &\eqref{eq:ch28-dual-energy-of-ma}\\
$\varphi$ bounded, model
 &$\MA_\omega(\varphi)$ atomic on divisorial points
 &proved
 &\Cref{prop:ch24-model-support}\\
$\varphi$ continuous
 &$\mu$ carried by one dual complex
 &sufficient, not necessary
 &\Cref{thm:bfj-supported-cy},
  \Cref{ex:ch28-unbounded-dual-energy}\\
$\varphi$ unbounded, finite energy
 &$\mu$ of finite dual energy, support not skeletal
 &example
 &\Cref{ex:ch28-unbounded-dual-energy}\\
$\varphi$ of infinite energy
 &non-pluripolar measure can lose mass
 &example
 &\Cref{ex:ch26-infinite-energy}\\
--- (no potential)
 &$\delta_x$, $x$ rational: $E^\vee_\omega=+\infty$
 &example
 &\Cref{ex:ch28-infinite-dual-energy}\\
decreasing net $\varphi_\alpha\searrow\varphi$
 &weak convergence of measures
 &proved
 &\Cref{prop:e1-monotone-continuity}\\
maximizing sequence for $F_\mu$
 &strong convergence to $\mu$
 &proved modulo two quoted estimates
 &\Cref{thm:finite-energy-measure-stability}\\
$\cE^1/\mathbf R$ with transported $\delta_\omega$
 &$\mathcal M^1_V$ complete for $\delta_\omega$
 &complete via the bijection
 &\Cref{thm:finite-energy-measure-completion},
  \Cref{rem:ch28-not-a-metric-on-potentials}\\
\bottomrule
\end{tabular}
\end{center}

Skeletal support is sufficient for continuity of the potential, and
\Cref{ex:ch28-unbounded-dual-energy} also supplies a continuous
nonskeletal example showing that it is not necessary.  Uniform
integrability on all of $\PSH_0$ is stronger than finite dual energy.
The transported complete quasi-metric is not identified here with the
potential gauges of \Cref{chap:bfj-energy}.

\begin{remark}[Comparison with the complex and the trivially valued cases]
\label{rem:ch28-cy-comparison}
The complex finite-energy variational theory gives a bijection modulo
constants in an ample K\"ahler class; see \cite{BBGZ13} and the big-class
extensions \cite{BEGZ10,BBEGZ16}.  The BFJ argument here also yields
a finite-energy bijection, but uses its own global orthogonality theorem,
skeletal compactness, and quantitative energy estimates.  It does not
import a local Bedford--Taylor proof, or assert compactness of bounded
sets in a strong metric.  The trivially valued theorem
\cite[Theorem~12.8]{BJ22} and its precise envelope and completeness
equivalences are treated in
\Cref{chap:bj-finite-energy,chap:bj-strong-topologies}.
\end{remark}

\section{Stability with respect to the measure}
\label{sec:ch28-stability}

Three stability statements are proved, in increasing strength and decreasing
generality.  \Cref{thm:finite-energy-measure-stability} says that the
measures of a model maximizing sequence converge strongly for every
finite-energy measure.  \Cref{thm:ch28-finite-energy-existence} additionally
proves strong convergence of its normalized potentials.
\Cref{prop:fixed-skeleton-stability} says that on a fixed dual complex the
solutions converge; it applies only to admissible measures but controls
potentials.  \Cref{thm:ch28-curve-stability} is a quantitative estimate with
explicit constants, available on curves.

\subsection{Strong convergence of a maximizing sequence}
\label{sec:ch28-maximizing-net}

\begin{theorem}
\label{thm:finite-energy-measure-stability}
Assume orthogonality and let $\mu\in\mathcal M^1_V$.  If
$\varphi_j$ is a model sequence computing the supremum in
\eqref{eq:dual-energy-measure}, then
\[
                     \MA_\omega(\varphi_j)\longrightarrow\mu
\]
strongly.  In particular, strong measure convergence implies weak
convergence, but the converse need not hold.
\end{theorem}
\begin{proof}
Put $\mu_j:=\MA_\omega(\varphi_j)$ and measure the defect of maximality by
\[
 \varepsilon_j:=E_\omega^\vee(\mu)-
 \left(E_\omega(\varphi_j)-\int\varphi_j\,\rd\mu\right).
\]
Thus $\varepsilon_j\geq0$ and $\varepsilon_j\to0$.  We reproduce the
maximizing-sequence argument of
\cite[Theorem~2.22]{BJ25FiniteEnergy}, retaining our mass-$V$
normalization.

Let $f$ be a model function.  By
\cite[Corollary~1.13]{BJ25FiniteEnergy}, write $f=f_+-f_-$ with
$f_\pm$ model functions that are $C\omega$-psh for some $C>0$; obtaining
these representatives can require a further model.  Uniform
differentiability of equilibrium energy,
\cite[Proposition~2.18, equation~(2.19)]{BJ25FiniteEnergy}, therefore gives
a constant
$A_f$ independent of $j$ such that, for $0<t\leq1$,
\begin{equation}\label{eq:maximizing-uniform-differentiability}
 \Bigl|\,\widetilde E_\omega(\varphi_j+tf)-E_\omega(\varphi_j)
   -t\int f\,\rd\mu_j\,\Bigr|
 \leq A_f t^2.
\end{equation}
The estimate is two-sided, and it is the \emph{lower} bound
$\widetilde E_\omega(\varphi_j+tf)\geq
  E_\omega(\varphi_j)+t\int f\,\rd\mu_j-A_ft^2$
that is used below; an upper bound on
$\widetilde E_\omega(\varphi_j+tf)$ is already available from the
definition of $E^\vee_\omega(\mu)$ and would not combine with it.
Here
$\widetilde E_\omega(g):=E_\omega(P_\omega g)$ is the equilibrium-energy
extension.  Since $P_\omega(\varphi_j+tf)\leq\varphi_j+tf$, the definition
of $E_\omega^\vee(\mu)$ yields
\[
 \widetilde E_\omega(\varphi_j+tf)
 \leq E_\omega^\vee(\mu)+
      \int(\varphi_j+tf)\,\rd\mu
 =E_\omega(\varphi_j)+\varepsilon_j+
      t\int f\,\rd\mu.
\]
Combining this with
\eqref{eq:maximizing-uniform-differentiability} gives
\[
 t\int f\,\rd\mu_j
 \leq t\int f\,\rd\mu+\varepsilon_j+A_ft^2.
\]
First let $j\to\infty$ and then $t\downarrow0$; applying the same argument
to $-f$ gives
\[
                     \int f\,\rd\mu_j\longrightarrow
                     \int f\,\rd\mu.
\]
Model functions are uniformly dense in $C^0(X^{\An})$ and all these
measures have mass $V$, so $\mu_j\to\mu$ weakly.

It remains to prove convergence of the Legendre energies.  Fenchel equality
at a model Monge--Amp\`ere measure gives
\[
 E_\omega^\vee(\mu_j)
 =E_\omega(\varphi_j)-\int\varphi_j\,\rd\mu_j.
\]
The mixed-energy Cauchy--Schwarz estimate used in
\Cref{thm:finite-energy-measure-completion} supplies constants
$C>0$ and $\alpha=2^{-n}$, independent of $j,k$, such that
\begin{equation}\label{eq:maximizing-mixed-energy-control}
 \left|\int\varphi_j\,\rd(\mu_j-\mu_k)\right|
 \leq C\max\{\varepsilon_j,\varepsilon_k\}^{\alpha}.
\end{equation}
Indeed, \cite[Proposition~2.2, equation~(2.7)]{BJ25FiniteEnergy}
and quasi-symmetry bound $J_\omega(0,\varphi_j)$ on a cofinal tail by a
constant depending on $E_\omega^\vee(\mu)$ and the bounded defects.
The same equation bounds $J_\omega(\varphi_j,\varphi_k)$ by a constant
times $\varepsilon_j+\varepsilon_k$.  The integral estimate
\cite[Theorem~1.33, equation~(1.37)]{BJ25FiniteEnergy} then gives
\eqref{eq:maximizing-mixed-energy-control}, after absorbing a fixed factor into $C$.  More explicitly, the integral estimate first
bounds the left-hand side by a constant times
$\max\{\varepsilon_j,\varepsilon_k\}^{1/2}$; on the tail with defects
at most one this is at most a constant times the displayed
$\alpha$-power, since $\alpha\leq1/2$.
For fixed $j$, weak convergence of $\mu_k$ lets
$k\to\infty$ in the left-hand side of
\eqref{eq:maximizing-mixed-energy-control}; then
\[
 \left|\int\varphi_j\,\rd(\mu_j-\mu)\right|
 \leq C\varepsilon_j^\alpha.
\]
Consequently
\[
 E_\omega^\vee(\mu_j)-E_\omega^\vee(\mu)
 =-\varepsilon_j+
   \int\varphi_j\,\rd(\mu-\mu_j)\longrightarrow0.
\]
This is strong convergence by definition.  Its weak component is therefore
automatic, while weak convergence alone does not force the last displayed
energy convergence.
\end{proof}

This is the measure-side density mechanism.  The normalized potential
convergence in the standing setting follows separately from
\Cref{thm:ch28-finite-energy-existence}, using its compact energy level
and continuous pairing.  Mere weak convergence of arbitrary measures
does not control their dual energies.

\begin{remark}[Maximizing sequences and nonmetrizable compactness]
\label{rem:ch28-maximizing-sequence-exists}
A model maximizing sequence exists because the defining supremum is a
finite supremum of real numbers.  This fact requires no metrizability
of the potential space and does not supply a cofinal model
regularization sequence.  Compactness initially supplies convergent
subnets, rather than a sequential-compactness theorem.  Under
orthogonality, the uniform energy bound and uniqueness of all cluster
points in \Cref{thm:ch28-finite-energy-existence} prove convergence of
the entire normalized maximizing sequence.  No sequential-compactness
claim is needed.
\end{remark}

\begin{remark}[Which parts of the proof are quoted]
\label{rem:ch28-stability-quoted}
The uniform differentiability estimate is
\cite[Proposition~2.18, equation~(2.19)]{BJ25FiniteEnergy}; its model-test
decomposition uses Corollary~1.13 of that source.
The mixed-energy bound uses Proposition~2.2(2.7) to bound both the
reference Dirichlet energies and the pairwise defects, followed by
Theorem~1.33(1.37).  These quantitative inputs are stronger than
pointwise differentiability or a single Cauchy--Schwarz inequality.
The remaining steps above reproduce the maximizing-sequence argument of
\cite[Theorem~2.22]{BJ25FiniteEnergy}.
\end{remark}

\subsection{Stability of the solutions on a fixed dual complex}
\label{sec:ch28-fixed-skeleton}

\begin{proposition}
\label{prop:fixed-skeleton-stability}
Assume orthogonality.  Let $\mu_i,\mu$ have mass $V$, be supported on one fixed dual complex
$\Delta$, and satisfy $\mu_i\to\mu$ weakly.  Let $\varphi_i,\varphi$
be the solutions of \Cref{thm:bfj-supported-cy}, normalized by
supremum zero.  Then $\varphi_i\to\varphi$ in the BFJ weak topology and
uniformly on $\Delta$; their energies converge as well.
\end{proposition}

\begin{proof}
The normalized psh restrictions form a compact subset of $C^0(\Delta)$.
Weak convergence of the measures is therefore uniform on this compact test
family.  Hence their finite linear terms converge uniformly; on
$\cE^1$ the functionals $F_{\mu_i}$ differ uniformly by the same bound,
and outside $\cE^1$ they all take the extended value $-\infty$.  Every cluster point of their maximizers maximizes
$F_\mu$.  Uniqueness identifies it with $\varphi$, so the whole net
converges weakly; skeletal equicontinuity upgrades the restrictions to
uniform convergence.  The identity
$E_\omega(\varphi_i)=F_{\mu_i}(\varphi_i)+
\int\varphi_i\,\rd \mu_i$ gives convergence of energies.
\end{proof}

We expand the four assertions of that proof, since each rests on a statement
that has to be located.

\emph{Uniform convergence on a compact test family.}  Let
$\mathcal R:=\{\varphi|_\Delta:\varphi\in\PSH_0(X,\omega)\}$, a compact
subset of $C^0(\Delta)$ by compactness of $\PSH_0$ and continuity
of restriction; \Cref{cor:ch22-arzela-ascoli} gives the same
equicontinuity criterion.  Continuous functions on the closed subset
$\Delta$ extend to $X^{\An}$ by Tietze's theorem, so ambient weak
convergence tests all of $C^0(\Delta)$.  A net of
measures on $\Delta$ of bounded mass converging weakly converges uniformly
on a compact subset of $C^0(\Delta)$: given $\varepsilon>0$, cover
$\mathcal R$ by finitely many balls of radius $\varepsilon/(3V)$ with
centres $g_1,\ldots,g_N$, choose an index beyond which
$|\int g_m\,\rd\mu_i-\int g_m\,\rd\mu|<\varepsilon/3$ for every
$m\leq N$, and estimate an arbitrary $g\in\mathcal R$ through the
nearest centre.  Hence
\begin{equation}
\label{eq:ch28-uniform-on-test-family}
 \sup_{\varphi\in\PSH_0(X,\omega)}
 \Bigl|\int\varphi\,\rd\mu_i-\int\varphi\,\rd\mu\Bigr|
 \ \longrightarrow\ 0 ,
\end{equation}
so on $\cE^1\cap\PSH_0$ the difference
$F_{\mu_i}-F_\mu$ tends uniformly to zero.  On
$\PSH_0\setminus\cE^1$ all these functionals equal $-\infty$; their
pairings are finite there, but their energies are not.  Subsequent
comparisons use the uniform bound on the linear terms, without
subtracting two infinite values.

\emph{Cluster points are maximizers.}  Uniform convergence of an upper
semicontinuous family transfers maximizers to maximizers.  Precisely, let
$\varphi_{i_\beta}\to\chi$ be a convergent subnet, which exists by
\Cref{cor:ch22-net-compactness}.  There is a common skeletal pairing
bound $M$, and $F_{\mu_i}(\varphi_i)\geq0$ gives
$E_\omega(\varphi_i)\geq-M$.  Thus $\chi\in\mathcal K_M$ and all
maximizer values in the comparison below are finite.
For a competitor of infinite energy the inequality is immediate;
for any competitor
$\psi\in\PSH_0(X,\omega)$,
\[
 F_\mu(\chi)\ \geq\ \varlimsup_\beta F_{\mu}(\varphi_{i_\beta})
 =\varlimsup_\beta F_{\mu_{i_\beta}}(\varphi_{i_\beta})
 \ \geq\ \varlimsup_\beta F_{\mu_{i_\beta}}(\psi)=F_\mu(\psi),
\]
the first inequality by upper semicontinuity of $F_\mu$
(Step~1 of the proof of \Cref{thm:ch28-existence-admissible}), the two
equalities by \eqref{eq:ch28-uniform-on-test-family}, and the middle
inequality by maximality of $\varphi_{i_\beta}$ for $F_{\mu_{i_\beta}}$.

\emph{Identification of the cluster point.}  Orthogonality is in force,
since \Cref{thm:bfj-supported-cy} assumes it and the solutions
$\varphi_i,\varphi$ are the ones it produces.  By
\Cref{prop:variational-maximizer-solves}, $\chi$ solves
$\MA_\omega(\chi)=\mu$; by \Cref{thm:bfj-supported-regularity} it is
continuous; and by the uniqueness clause of
\Cref{thm:bfj-supported-cy}, together with the normalization
$\sup\chi=0$, we get $\chi=\varphi$.  Thus $\varphi$ is the only
possible limit of a convergent subnet.  In a compact space this forces
convergence of the whole net: if $(\varphi_i)$ did not converge to
$\varphi$, some cofinal subnet would avoid a fixed neighbourhood of
$\varphi$, and by \Cref{prop:net-criterion}\,(3) that subnet would itself
have a convergent subnet, whose limit would be a point of
$\PSH_0(X,\omega)$ different from $\varphi$ --- impossible, because the
argument above applies to every convergent subnet and because limits in
$\PSH_0(X,\omega)$ are unique: the BFJ topology is that of uniform
convergence on dual complexes, and a psh function is determined by its
skeletal restrictions by \Cref{thm:bfj-skeletal-characterization}, so the
topology is Hausdorff and \Cref{prop:net-criterion}\,(4) applies.

\emph{Convergence of energies.}  From
$E_\omega(\varphi_i)=F_{\mu_i}(\varphi_i)+\int\varphi_i\,\rd\mu_i$, the
first summand converges to $F_\mu(\varphi)$ --- it equals
$E^\vee_\omega(\mu_i)$ by
\Cref{thm:ch28-existence-admissible}\,\cref{itm:ch28-exist:value}, and
\eqref{eq:ch28-uniform-on-test-family} shows that
$\sup F_{\mu_i}\to\sup F_\mu$ --- while the second converges by
\Cref{lem:joint-weak-uniform}, applied on $\Delta$ to the uniformly
convergent restrictions $\varphi_i|_\Delta$ and the weakly convergent
measures $\mu_i$.  In particular $\mu_i\to\mu$ strongly in the sense of
\Cref{def:ch28-strong-measure-convergence}, and
$\varphi_i\to\varphi$ strongly in the sense of \Cref{sec:ch26-strong}.

The proposition intentionally makes no claim of uniform convergence on all
of $X^{\An}$.  In higher dimension a semipositive metric solving a
skeleton-supported equation need not factor through the retraction to that
skeleton.

\begin{warning}{Weak convergence of the measures is enough here, and only
here}
\Cref{prop:fixed-skeleton-stability} deduces strong convergence of the
measures from \emph{weak} convergence.  That is legitimate because all the
measures are carried by one compact set $\Delta$ on which the normalized
psh functions form an equicontinuous family; the implication is false in
general, as \Cref{ex:ch28-infinite-dual-energy} shows, where a weakly
convergent net of Monge--Amp\`ere measures has a limit of infinite dual
energy.  The hypothesis that the \emph{same} dual complex carries every
$\mu_i$ cannot be weakened to ``each $\mu_i$ is carried by some dual
complex'': the union of an infinite family of dual complexes need not lie in
one of them, and \Cref{prop:ch28-admissible-calculus}\,%
\cref{itm:ch28-admcalc:union} covers only finite unions.
\end{warning}

\subsection{A quantitative estimate on curves}
\label{sec:ch28-curve-stability}

On a curve the whole variational problem is a Dirichlet problem on a metric
graph, and stability becomes an explicit inequality with a constant
depending only on the total length of the graph.  We prove it here; the
worked numbers are in \Cref{sec:ch28-curves}.

Throughout this subsection $X$ is a smooth projective geometrically
connected curve, $\mathcal X$ is a split regular semistable model, with nodes of residue
degree one, and skeleton
$\Gamma$ carrying its modulus lengths, and $r:X^{\An}\to\Gamma$ is the
retraction.  The form $\omega$ comes from a relatively nef model on
$\mathcal X$ of a line bundle $L$ of degree $V>0$; hence its curvature
measure $\mu_0=\MA_\omega(0)$ is carried by $\Gamma$ and has mass $V$.
These additional curve hypotheses ensure the unweighted formulas
below; compare \Cref{sec:ch31-curve-setup,prop:graph-ma-solution}.
They are not part of the general standing hypotheses of this chapter.  We write
\begin{equation}
\label{eq:ch28-dirichlet-form}
 D_\Gamma(f):=\int_\Gamma|f'|^2\,\rd s\ \in[0,+\infty],
 \qquad
 L(\Gamma):=\text{total length of }\Gamma ,
\end{equation}
for the Dirichlet energy of a piecewise affine, or more generally absolutely
continuous, function on $\Gamma$.

\begin{lemma}
\label{lem:ch28-graph-sobolev}
Let $\Gamma$ be a compact connected metric graph of total length
$L(\Gamma)$ and let $f$ be continuous on $\Gamma$ and piecewise
$C^1$.  Then
\begin{equation}
\label{eq:ch28-graph-sobolev}
 \operatorname{osc}_\Gamma f
 \ \leq\ L(\Gamma)^{1/2}\,D_\Gamma(f)^{1/2}.
\end{equation}
\end{lemma}

\begin{proof}
Let $x,y\in\Gamma$ and let $\gamma$ be an injective path from $x$ to
$y$ inside $\Gamma$; such a path exists because $\Gamma$ is connected,
and its length $\ell(\gamma)$ is at most $L(\Gamma)$ because an
injective path traverses each edge at most once.  By the fundamental theorem
of calculus along $\gamma$ and the Cauchy--Schwarz inequality,
\[
 |f(x)-f(y)|
 \ \leq\ \int_\gamma|f'|\,\rd s
 \ \leq\ \ell(\gamma)^{1/2}
   \Bigl(\int_\gamma|f'|^2\,\rd s\Bigr)^{1/2}
 \ \leq\ L(\Gamma)^{1/2}D_\Gamma(f)^{1/2}.
\]
Taking the supremum over $x,y$ gives \eqref{eq:ch28-graph-sobolev}.
\end{proof}

\begin{lemma}
\label{lem:ch28-green-dirichlet}
Let $\mu_1,\mu_2$ be positive Radon measures of mass $V$ on $\Gamma$, and
let $u_1,u_2$ be the Green-kernel solutions supplied by
\Cref{prop:graph-ma-solution}.  Then $u:=u_1-u_2$ is absolutely continuous
on every edge, $D_\Gamma(u)<+\infty$, and
\begin{equation}
\label{eq:ch28-green-dirichlet}
 D_\Gamma(u)=-\int_\Gamma u\,\rd(\mu_1-\mu_2).
\end{equation}
\end{lemma}

\begin{proof}
Choose atomic rational measures $\mu_{a,m}\to\mu_a$ weakly, of mass $V$,
as in the proof of \Cref{prop:graph-ma-solution}, and use one base point to
normalize all Green potentials $u_{a,m}$ and $u_a$.  That proof gives
$u_{a,m}\to u_a$ uniformly.  Put $w_m:=u_{1,m}-u_{2,m}$ and
$\nu_m:=\mu_{1,m}-\mu_{2,m}$.  The functions $w_m$ are piecewise affine,
so \eqref{eq:graph-green} gives
\[
 D_\Gamma(w_m-w_k)
 =-\int_\Gamma(w_m-w_k)\,\rd(\nu_m-\nu_k)
 \leq4V\lVert w_m-w_k\rVert_\infty.
\]
Thus the edge derivatives $w_m'$ are Cauchy in the direct sum of the edge
$L^2$ spaces.  Their limit is the weak edge derivative of the uniform
limit $u$, so $u$ is edgewise absolutely continuous and
$D_\Gamma(w_m)\to D_\Gamma(u)<+\infty$.  Finally, another application of
\eqref{eq:graph-green} and the joint weak--uniform convergence lemma
\Cref{lem:joint-weak-uniform} give
\[
 D_\Gamma(u)
 =\lim_mD_\Gamma(w_m)
 =-\lim_m\int_\Gamma w_m\,\rd\nu_m
 =-\int_\Gamma u\,\rd(\mu_1-\mu_2),
\]
which proves \eqref{eq:ch28-green-dirichlet}.  Adding constants to either
solution changes neither side because $\mu_1-\mu_2$ has mass zero.
\end{proof}

\begin{theorem}
\label{thm:ch28-curve-stability}
\index{stability!with respect to the measure}
Let $\mu_1,\mu_2$ be positive measures of mass $V$ carried by
$\Gamma$, and let $u_1,u_2$ be the corresponding solutions of
$\MA_\omega(u_a)=\mu_a$ furnished by
\Cref{prop:graph-ma-solution}, normalized by
$\sup_{X^{\An}}u_a=0$.  Put $u:=u_1-u_2$ and let
$\lVert\mu_1-\mu_2\rVert$ denote the total variation of the signed measure
$\mu_1-\mu_2$ on $\Gamma$.  Then
\begin{align}
\label{eq:ch28-curve-stability-dirichlet}
 D_\Gamma(u)&\ \leq\ \tfrac14L(\Gamma)\,
       \lVert\mu_1-\mu_2\rVert^2,\\
\label{eq:ch28-curve-stability-sup}
 \lVert u_1-u_2\rVert_{C^0(X^{\An})}
   &\ \leq\ \tfrac12L(\Gamma)\,
       \lVert\mu_1-\mu_2\rVert ,\\
\label{eq:ch28-curve-stability-dual}
 \bigl|E^\vee_\omega(\mu_1)-E^\vee_\omega(\mu_2)\bigr|
   &\ \leq\ \tfrac18L(\Gamma)\lVert\mu_1-\mu_2\rVert^2
      +\tfrac12L(\Gamma)^{1/2}\lVert\mu_1-\mu_2\rVert
        \min_aD_\Gamma(u_a)^{1/2}.
\end{align}
In particular $\mu\mapsto u_\mu$ is Lipschitz from the total-variation
norm to the uniform norm, with constant $L(\Gamma)/2$, and
$E^\vee_\omega$ is locally Lipschitz on the set of measures carried by
$\Gamma$.
\end{theorem}

\begin{proof}
\emph{Step 1: the Dirichlet identity.}  By
\eqref{eq:curve-ma-laplacian} the two solutions satisfy
$\Delta_\Gamma u_a=\mu_a-\mu_0$ on $\Gamma$, so
$\Delta_\Gamma u=\mu_1-\mu_2$.  The finite-energy identity
\Cref{lem:ch28-green-dirichlet} gives
\begin{equation}
\label{eq:ch28-stability-dirichlet-identity}
 D_\Gamma(u)=-\int_\Gamma u\,\rd(\mu_1-\mu_2),
\end{equation}
all terms being finite.  Continuity alone would not justify the derivative
integral; this is why \Cref{lem:ch28-green-dirichlet} was separated out.

\emph{Step 2: the estimate.}  Since $\mu_1-\mu_2$ has total mass zero, the
right-hand side of
\eqref{eq:ch28-stability-dirichlet-identity} is unchanged when $u$ is
replaced by $u-c$ for any constant $c$; choosing $c$ to be the
midpoint of the range of $u$ on $\Gamma$ gives
\[
 D_\Gamma(u)\ \leq\ \tfrac12\operatorname{osc}_\Gamma(u)\,
 \lVert\mu_1-\mu_2\rVert
 \ \leq\ \tfrac12L(\Gamma)^{1/2}D_\Gamma(u)^{1/2}
   \lVert\mu_1-\mu_2\rVert ,
\]
using \Cref{lem:ch28-graph-sobolev}.  If $D_\Gamma(u)=0$ there is nothing
to prove; otherwise dividing by $D_\Gamma(u)^{1/2}$ gives
$D_\Gamma(u)^{1/2}\leq\tfrac12L(\Gamma)^{1/2}
\lVert\mu_1-\mu_2\rVert$, whence
\eqref{eq:ch28-curve-stability-dirichlet}, and
\[
 \operatorname{osc}_\Gamma(u)\ \leq\ L(\Gamma)^{1/2}D_\Gamma(u)^{1/2}
 \ \leq\ \tfrac12 L(\Gamma)\lVert\mu_1-\mu_2\rVert .
\]

\emph{Step 3: from the graph to $X^{\An}$.}  By
\Cref{prop:graph-ma-solution} both solutions factor through the retraction,
$u_a=U_a\circ r$ with $U_a$ continuous on $\Gamma$; this is the
one-dimensional retraction statement, and it is false in higher dimension.
Hence $u=U\circ r$ with $U=U_1-U_2$, and
$\lVert u\rVert_{C^0(X^{\An})}=\lVert U\rVert_{C^0(\Gamma)}$.  The
normalization $\sup_{X^{\An}}u_a=0$ and
\Cref{prop:ch21-sup-on-skeleton} give $\max_\Gamma U_a=0$ for
$a=1,2$; consequently $U$ attains a value $\leq0$ at a maximum point
of $U_2$ and a value $\geq0$ at a maximum point of $U_1$, so
$\lVert U\rVert_{C^0(\Gamma)}\leq\operatorname{osc}_\Gamma U$.  Combining
with Step~2 gives \eqref{eq:ch28-curve-stability-sup}.

\emph{Step 4: the dual energies.}  By
\Cref{prop:ch28-curve-dual-energy} below,
$E^\vee_\omega(\mu_a)=\tfrac12D_\Gamma(u_a)$.  Therefore
\[
 \begin{aligned}
 \bigl|E^\vee_\omega(\mu_1)-E^\vee_\omega(\mu_2)\bigr|
 &=\tfrac12\bigl|D_\Gamma(u_1)-D_\Gamma(u_2)\bigr|\\
 &=\tfrac12\bigl|D_\Gamma(u_1)^{1/2}-D_\Gamma(u_2)^{1/2}\bigr|\\
 &\qquad\cdot
   \bigl(D_\Gamma(u_1)^{1/2}+D_\Gamma(u_2)^{1/2}\bigr).
 \end{aligned}
\]
The first factor is at most $D_\Gamma(u_1-u_2)^{1/2}$, because
$f\mapsto D_\Gamma(f)^{1/2}$ is a seminorm and hence satisfies the reverse
triangle inequality; and by
\eqref{eq:ch28-curve-stability-dirichlet} it is at most
$\tfrac12L(\Gamma)^{1/2}\lVert\mu_1-\mu_2\rVert$.  For the second factor, write
$D_\Gamma(u_1)^{1/2}+D_\Gamma(u_2)^{1/2}
\leq 2\min_aD_\Gamma(u_a)^{1/2}+D_\Gamma(u)^{1/2}$, again by the triangle
inequality, and bound $D_\Gamma(u)^{1/2}$ as before.  This is
\eqref{eq:ch28-curve-stability-dual}.
\end{proof}

\begin{remark}[What the constant measures, and what replaces it]
\label{rem:ch28-curve-constant}
The coefficient $L(\Gamma)/2$ in
\eqref{eq:ch28-curve-stability-sup} uses the total length of the skeleton in
its modulus metric.  Refining the semistable model merely subdivides edges
and changes neither this length nor the coefficient, so the sharp factor
$1/2$ is stable under subdivision.  Under a finite scalar extension one
must use the actual modulus metric on the resulting skeleton together with
the measure normalization of \Cref{thm:base-change-ma}; no blanket
multiplication of lengths by the ramification index is asserted when the
absolute value is kept fixed.

In dimension $n\geq2$, the measure-potential pairing is the mixed
energy $I(\varphi_1,\varphi_2)$, a sum of $n$ nonnegative Dirichlet
coefficients.  It need not vanish when the prescribed measures differ.
The graph Sobolev estimate has no higher-dimensional counterpart
established here, so this identity alone does not provide the same
supremum-norm bound.  What survives in
general is \Cref{prop:fixed-skeleton-stability}, which is qualitative, and
\Cref{thm:finite-energy-measure-stability}, which is quantitative but
concerns the measures rather than the potentials.
\end{remark}

\section{Worked calculations}
\label{sec:ch28-curves}

The calculations below are carried to explicit numbers.  Together they test
every normalization of the chapter: the mass $V$ in
\eqref{eq:dual-energy-measure}, the sign of the Legendre transform, the
identity \eqref{eq:ch28-dual-energy-of-ma}, the coercivity constant
\eqref{eq:ch28-skeleton-M-bound}, and the stability estimate
\eqref{eq:ch28-curve-stability-sup}.  All of them except
\Cref{ex:ch28-divisorial-green} take place on a metric graph, where the
theory is a Hilbert-space computation.

\subsection{The dual energy on a curve skeleton}
\label{sec:ch28-curve-dual-energy}

\begin{proposition}
\label{prop:ch28-curve-dual-energy}
\index{Dirichlet energy!on a metric graph}\index{Green function!energy of a measure}
Retain the geometrically connected semistable curve setting fixed at the
start of \Cref{sec:ch28-curve-stability}.  Let $\mu$ be a positive measure
of mass $V$ carried by $\Gamma$, and let $u$ be the solution of
$\MA_\omega(u)=\mu$ supplied by \Cref{prop:graph-ma-solution}.  Then
\begin{equation}
\label{eq:ch28-curve-dual-energy}
 E^\vee_\omega(\mu)
 =\tfrac12\int_\Gamma|u'|^2\,\rd s
 =\tfrac12 D_\Gamma(u)
 =\tfrac12\int_\Gamma u\,\rd(\mu_0-\mu)
 \ \in[0,+\infty).
\end{equation}
In particular $E^\vee_\omega(\mu)=0$ if and only if $\mu=\mu_0$.
\end{proposition}

\begin{proof}
The function $u$ is continuous, hence bounded, hence lies in
$\cE^1(X,\omega)$ by
\Cref{thm:ch26-e1-structure}\,\cref{itm:ch26-e1:bounded}, and
$\MA_\omega(u)=\mu$.  \Cref{prop:ch28-dual-energy-of-ma} therefore gives
\[
 E^\vee_\omega(\mu)=E_\omega(u)-\int u\,\rd\mu .
\]
Choose the atomic approximants $u_m$ used in the proof of
\Cref{lem:ch28-green-dirichlet}, now with $\mu_1=\mu$,
$\mu_2=\mu_0$ and $u_2=0$.  Formula \eqref{eq:ch26-graph-energy} applies
to $u_m$.  Uniform convergence, the energy Lipschitz estimate
\eqref{eq:ch26-energy-lipschitz}, and the strong derivative convergence proved in
\Cref{lem:ch28-green-dirichlet} give
\[
 E_\omega(u)=\int_\Gamma u\,\rd\mu_0-\tfrac12 D_\Gamma(u).
\]
Hence
\[
 E^\vee_\omega(\mu)
 =\int_\Gamma u\,\rd(\mu_0-\mu)-\tfrac12D_\Gamma(u).
\]
The identity \eqref{eq:ch28-green-dirichlet}, again with
$\mu_1=\mu$ and $\mu_2=\mu_0$, gives
$\int_\Gamma u\,\rd(\mu_0-\mu)=D_\Gamma(u)$.  Substituting,
\[
 E^\vee_\omega(\mu)=D_\Gamma(u)-\tfrac12D_\Gamma(u)
 =\tfrac12D_\Gamma(u),
\]
which is \eqref{eq:ch28-curve-dual-energy}, the third expression being the
same computation read backwards.  Finally $D_\Gamma(u)=0$ forces $u$
constant on the connected graph $\Gamma$, hence
$\mu=\mu_0+\Delta_\Gamma u=\mu_0$; and conversely $u=0$ solves
$\MA_\omega(u)=\mu_0$.
\end{proof}

\begin{remark}[The classical energy of a measure]
\label{rem:ch28-classical-energy}
Formula \eqref{eq:ch28-curve-dual-energy} identifies the Legendre dual
energy of $\mu$ with one half of the Dirichlet energy of its Green
potential relative to $\mu_0$.  Writing $g_\Gamma$ for the Green kernel
of \Cref{prop:graph-ma-solution} and $\nu=\mu-\mu_0$, which has total mass
zero, one has $u(x)=\int g_\Gamma(x,y)\,\rd\nu(y)$ and therefore
\[
 E^\vee_\omega(\mu)
 =-\tfrac12\iint_{\Gamma\times\Gamma}
    g_\Gamma(x,y)\,\rd\nu(x)\,\rd\nu(y),
\]
the classical energy of the signed measure $\nu$, with the sign that makes
it nonnegative.  This is the precise sense in which the non-Archimedean dual
energy generalizes the logarithmic energy of classical potential theory; the
higher-dimensional theory replaces the quadratic form by the mixed
intersection inequalities of \Cref{thm:ch26-mixed-energy-identities}.
\end{remark}

\subsection{A solved equation on a four-component chain}
\label{sec:ch28-chain}

\begin{example}[{The chain $\Gamma=[0,3]$: a solved equation, verified}]
\label{ex:ch28-chain-dual-energy}
Continue \Cref{ex:ch20-chain-three} and \Cref{ex:ch26-chain-energy}:
$\Gamma=[0,3]$ subdivided at the integers, all four vertices
$x_0,\ldots,x_3$ and three edges of length one, so $L(\Gamma)=3$; the
reference measure is $\mu_0=\delta_{x_0}+\delta_{x_3}$ and $V=2$.  A
model function is a piecewise affine $u$ with vertex values
$(u_0,u_1,u_2,u_3)$, and by \eqref{eq:ch31-graph-laplacian} its
Monge--Amp\`ere measure is the atomic measure whose mass at $x_v$ is
\begin{equation}
\label{eq:ch28-chain-masses}
 \begin{aligned}
 m_0&=1+(u_1-u_0),
 &m_1&=(u_0-u_1)+(u_2-u_1),\\
 m_2&=(u_1-u_2)+(u_3-u_2),
 &m_3&=1+(u_2-u_3),
 \end{aligned}
\end{equation}
and $u$ is $\omega$-psh exactly when all four entries are nonnegative,
which is \eqref{eq:ch20-chain-constraints}.  By
\eqref{eq:ch26-chain-energy},
\begin{equation}
\label{eq:ch28-chain-energy}
 E_\omega(u)=u_0+u_3
   -\tfrac12\bigl((u_1-u_0)^2+(u_2-u_1)^2+(u_3-u_2)^2\bigr).
\end{equation}

\emph{The equation $\MA_\omega(u)=2\delta_{x_2}$.}  Setting
\eqref{eq:ch28-chain-masses} equal to $(0,0,2,0)$ and solving from the
left: $1+(u_1-u_0)=0$ gives $u_1-u_0=-1$; then
$(u_0-u_1)+(u_2-u_1)=0$ gives $u_2-u_1=-1$; then
$(u_1-u_2)+(u_3-u_2)=2$ gives $u_3-u_2=1$; and the fourth equation
$1+(u_2-u_3)=1-1=0$ is automatically satisfied, as it must be, the two
sides having equal total mass $2=V$.  Normalizing by
$\max_vu_v=0$ gives the unique solution
\begin{equation}
\label{eq:ch28-chain-solution}
 u=(u_0,u_1,u_2,u_3)=(0,-1,-2,-1).
\end{equation}

\emph{Verification.}  The four masses \eqref{eq:ch28-chain-masses} are
\[
 \bigl(1-1,\ 1-1,\ 1+1,\ 1-1\bigr)=(0,0,2,0),
\]
all nonnegative, so $u$ is indeed $\omega$-psh, and
$\MA_\omega(u)=2\delta_{x_2}$ of total mass $2=V$.  The three edge
slopes are $-1,-1,+1$.

\emph{Energy and dual energy.}  From \eqref{eq:ch28-chain-energy},
\[
 E_\omega(u)=0+(-1)-\tfrac12(1+1+1)=-\tfrac52 ,
\]
and $\int u\,\rd\mu=2u_2=-4$, so
\[
 F_\mu(u)=E_\omega(u)-\int u\,\rd\mu=-\tfrac52+4=\tfrac32 .
\]
Independently, \Cref{prop:ch28-curve-dual-energy} gives
\[
 E^\vee_\omega(2\delta_{x_2})=\tfrac12D_\Gamma(u)
 =\tfrac12\bigl(1^2+1^2+1^2\bigr)=\tfrac32 ,
\]
in agreement.  The two computations use different formulas ---
\eqref{eq:ch28-dual-energy-of-ma} and
\eqref{eq:ch28-curve-dual-energy} --- and their agreement is a check on the
sign conventions of \Cref{app:normalizations}.

\emph{Three more measures.}  The same procedure gives the following table;
every solution is normalized by $\max_vu_v=0$, and the last column is
computed twice, as $E_\omega(u)-\int u\,\rd\mu$ and as
$\tfrac12D_\Gamma(u)$.
\begin{center}
\small
\begin{tabular}{c|c|c|c|c}
\toprule
$\mu$&solution $u$&$E_\omega(u)$&$\int u\,\rd\mu$
 &$E^\vee_\omega(\mu)$\\
\midrule
$\delta_{x_0}+\delta_{x_3}$&$(0,0,0,0)$&$0$&$0$&$0$\\
$\delta_{x_1}+\delta_{x_2}$&$(0,-1,-1,0)$&$-1$&$-2$&$1$\\
$\delta_{x_2}+\delta_{x_3}$&$(0,-1,-2,-2)$&$-3$&$-4$&$1$\\
$2\delta_{x_2}$&$(0,-1,-2,-1)$&$-\tfrac52$&$-4$&$\tfrac32$\\
\bottomrule
\end{tabular}
\end{center}
For the second row the slopes are $(-1,0,1)$, so
$\tfrac12D_\Gamma=\tfrac12(1+0+1)=1$; for the third they are
$(-1,-1,0)$, so $\tfrac12D_\Gamma=1$; the first row is the reference,
where $u=0$ and both expressions vanish, confirming the equality case of
\Cref{prop:ch28-curve-dual-energy}.

\emph{The slope constraints.}  Writing $s_j:=u_j-u_{j-1}$ for the three
edge slopes, the four nonnegativity conditions
\eqref{eq:ch28-chain-masses} read
\[
 1+s_1\geq0,\qquad s_2-s_1\geq0,\qquad s_3-s_2\geq0,\qquad 1-s_3\geq0 ,
\]
that is
\begin{equation}
\label{eq:ch28-chain-slope-window}
 -1\ \leq\ s_1\ \leq\ s_2\ \leq\ s_3\ \leq\ 1 .
\end{equation}
For model functions determined on this chain these inequalities
characterize the $\omega$-psh functions: they are convex and affine on
its three edges, with slopes in $[-1,1]$.  This is
\eqref{eq:ch20-chain-constraints} in slope coordinates.

\emph{The coercivity constant, and how lossy it is.}  By
\eqref{eq:ch28-chain-slope-window} every normalized $u$ satisfies
$\operatorname{osc}_\Gamma u\leq|s_1|+|s_2|+|s_3|\leq L(\Gamma)=3$, hence
$-3\leq u\leq0$ on $\Gamma$.  The same bound holds for the
restriction of any normalized $\omega$-psh function: the curve
retraction and slope formulas give a convex restriction with all
one-sided slopes in $[-1,1]$, and the retraction inequality shows that
its maximum on $\Gamma$ is its global supremum.  This follows by model
regularization from \Cref{prop:ch21-curve-retraction},
\eqref{eq:ch21-retraction-inequality}, and
\eqref{eq:ch21-curve-psh-test}.  Therefore
\eqref{eq:ch28-uniform-integrability} gives $M(\mu)\leq3V=6$ for every
$\mu$ carried by $\Gamma$, which is the bound
\eqref{eq:ch28-skeleton-M-bound} with $C=1$ and
$\delta(\Gamma)=3$.  On the other hand
\eqref{eq:ch28-chain-slope-window} gives
$D_\Gamma(u)=s_1^2+s_2^2+s_3^2\leq3$ for these model functions.
For an arbitrary positive measure carried by $\Gamma$, its graph
solution has distributionally increasing derivative with values in
$[-1,1]$, by the endpoint mass constraints and positivity of the
interior Laplacian.  Hence $\int_\Gamma|u'|^2\,\rd s\leq3$ as well.
Thus by
\Cref{prop:ch28-curve-dual-energy} the dual energy of any measure carried by
$\Gamma$ is at most $\tfrac32$; and the value $\tfrac32$ is attained,
for instance at $\mu=2\delta_{x_0}$, whose equations
$\eqref{eq:ch28-chain-masses}=(2,0,0,0)$ give
$s_1=s_2=s_3=1$ and hence, after normalizing $\max_vu_v=0$, the solution
\[
 u=(-3,-2,-1,0),
 \qquad
 E_\omega(u)=(-3)+0-\tfrac32=-\tfrac92,
 \qquad
 E^\vee_\omega(2\delta_{x_0})=\tfrac32 .
\]
The check through \eqref{eq:ch28-dual-energy-of-ma} is
$E_\omega(u)-\int u\,\rd\mu=-\tfrac92-2\cdot(-3)=\tfrac32$, as it must be.
So the coercivity bound $M(\mu)\leq6$ overshoots the true maximum
$\tfrac32$ of the dual energy by a factor of four; the loss comes from
replacing the Dirichlet integral by the crude bound
$\operatorname{osc}_\Gamma u\cdot V$.  This is typical: coercivity
constants are qualitatively correct and quantitatively wasteful.
\end{example}

\subsection{Stability with numbers}
\label{sec:ch28-stability-numbers}

\begin{example}[An affine family, and the stability constant]
\label{ex:ch28-affine-family}
Let $\Gamma=[x_0,x_1]$ be the two-vertex skeleton of
\Cref{ex:two-vertex-ma}, of length $L(\Gamma)=1$, with
$\mu_0=\delta_{x_0}+\delta_{x_1}$ and $V=2$.  Parametrize $\Gamma$ by
$t\in[0,1]$.  For $\varepsilon\in[0,1]$ set
\[
 \mu_\varepsilon:=(1+\varepsilon)\delta_{x_0}
                 +(1-\varepsilon)\delta_{x_1}
 \ \in\ \mathcal M_2(X).
\]

\emph{The solution.}  By \eqref{eq:two-component-ma} with
$m_0=m_1=1$, a function with $u(x_0)=a$, $u(x_1)=b$ has
$\MA_\omega(u)=(1+b-a)\delta_{x_0}+(1+a-b)\delta_{x_1}$.  Setting this
equal to $\mu_\varepsilon$ gives $b-a=\varepsilon$, and the
normalization $\max u=0$ gives
\begin{equation}
\label{eq:ch28-affine-solution}
 u_\varepsilon(t)=\varepsilon(t-1),
 \qquad
 u_\varepsilon(x_0)=-\varepsilon,\quad u_\varepsilon(x_1)=0 .
\end{equation}
The slope is $\varepsilon\in[0,1]$, so $u_\varepsilon$ is
$\omega$-psh by \Cref{ex:two-vertex-ma}, whose admissible region is
$|a-b|\leq1$.

\emph{The three quantities.}  From \eqref{eq:ch26-graph-energy},
\[
 E_\omega(u_\varepsilon)
 =u_\varepsilon(x_0)+u_\varepsilon(x_1)-\tfrac12\varepsilon^2
 =-\varepsilon-\tfrac{\varepsilon^2}2 ,
\]
and
$\int u_\varepsilon\,\rd\mu_\varepsilon
=(1+\varepsilon)(-\varepsilon)+(1-\varepsilon)\cdot0
=-\varepsilon-\varepsilon^2$, so
\[
 F_{\mu_\varepsilon}(u_\varepsilon)
 =-\varepsilon-\tfrac{\varepsilon^2}2+\varepsilon+\varepsilon^2
 =\frac{\varepsilon^2}2 ,
\]
which agrees with \Cref{prop:ch28-curve-dual-energy}:
$E^\vee_\omega(\mu_\varepsilon)=\tfrac12D_\Gamma(u_\varepsilon)
=\tfrac12\varepsilon^2$.

\emph{The stability estimate.}  For
$\varepsilon,\varepsilon'\in[0,1]$,
\[
 u_\varepsilon-u_{\varepsilon'}
 =(\varepsilon-\varepsilon')(t-1),
 \qquad
 \lVert u_\varepsilon-u_{\varepsilon'}\rVert_{C^0}
 =|\varepsilon-\varepsilon'| ,
\]
the supremum being attained at $t=0$; while
$\mu_\varepsilon-\mu_{\varepsilon'}
=(\varepsilon-\varepsilon')(\delta_{x_0}-\delta_{x_1})$ has total variation
$\lVert\mu_\varepsilon-\mu_{\varepsilon'}\rVert
=2|\varepsilon-\varepsilon'|$.  The estimate
\eqref{eq:ch28-curve-stability-sup} predicts
\[
 \lVert u_\varepsilon-u_{\varepsilon'}\rVert_{C^0}
 \ \leq\ \tfrac12L(\Gamma)
       \lVert\mu_\varepsilon-\mu_{\varepsilon'}\rVert
 =|\varepsilon-\varepsilon'| ,
\]
and equality holds.  Thus the coefficient $L(\Gamma)/2$ is already sharp
on a one-edge skeleton.  The Dirichlet form obeys
\eqref{eq:ch28-stability-dirichlet-identity} exactly:
\[
 \begin{aligned}
 D_\Gamma(u_\varepsilon-u_{\varepsilon'})
 &=(\varepsilon-\varepsilon')^2,\\
 -\int_\Gamma(u_\varepsilon-u_{\varepsilon'})\,
  \rd(\mu_\varepsilon-\mu_{\varepsilon'})
    &=-(\varepsilon-\varepsilon')(-1)(\varepsilon-\varepsilon')\\
 &=(\varepsilon-\varepsilon')^2 .
 \end{aligned}
\]
Finally
$|E^\vee_\omega(\mu_\varepsilon)-E^\vee_\omega(\mu_{\varepsilon'})|
=\tfrac12|\varepsilon^2-\varepsilon'^2|\to0$ as
$\varepsilon'\to\varepsilon$, so $\mu_{\varepsilon'}\to\mu_\varepsilon$
strongly in the sense of \Cref{def:ch28-strong-measure-convergence}, in
accordance with \Cref{prop:fixed-skeleton-stability}.

\emph{The two energies are not comparable from below.}  As
$\varepsilon\downarrow0$,
\[
 E^\vee_\omega(\mu_\varepsilon)=\frac{\varepsilon^2}2,
 \qquad
 -E_\omega(u_\varepsilon)=\varepsilon+\frac{\varepsilon^2}2,
 \qquad
 \frac{E^\vee_\omega(\mu_\varepsilon)}{-E_\omega(u_\varepsilon)}
 =\frac{\varepsilon}{2+\varepsilon}\ \longrightarrow\ 0 .
\]
This is the assertion made in \Cref{rem:ch28-reading-dual-of-ma}: the lower
bound $0$ in \eqref{eq:ch28-dual-energy-bracket} cannot be improved to any
positive multiple of $-E_\omega(\varphi)$.  For contrast, on a chain of
$N$ unit edges with $\mu_0=\delta_{x_0}+\delta_{x_N}$ and $V=2$, the
measure $\mu=2\delta_{x_0}$ has solution
$u=(-N,-N+1,\ldots,-1,0)$, all slopes equal to $1$, whence
\[
 D_\Gamma(u)=N,
 \qquad
 E^\vee_\omega(2\delta_{x_0})=\frac N2 ,
 \qquad
 E_\omega(u)=u_0+u_N-\frac N2=-\frac{3N}2 ,
\]
and the ratio $E^\vee_\omega/(-E_\omega)=\tfrac13$ is bounded away from
both $0$ and the upper bound $n=1$.  The family above realizes
$\varepsilon/(2+\varepsilon)$, hence only $(0,\tfrac13]$; the extreme
value $1$ occurs too.  For $\mu=2\,\rd s$ on $\Gamma=[0,1]$ the
solution is $u(t)=t^2-t$, with $D_\Gamma(u)=\tfrac13$ and
$\int u\,\rd\mu_0=u(0)+u(1)=0$, so
$E^\vee_\omega/(-E_\omega)=1$.  Interpolating with
$u_{\lambda,s}(t)=\lambda(t^2-t)+s(t-1)$, where $\lambda,s\geq0$ and
$\lambda+s\leq1$, gives ratio
$\bigl(\tfrac12(\tfrac{\lambda^2}3+s^2)\bigr)
 /\bigl(\tfrac12(\tfrac{\lambda^2}3+s^2)+s\bigr)$,
so the whole range $(0,1]$ occurs already on curves.
\end{example}

\subsection{A finite-energy measure not carried by any dual complex}
\label{sec:ch28-not-skeletal}

\begin{example}[An unbounded potential and its measure]
\label{ex:ch28-unbounded-dual-energy}
Keep the setting of \Cref{ex:ch26-unbounded-finite-energy}: $K$ complete
discretely valued, $X=\mathbf P^1_K$, $L=\cO(1)$, $\omega$ the
curvature of the model metric of $(\mathbf P^1_R,\cO(1))$, so
$\mu_0=\delta_\eta$ and $V=1$; $R_0=[\eta,0]$ the ray with parameter
$t$; and $\varphi=u(\mathbf t)$ the $\omega$-psh function determined by
the piecewise affine $u$ with $u(0)=0$ and slope $-1/k$ on
$[k-1,k]$, so that $u(k)=-H_k$.  There it was computed that
\[
 \begin{aligned}
 \MA_\omega(\varphi)
   &=\sum_{k\geq1}\frac1{k(k+1)}\delta_{t=k}=: \mu,\\
 E_\omega(\varphi)&=-\frac{\pi^2}{12},\\
 \int(-\varphi)\MA_\omega(\varphi)&=\frac{\pi^2}6 .
 \end{aligned}
\]

\emph{Its dual energy.}  By \eqref{eq:ch28-dual-energy-of-ma},
\begin{equation}
\label{eq:ch28-unbounded-dual-energy}
 E^\vee_\omega(\mu)
 =E_\omega(\varphi)+\int(-\varphi)\MA_\omega(\varphi)
 =-\frac{\pi^2}{12}+\frac{\pi^2}{6}
 =\frac{\pi^2}{12}\approx0.8225 .
\end{equation}
So $\mu\in\mathcal M^1_1(\mathbf P^1,\omega)$, and the bracket
\eqref{eq:ch28-dual-energy-bracket} reads
$0\leq\pi^2/12\leq n(-E_\omega(\varphi))=\pi^2/12$ with $n=1$: the upper
bound is attained, so the constant $n$ in
\eqref{eq:ch28-dual-energy-bracket} is sharp for $n=1$.

\emph{It is not carried by a dual complex.}  The measure is carried by
$\{t=k:k\geq1\}$, an infinite set of divisorial points on $R_0$
accumulating at the rational point $z=0$.  Its topological support also
contains that limit point, although it has no atom there.  A dual complex of an
SNC model of $\mathbf P^1$ is a finite metric graph; its intersection
with $R_0$ lies in a bounded segment $[0,m]$, so it contains only
finitely many of the points $t=k$.  Hence $\mu$ is carried by no dual
complex, and \Cref{lem:ch28-skeleton-admissible} does not apply.

\emph{It is nevertheless solvable.}  By
\Cref{prop:ch28-image-is-solvable}, $F_\mu$ attains its maximum at
$\varphi$ itself, with value $\pi^2/12$.  The solution $\varphi$ is
\emph{unbounded}, so it is not continuous, and
\Cref{thm:bfj-supported-cy} does not apply to $\mu$ --- consistently, since
$\mu$ violates that theorem's support hypothesis.  This is the announced
example separating the two sufficient conditions of
\Cref{sec:ch28-coercivity} and \Cref{sec:ch28-image-solvable}.

\emph{Is $\mu$ admissible?}  No.  Uniform integrability fails:
testing \eqref{eq:ch28-uniform-integrability} against the normalized psh
functions $\lambda\varphi_0^{\langle m\rangle}$ of
\Cref{ex:ch28-infinite-dual-energy} gives
\[
 \begin{aligned}
 -\int\lambda\varphi_0^{\langle m\rangle}\,\rd\mu
 &=\lambda\sum_{k\geq1}\frac{\min\{k,m\}}{k(k+1)}\\
 &=\lambda\Bigl(\sum_{k=1}^{m}\frac1{k+1}
    +m\sum_{k>m}\frac1{k(k+1)}\Bigr)\\
 &=\lambda\Bigl(H_{m+1}-1+\frac{m}{m+1}\Bigr)
 \ \xrightarrow[m\to\infty]{}\ +\infty
 \end{aligned}
\]
for each fixed $\lambda\in(0,1]$, since $H_{m+1}\to\infty$; here
$\sum_{k>m}\frac1{k(k+1)}=\frac1{m+1}$ by telescoping.
Hence $M(\mu)=+\infty$, while $E^\vee_\omega(\mu)=\pi^2/12<\infty$.
The logarithmic test $\varphi_0(t)=-t$ itself has
$\int\varphi_0\,\rd\mu=-\infty$, so a finite real pairing on all of
$\PSH_0$ fails as well.
This shows that the implication
$\text{uniformly integrable}\Rightarrow\text{finite dual energy}$ of
\Cref{lem:ch28-coercivity} is strict: uniform integrability is a genuinely
stronger condition than finiteness of the dual energy.  It is useful for
the shorter compactness argument on the whole normalized psh space;
general finite-energy existence uses the quantitative energy-level
argument instead.

\emph{Continuity does not require skeletal support.}  On the same ray
take $\tau(t)=-1+e^{-t}$, extend it constantly on branches attached to
the ray, set it equal to $0$ outside the open unit disk, and set
$\tau=-1$ at the rational endpoint $z=0$.  Equivalently,
$\tau(z)=-1+\min\{1,|z|\}$, so it is continuous on $\mathbf P^{1,\An}$.
Convex piecewise affine interpolants with rational breakpoints, slopes
in $[-1,0]$, and constant tails converge uniformly to $\tau$; they are
$\omega$-psh by the curve slope formula.  Thus $\tau$ is continuous
$\omega$-psh.  Its slope is $-e^{-t}$, whose value $-1$ at the Gauss
point cancels the reference atom and whose limit at the rational
endpoint is zero.  Consequently
\[
 \MA_\omega(\tau)=e^{-t}\,\rd t\quad\text{on }R_0,\qquad
 E_\omega(\tau)=-\tfrac14,\qquad
 E_\omega^\vee(\MA_\omega(\tau))=\tfrac14.
\]
These energy values follow from
$\tfrac12\int_0^\infty e^{-2t}\,\rd t=\tfrac14$ and
$\tau(\eta)=0$; the pairing is $-\tfrac12$.
Every tail of $R_0$ carries positive measure, whereas every dual
complex intersects it in a bounded subset.  This measure is therefore
carried by no dual complex despite having a continuous solution.
\end{example}

\subsection{Two structural checks}
\label{sec:ch28-two-checks}

\begin{example}[A Dirac mass at a divisorial point]
\label{ex:ch28-divisorial-green}
Assume orthogonality, let $x$ be divisorial, and define
\begin{equation}
\label{eq:green-envelope-divisorial}
 \varphi_x:=\sup\{\psi\in\PSH(X,\omega)\mid\psi(x)\leq0\}.
\end{equation}
The precise identification is \cite[Lemma~8.5 and equation~(8.7)]{BFJ15}:
this pointwise supremum is the continuous solution, normalized by
$\varphi_x(x)=0$, and
\[
                      \MA_\omega(\varphi_x)=V\delta_x.
\]
This uses the supported existence and domination arguments; a constraint
at one point is not itself a continuous obstacle to which the envelope
differentiability theorem can be applied directly.
Finite generation of the section algebra described in
\cite[Remark~8.8]{BFJ15} is sufficient for a rational model Green function.
Over $K=\mathbf C((\pi))$ with a rational semipositive form,
\cite[Theorem~A(ii)]{BJGreen24} gives divisorial examples whose Green
functions are not $\mathbf Q$-PL.  Those examples are nevertheless
$\mathbf R$-PL, so they do not show failure of real-model regularity.
Theorem~B of that source gives a finite atomic measure on
$\mathbf P^3_K$ whose continuous solution is invariant under no SNC
retraction, and hence is not $\mathbf R$-PL.

Three consequences for the present chapter.  First, $V\delta_x$ is carried
by any dual complex containing $x$, hence is admissible by
\Cref{lem:ch28-skeleton-admissible}, so
\Cref{thm:ch28-existence-admissible} and
\Cref{thm:bfj-supported-cy} apply and produce a unique continuous normalized
solution; the example identifies it, granted orthogonality, with the
envelope \eqref{eq:green-envelope-divisorial}.  Second,
\Cref{prop:ch28-dual-energy-of-ma} evaluates its dual energy in closed form:
\[
 E^\vee_\omega(V\delta_x)=E_\omega(\varphi_x)-V\varphi_x(x).
\]
Both terms are computable once a normalization of $\varphi_x$ has been
fixed, and the expression is invariant under adding a constant, since
$E_\omega$ shifts by $cV$ and the second term by $cV$ as well.  With
the Green normalization $\varphi_x(x)=0$ of \cite{BJGreen24}, it reduces to
$E_\omega(\varphi_x)$, which is then $\geq0$ because $\varphi_x\geq0$
--- the constant $0$ being a competitor in
\eqref{eq:green-envelope-divisorial} --- and $E_\omega$ is increasing by
\Cref{thm:extended-energy-basic}.  Third, Theorem~A excludes an assumption of rational-model
regularity, while the finite-atomic example of Theorem~B excludes even
real-model regularity for all supported solutions.  The variational
proof therefore works with continuous potentials rather than assuming
that its maximizer is a model function.
\end{example}

\begin{example}[The graph case]
\label{ex:ch28-graph-case}
On a compact skeleton $\Gamma$ of a curve, write
$\MA_\omega(u)=\omega+\Delta u$.  Solving
$\MA_\omega(u)=\mu$ is the linear equation
\[
                          \Delta u=\mu-\omega,
\]
whose right side has total mass zero.  The Green kernel gives a continuous
solution unique modulo constants.  In this dimension the Legendre energy is
the classical energy of a measure and the variational functional is a
strictly concave quadratic form.  The higher-dimensional theory replaces
this Hilbert-space calculation by mixed intersection inequalities.

The precise statements behind this paragraph are, in order:
\Cref{prop:graph-ma-solution} for existence and uniqueness of the graph
solution; \Cref{prop:ch28-curve-dual-energy} and
\Cref{rem:ch28-classical-energy} for the identification of the dual energy
with the classical energy; \Cref{prop:ch28-curve-uniqueness} for uniqueness
of the analytic solution; and
\Cref{thm:ch26-mixed-energy-identities} together with
\Cref{thm:ch20-mixed-hodge} for the higher-dimensional substitute.  Strict
concavity of $F_\mu$ on the class of potentials factoring through this fixed skeleton,
modulo constants, holds because
\eqref{eq:ch26-graph-energy} exhibits $E_\omega$ as an affine functional
minus $\tfrac12D_\Gamma$, and $D_\Gamma$ is a positive definite
quadratic form modulo constants; that is exactly the equality clause of
\Cref{prop:ch20-relative-hodge} in dimension one, and it is what fails to be
available for $n\geq2$ with merely nef classes.
\end{example}

\section{Chapter dictionary}
\label{sec:ch28-dictionary}

\subsection{Objects and their variational meaning}
\label{sec:ch28-dict-objects}

\begin{center}
\small
\begin{tabular}{p{.26\textwidth}|p{.32\textwidth}|p{.34\textwidth}}
\toprule
\textbf{Object}&\textbf{Variational meaning}&\textbf{Established at}\\
\midrule
$E^\vee_\omega(\mu)$
 &concave conjugate of $E_\omega$; the value $\sup F_\mu$
 &\Cref{def:dual-energy-measure},
  \Cref{prop:ch28-sup-is-dual-energy}\\
$\mu(X^{\An})=V$
 &forced by translation invariance, not imposed
 &\Cref{lem:ch28-mass-forced}\\
$\mathcal M^1_V(X,\omega)$
 &convex set with weakly compact dual-energy sublevels
 &\Cref{prop:ch28-dual-order}\\
$E^\vee_\omega(\MA_\omega\varphi)
 =E_\omega(\varphi)-\int\varphi\,\MA_\omega(\varphi)$
 &Fenchel equality at a solution
 &\Cref{prop:ch28-dual-energy-of-ma},
  \Cref{cor:finite-energy-variational-characterization}\\
$J_\omega(\varphi,\psi)$
 &the defect in the Fenchel inequality; $=J_\varphi(\psi)$
 &\Cref{prop:ch28-fenchel-defect}\\
uniform integrability $M(\mu)<\infty$
 &coercivity of $F_\mu$ at rate one
 &\Cref{def:ch28-admissible-measure},
  \Cref{lem:ch28-coercivity}\\
continuous pairing
 &upper semicontinuity of $F_\mu$ on $\PSH_0$
 &\Cref{def:ch28-admissible-measure}\\
support on a dual complex
 &admissibility, with
  $M(\mu)\leq C\delta(\Delta_{\mathcal X})V$
 &\Cref{lem:ch28-skeleton-admissible}\\
maximizer of $F_\mu$
 &solution of $\MA_\omega(\varphi)=\mu$
 &\Cref{prop:variational-maximizer-solves},
  \Cref{thm:ch28-existence-admissible}\\
strong measure convergence
 &weak convergence $+$ convergence of dual energies
 &\Cref{def:ch28-strong-measure-convergence},
  \Cref{thm:finite-energy-measure-stability}\\
$\delta_\omega$ on measures
 &complete quasi-metric for the strong topology
 &\Cref{thm:finite-energy-measure-completion} (quoted)\\
$E^\vee_\omega(\mu)=\tfrac12D_\Gamma(u)$
 &classical energy of a measure on a graph
 &\Cref{prop:ch28-curve-dual-energy}\\
\bottomrule
\end{tabular}
\end{center}

\subsection{The four statements that must never be merged}
\label{sec:ch28-dict-four}

\begin{center}
\small
\begin{tabular}{p{.22\textwidth}|p{.38\textwidth}|p{.32\textwidth}}
\toprule
\textbf{Name}&\textbf{Statement}&\textbf{Status here}\\
\midrule
Finite dual energy
 &$E^\vee_\omega(\mu)<\infty$, that is, $F_\mu$ bounded above
 &definition; \Cref{prop:ch28-sup-is-dual-energy}\\
Uniform integrability
 &$\sup\{-\int\varphi\,\rd\mu:\varphi\in\PSH_0\}<\infty$
 &strictly stronger:
  \Cref{ex:ch28-unbounded-dual-energy}\\
Existence of a maximizer
 &$\sup F_\mu$ attained in $\cE^1$
 &proved for all finite-energy $\mu$ under orthogonality:
  \Cref{thm:ch28-finite-energy-existence}\\
Continuity of the solution
 &the maximizer lies in $C^0(X^{\An})$
 &a separate theorem, using support again:
  \Cref{thm:bfj-supported-regularity}\\
\bottomrule
\end{tabular}
\end{center}

Uniform integrability implies finite dual energy.  Under orthogonality,
finite dual energy already implies existence of a maximizer by
\Cref{thm:ch28-finite-energy-existence}.  Uniform integrability alone
does not supply the continuous pairing needed for the shorter proof on
all of $\PSH_0$.  Neither finite dual energy nor attainment implies
continuity: the unbounded example above proves this.  The separate
supported regularity theorem gives a sufficient condition for
continuity, and the bounded contrast above shows that its support
condition is not necessary.

\subsection{Numbers to remember}
\label{sec:ch28-dict-numbers}

\begin{center}
\small
\begin{tabular}{p{.40\textwidth}|p{.24\textwidth}|p{.28\textwidth}}
\toprule
\textbf{Quantity}&\textbf{Value}&\textbf{Where}\\
\midrule
mass of every measure in $\mathcal M^1_V$
 &$V=(\{\omega\}^n)$
 &\Cref{lem:ch28-mass-forced}\\
dual energy of the reference measure $\MA_\omega(0)$
 &$0$
 &\Cref{lem:ch28-mass-forced},
  \Cref{prop:ch28-curve-dual-energy}\\
bracket for $E^\vee_\omega(\MA_\omega\varphi)$, $\sup\varphi=0$
 &$0\leq\ \cdot\ \leq n(-E_\omega(\varphi))$
 &\eqref{eq:ch28-dual-energy-bracket}\\
coercivity constant for skeletal $\mu$
 &$M(\mu)\leq C\delta(\Delta_{\mathcal X})V$
 &\eqref{eq:ch28-skeleton-M-bound}\\
Hölder exponent in
 \eqref{eq:maximizing-mixed-energy-control}
 &$\alpha=2^{-n}$
 &\Cref{thm:finite-energy-measure-stability}\\
curve stability constant
 &$L(\Gamma)/2$
 &\eqref{eq:ch28-curve-stability-sup}\\
chain solution of $\MA_\omega(u)=2\delta_{x_2}$
 &$(0,-1,-2,-1)$, $E_\omega=-\tfrac52$
 &\Cref{ex:ch28-chain-dual-energy}\\
its dual energy
 &$\tfrac32$
 &\Cref{ex:ch28-chain-dual-energy}\\
$E^\vee_\omega$ on the chain, three atomic measures
 &$0,\ 1,\ 1,\ \tfrac32$
 &\Cref{ex:ch28-chain-dual-energy}\\
affine family: $E^\vee_\omega(\mu_\varepsilon)$
 &$\varepsilon^2/2$
 &\Cref{ex:ch28-affine-family}\\
its true and proved stability coefficient
 &$L(\Gamma)/2=\tfrac12$
 &\Cref{ex:ch28-affine-family}\\
dual energy of $\MA_\omega(\varphi)$,
 $\varphi$ the harmonic-slope potential
 &$\pi^2/12\approx0.8225$
 &\eqref{eq:ch28-unbounded-dual-energy}\\
dual energy of a Dirac at a rational point of $\mathbf P^1$
 &$+\infty$
 &\Cref{ex:ch28-infinite-dual-energy}\\
\bottomrule
\end{tabular}
\end{center}

\section*{Source notes}

\enlargethispage{2\baselineskip}

\paragraph{The original theory.}
The original variational construction, finite-energy uniqueness, capacity
regularization, and supported existence theorem are
\cite[\S\S6--8]{BFJ15}; the exact statement is its Theorem~8.1.  The
removal of algebraizability is \cite{BGGJKM20}.  The measure-side Legendre
formalism and strong finite-energy completion are presented in
\cite{BJ25FiniteEnergy}.  No assertion in this chapter says that an
arbitrary finite-energy measure has a continuous solution.
Over an arbitrary non-Archimedean field the finite-energy space can still be
defined, but extension of the Monge--Amp\`ere operator is explicitly left
open in \cite[Remark~4.4]{PS23Hybrid}; see
\Cref{prob:frontier-general-ma}.

\paragraph{The dating of the individual hypotheses.}
The table of \Cref{sec:ch28-scope} may be read as a chronology.  The
existence and continuity theorem for measures supported on a dual complex,
over a discretely valued field of residue characteristic zero and for an
algebraizable $X$, is \cite{BFJ15} (2015).  The removal of
algebraizability, by differentiability of non-Archimedean volumes, is
\cite[Theorem~6.3.3 and Corollary~6.3.4]{BGGJKM20} (2020); the underlying
volume functional is developed in \cite{BGM21}.  Continuity of envelopes and
the Calabi--Yau theorem in equal characteristic $p$, under geometric-origin
and resolution hypotheses, are \cite[Theorems~1.4, 8.2, 9.2 and~9.3]{GJKM19}
(2019).  The same paper proves ample metric-envelope continuity on
curves over arbitrary complete nontrivially valued fields,
\cite[Theorem~1.1]{GJKM19}; the separate skeletal-supported solution
statement is \Cref{thm:curves-arbitrary-field-continuity}.
Mixed characteristic,
again under resolution hypotheses, is \cite{FGK25} (2025).  The trivially
valued theory, in which the corresponding surjectivity statement is
equivalent to the envelope property, is \cite[\S\S7--12]{BJ22} (2022) and
\Cref{part:bj} of this book.  The synthetic, axiomatic treatment of
finite-energy measures and their strong topology, from which
\Cref{thm:finite-energy-measure-completion} is quoted, is
\cite{BJ25FiniteEnergy} (2025).  The synthetic measure theory applies to arbitrary projective Berkovich
spaces; it does not by itself supply the full pointwise singular
operator in all those settings.  The discretely valued hypotheses
here concern the pointwise BFJ framework, rather than the whole domain
of that measure theory.

\paragraph{The complex source of the method.}
The variational solution of a Monge--Amp\`ere equation by maximizing
$E-\int\cdot\,\rd\mu$ over a finite-energy class, the identification of the
dual energy as a Legendre transform, and the coercivity/compactness scheme of
\Cref{sec:ch28-coercivity} are those of Berman--Boucksom--Guedj--Zeriahi
\cite{BBGZ13}, resting on the equilibrium-energy differentiability of
Berman--Boucksom \cite{BB10}.  The finite-energy formalism in a big class,
including the measure-side energy and its topology, is \cite{BEGZ10}, and the
corresponding stability theory is \cite{BBEGZ16}; the reader who knows those
papers will recognize \Cref{prop:ch28-dual-energy-of-ma},
\Cref{prop:ch28-fenchel-defect} and \Cref{thm:ch28-existence-admissible} as
the non-Archimedean transcriptions of standard complex statements.  The non-Archimedean proof requires separate verification of the
operator, compactness, and orthogonality inputs.  Neither theory gives
a general supremum-norm bound for singular $\cE^1$ potentials merely
from their Dirichlet energy, nor does a complete strong metric imply
compactness of its bounded sets.  On a fixed compact curve skeleton,
the graph formula $E^\vee=\tfrac12D_\Gamma$ is the classical
Dirichlet energy of the corresponding Green potential.

\paragraph{The presentational choices made here.}
Four things are organized differently from the sources.

First, the dual energy is defined by a supremum over \emph{model} psh
functions and the enlargement of the test class is proved
(\Cref{prop:ch28-test-class}) rather than assumed.  This keeps
\Cref{prop:dual-energy-elementary} --- convexity and weak lower
semicontinuity --- available with a two-line proof, and it isolates
decreasing regularization \textup{(D3)} as the only input needed for the
enlargement.

Second, coercivity is axiomatized.  The two conditions of
\Cref{def:ch28-admissible-measure} are stated separately because they do
different work: uniform integrability produces the energy bound
\eqref{eq:ch28-sublevel-inclusion}, continuity of the pairing produces upper
semicontinuity of $F_\mu$.  Support on a dual complex implies both
(\Cref{lem:ch28-skeleton-admissible}), which is how the original theorem is
recovered; but the two implications use different parts of
\Cref{thm:bfj-skeletal-equicontinuity}, and separating them makes visible
that neither condition is implied by finiteness of the dual energy
(\Cref{ex:ch28-unbounded-dual-energy}).

Third, uniqueness is dissected.  Vanishing of $I$, the bounded
Dirichlet coefficients, and the corresponding mixed-measure differences
is proved in \Cref{prop:ch28-solution-set}.  The remaining input is the
reference-energy and ample-Hodge argument of
\cite[Proposition~8.2]{BFJ15}, described in
\Cref{rem:ch28-uniqueness-gap}.  The continuous curve proof uses the
maximum principle directly.

Fourth, the measure-side quasi-metric is quoted and kept distinct from
the maximum-based potential gauge, which fails the triangle inequality.
Its well-definedness follows from energy continuity estimates,
independently of potential uniqueness.  The finite-energy bijection
allows transport of that complete quasi-metric, without identifying it
with the gauges of \Cref{chap:bfj-energy}.

\paragraph{Nets rather than sequences.}
The compactness argument of \Cref{thm:ch28-existence-admissible} is run with
nets, and \Cref{rem:ch28-net-form} explains why: $\PSH_0(X,\omega)$ is
compact by \Cref{thm:bfj-psh-compactness} but is not known to be metrizable,
and \Cref{lem:ch21-no-cofinal-sequence} obstructs the usual reduction.  On
the other hand \Cref{rem:ch28-maximizing-sequence-exists} points out that a
maximizing \emph{sequence} always exists, because the supremum defining
$E^\vee_\omega(\mu)$ is a supremum of real numbers; the potential convergence argument uses compactness through subnets
and identifies all cluster points by uniqueness.  Thus
\Cref{thm:ch28-finite-energy-existence} proves convergence of a
normalized maximizing sequence without a sequential-compactness
theorem.  It does not provide a cofinal model regularization sequence.
The general net policy is \Cref{rem:net-vs-sequence-checklist}.

\paragraph{Green functions and the sharpness of skeletal support.}
Two constructions of Boucksom--Jonsson bound the ambitions of the theory
from opposite sides.  Their non-Archimedean Green functions
\cite[Theorem~A]{BJGreen24} show that the solution attached to a Dirac mass
at a divisorial point --- the envelope
\eqref{eq:green-envelope-divisorial} --- need not be a $\mathbf Q$-PL
function; the examples in Theorem~A(ii) are nevertheless real-PL.
They therefore rule out rational-model regularity, rather than all
real-model regularity.  Their
\cite[Theorem~B]{BJGreen24} constructs a finite atomic measure supported on a
dual complex whose solution is invariant under no SNC retraction, which is
why Step~3 of the proof of \Cref{thm:bfj-supported-cy} is a Dini-and-capacity
argument and not a reduction to the skeleton.  The corresponding failure of
retractions to preserve semipositivity in dimension $\geq2$ is the appendix
to \cite{GJKM19}.

\paragraph{Scope.}
The finite-energy bijection proved here uses the standing BFJ
pointwise operator framework and orthogonality.  It does not extend
that singular operator to arbitrary complete non-Archimedean fields.
The broader question is \Cref{prob:frontier-general-ma}, with the
precise scope warning \cite[Remark~4.4]{PS23Hybrid}.
Continuous supported solutions, finite-energy existence, and
extensions of the operator require their respective displayed
hypotheses.  The chapter is written so that each of them can be located.
